\chapter{Foundations of Angular Electromagnetic Structure}
\label{ch:01}
\setcounter{section}{-1}
% R10.3 — Figure 1.12 layout patch (inherited from R10.2)
% R10.2 — local layout patch for Fig.~1.12 (build-package only)
\makeatletter
\renewcommand{\VolIPhysReactiveVsRadiative}{%
\begin{tikzpicture}[voliphys]
  \draw[volishelld, fill=black!4] (-0.35,-0.35) rectangle (0.35,0.35);
  \draw[volisrc] (0,0) circle (1.5pt);
  \node[volilabel, font=\scriptsize, align=center] at (0,-0.78) {bounded\\source};
  \draw[volishelld] (0,0) circle (0.58);
  \foreach \a in {20,80,140,200,260,320} {
    \draw[gray!40!black, -{Latex[length=1.6mm,width=1mm]}, line width=0.55pt]
      ({0.42*cos(\a)},{0.42*sin(\a)}) arc [start angle=\a, end angle=\a+50, radius=0.42];
  }
  \node[volilabel, font=\scriptsize, align=center] at (-1.75,0.15) {reactive\\$\mathbf{P}$};
  \draw[volishell] (0,0) circle (1.38);
  \node[volilabel, font=\scriptsize] at (1.65,-0.35) {$\Gamma_R$};
  \foreach \a in {0,45,90,135,180,225,270,315} {
    \draw[volipoynt] ({1.38*cos(\a)},{1.38*sin(\a)}) -- ({1.78*cos(\a)},{1.78*sin(\a)});
  }
  \node[volilabel, font=\scriptsize, align=center] at (0,2.15)
    {$\oint_{\Gamma_R}\Re\{\mathbf{S}_\mathrm{c}\!\cdot\!\hat{\mathbf{n}}\}\,dS$};
\end{tikzpicture}%
}
\makeatother



\section{Introduction and Scope}
\label{sec:ch01.0}

Beth's torsion experiment asks a mechanical question about light. A narrow, circularly polarized beam falls on a quartz wave plate on a fine suspension; the plate can turn even when the transverse intensity appears uniform about the beam axis \cite{beth1936}. The suspension does not read a phase drawing: it registers torque transmitted through the dielectric interface as a Maxwell traction integral. That distinction---between how a field is \emph{organized} in angle and what \emph{action} it exerts on matter---is the thread of the chapter. Three physical levels must therefore be kept apart: the \emph{angular organization} of $(\E,\H)$ on a chart or shell (not a conserved quantity by itself); angular momentum \emph{associated with} the field in a declared volume under the balances of section~\ref{sec:ch01.1}; and signed angular-momentum \emph{flux} $\boldsymbol{\Phi}_{J,\mathrm{out}}$ crossing a chosen closed surface when the stress integral is well posed \cite{stratton1941,beth1936}.

With impressed sources, constitutive response, and boundary data specified, Maxwell's equations determine the field pair $(\E,\H)$ \cite{maxwell1873}. Energy transport follows the Poynting vector $\mathbf{s}=\E\times\H$; vacuum electromagnetic momentum density is $\mathbf{g}_{\mathrm{em}}=\varepsilon_0\E\times\B$, while the symmetric Maxwell stress $\mathbf{T}$ supplies momentum-flux density and surface traction $\mathbf{t}=\mathbf{T}\cdot\hat{\mathbf{n}}$ \cite{maxwell1873,poynting1884,stratton1941}. Angular momentum about an origin combines field storage in $\mathbf{r}\times\mathbf{g}_{\mathrm{em}}$ with stress-moment transport through $\mathbf{r}\times\mathbf{T}$ \cite{beth1936,stratton1941}. \emph{Angular structure} describes how that pair is organized under rotation of physical space---for instance through an azimuthal phase $e^{\jj m\phi}$ in a transverse pattern. Mechanical torque on matter is fixed by cycle-averaged traction $\mathbf{t}=\Tavg\cdot\hat{\mathbf{n}}$ and by the stress-moment balances of Subsection~\ref{sec:ch01.1.5}, not by the index $m$ alone \cite{allen1992,stratton1941}.

The same $(\E,\H)$ pair supports four observables that answer different physical questions and must not be collapsed \cite{stratton1941}: time-averaged real power $\Re(\mathbf{S}_\mathrm{c}\cdot\hat{\mathbf{n}})\,dS$ through an element of area; reactive harmonic exchange $\Im(\mathbf{S}_\mathrm{c}\cdot\hat{\mathbf{n}})\,dS$; mechanical traction $\mathbf{t}=\Tavg\cdot\hat{\mathbf{n}}$ from the Maxwell stress; and outward stress-moment flux carried by $\Mavg=\mathbf{r}\times\Tavg$. On an interface, in a structured transverse mode, or on a near-zone shell where $\Im(\mathbf{S}_\mathrm{c}\cdot\hat{\mathbf{n}})$ dominates the local normal exchange, a single local reading may emphasize one observable while obscuring the others. Matching spiral isophases in $\E$ do not determine $\oint \Mavg\cdot\hat{\mathbf{n}}\,dS$, because the traction integrand couples $\E$ to $\H$ through the constitutive law; two fields can share $\arg E$ yet yield different stress moments \cite{stratton1941,allen1992}.

\begin{figure}[t]
\centering
\VolIPhysFig{\VolIPhysThreeRadiationConsequences}
\caption{Three Maxwellian consequences of radiation: linear momentum delivered to an absorber; mechanical torque on a material interface under circular polarization (not the vanishing symmetric stress moment on a centred transverse disk for a uniform plane wave); and equal counter-rotating orbital components whose signed modal transport cancels despite a structured intensity pattern \cite{poynting1884,beth1936,allen1992}.}
\label{fig:ch1.three-radiation-consequences}
\end{figure}

Steady radiation problems in this volume are analysed in the harmonic domain. Time dependence is taken as $\exp(+\jj\omega t)$ with $\omega$ in rad\,s$^{-1}$; $\E(\mathbf{r})$ is the complex phasor, so $\Re\{\E e^{\jj\omega t}\}$ is the instantaneous field. Quadratic observables---power, stress, and torque---are cycle averages of bilinear products and are therefore assembled from $\E$ and $\H^*$ \cite{stratton1941,harrington2001}.

For such fields the \emph{complex Poynting vector} is
\begin{equation}
\label{eq:ch1.intro-poynting}
\mathbf{S}_\mathrm{c}=\tfrac{1}{2}\E\times\H^*,
\end{equation}
where $(\cdot)^*$ denotes complex conjugation \cite{poynting1884,stratton1941}.
Time-averaged real power flux density is $\langle\mathbf{S}\rangle=\Re(\mathbf{S}_\mathrm{c})$;
$\mathbf{S}_\mathrm{c}$ itself is \emph{not} the cycle-averaged Poynting vector.
The divergence $\divg\mathbf{S}_\mathrm{c}$ is net outward complex power per unit volume:
its real part balances impressed work on sources, and its imaginary part tracks harmonic
exchange between electric and magnetic stored energy (section~\ref{sec:ch01.1.4})
\cite{stratton1941}. Both $\mathbf{S}_\mathrm{c}$ and $\langle\mathbf{S}\rangle$ have
dimensions W\,m$^{-2}$; $\langle\mathbf{S}\rangle\cdot\hat{\mathbf{n}}\,dS$ is active
power per unit area crossing a surface element with outward normal $\hat{\mathbf{n}}$.

In vacuum, electromagnetic momentum is stored in $\mathbf{g}_{\mathrm{em}}=\varepsilon_0\E\times\B$ and transported through the divergence of the Maxwell stress tensor \cite{maxwell1873,stratton1941}. In real time, the instantaneous tensor $\Tinst$ enters
\begin{equation}
\label{eq:ch1.intro-momentum-balance}
\frac{\partial\mathbf{g}_{\mathrm{em}}}{\partial t}+\mathbf{f}_{\mathrm{src}}
=\divg\Tinst,
\end{equation}
with $\mathbf{g}_{\mathrm{em}}=\varepsilon_0\,\E\times\B$ and $\mathbf{f}_{\mathrm{src}}=\rho\E+\Jimp\times\B$ using the source pair $(\rho,\Jimp)$ of Eqs.~\eqref{eq:ch1.maxwell-local} \cite{maxwell1873,stratton1941}. On phasor fields, cycle-averaged traction uses $\Tavg$ of Eq.~\eqref{eq:ch1.maxwell-stress} and $\mathbf{t}=\Tavg\cdot\hat{\mathbf{n}}$ \cite{stratton1941}. Let $V_m$ enclose the material whose force and torque are sought; $\hat{\mathbf{n}}$ is outward from $V_m$. At harmonic steady state, $\mathbf{F}_{\mathrm{matter}}=\oint_{\partial V_m}\Tavg\cdot\hat{\mathbf{n}}\,dS$.

Angular momentum about the origin is carried by the stress-moment density $\Mavg=\mathbf{r}\times\Tavg$, with instantaneous counterpart $\Minst=\mathbf{r}\times\Tinst$. Locally, source exchange contributes $\boldsymbol{\tau}_{\mathrm{src}}=\mathbf{r}\times\mathbf{f}_{\mathrm{src}}$; under Route~A continuum closure (Subsection~\ref{sec:ch01.1.4}), $\boldsymbol{\tau}_{\mathrm{matter}}$ coincides with $\boldsymbol{\tau}_{\mathrm{src}}$ on the material support.

Integrating traction over a closed surface gives the stress moment $\boldsymbol{\mathcal{T}}_J=\oint_{\partial V_m}\mathbf{r}\times(\Tavg\cdot\hat{\mathbf{n}})\,dS$. Outward electromagnetic angular-momentum transport is $\boldsymbol{\Phi}_{J,\mathrm{out}}=-\boldsymbol{\mathcal{T}}_J$, as in Eqs.~\eqref{eq:ch1.stress-moment-integral} and~\eqref{eq:ch1.outward-am-flux}. In harmonic steady state on a source-free enclosing surface, with vanishing storage rate inside $V_m$, mechanical torque balances transport: $\mathbf{T}_{\mathrm{matter}}=\boldsymbol{\mathcal{T}}_J=-\boldsymbol{\Phi}_{J,\mathrm{out}}$ \cite{beth1936,stratton1941}.

A uniform plane wave in vacuum separates radiation pressure from angular transport in the simplest setting. Take
$\E=E_0 e^{-\jj k z}\hat{\mathbf{x}}$ and, from $\curl\E=-\jj\omega\mu_0\H$,
$\H=(E_0/\eta_0)\,e^{-\jj k z}\hat{\mathbf{y}}$ with
$\eta_0=\sqrt{\mu_0/\varepsilon_0}$ and $k=\omega\sqrt{\varepsilon_0\mu_0}$
\cite{stratton1941}. Then
\begin{equation}
\label{eq:ch1.intro-plane-wave-preview}
\mathbf{S}_\mathrm{c}=\tfrac12\E\times\H^*=\tfrac{|E_0|^2}{2\eta_0}\,\hat{\mathbf{z}},
\end{equation}
so the power through a disk of radius $a$ normal to $\hat{\mathbf{z}}$ is
$P=\pi a^2 |E_0|^2/(2\eta_0)$. Let the field control volume be the vacuum region
$z<0$, so that at the plane $z=0$ the outward unit normal is $\hat{\mathbf{n}}=+\hat{\mathbf{z}}$.
Momentum conservation gives radiation pressure $p_{\mathrm{abs}}=\langle S_z\rangle/c=
|E_0|^2/(2\eta_0 c)$ on a perfect absorber, directed along $+\hat{\mathbf{z}}$ into the
material \cite{stratton1941}. A perfect reflector at normal incidence receives
$p_{\mathrm{refl}}=2\langle S_z\rangle/c$. The factor two follows from momentum reversal at a perfect reflector: the mirror receives the incident normal flux plus the reversed outgoing flux, $|S|/c+|{-S}|/c=2|S|/c$ \cite{maxwell1873,stratton1941}. The vacuum slab $V_f$ ($z<0$) fixes $\hat{\mathbf{n}}=+\hat{\mathbf{z}}$ at $z=0$, whereas $V_m$ encloses the absorber whose mechanical reaction is required \cite{maxwell1873,stratton1941}. With Eq.~\eqref{eq:ch1.maxwell-stress}, $\langle T_{zz}\rangle=-\varepsilon_0|E_0|^2/2$ on the vacuum side. Transverse symmetry about $\hat{\mathbf{z}}$ then gives $\dot J_{z,\mathrm{out}}=\mathbf{0}$ even when $P>0$ \cite{beth1936}: power crosses the disk while the symmetric plane wave carries no net axial angular-momentum flux. An azimuthal factor $e^{\jj m\phi}$ reshapes the transverse pattern but does not, without the paired $\H$ from Maxwell's law, predetermine the stress-flux sign \cite{allen1992,stratton1941}.

Finite-support impressed current $\Jimp$, constitutive laws, boundary conditions, and an
outgoing branch together fix a Maxwell field $(\E,\H)$ \cite{maxwell1873,stratton1941}.
On a radiative shell $\Gamma_R$, conserved outward power is
$\oint \Re(\mathbf{S}_\mathrm{c}\cdot\hat{\mathbf{n}})\,dS$ and electromagnetic
angular-momentum transport is $-\oint \Mavg\cdot\hat{\mathbf{n}}\,dS=\boldsymbol{\Phi}_{J,\mathrm{out}}$
under Eq.~\eqref{eq:ch1.outward-am-flux} \cite{beth1936}. Reactive near-zone shells may
carry large $|\Im(\mathbf{S}_\mathrm{c}\cdot\hat{\mathbf{n}})|$ while net signed angular
flux on $\Gamma_R$ cancels by symmetry; Subsection~\ref{sec:ch01.6.2} exhibits such a pair
\cite{allen1992,chu1948}.

\begin{figure}[t]
\centering
\VolIPhysFig{\VolIPhysReactiveVsRadiative}
\caption{Bounded source on finite support, large net reactive power in the near zone (dashed shell), and outward power flux through enclosing surface $\Gamma_R$ (solid shell). Reactive shells can carry large $|\Im(\mathbf{S}_\mathrm{c}\cdot\hat{\mathbf{n}})|$ while radiative shells carry the conserved outward power \cite{poynting1884,chu1948}.}
\label{fig:ch1.master-philosophy}
\end{figure}

Figure~\ref{fig:ch1.master-philosophy} fixes the reading surface before the integrals are interpreted. Reactive exchange can dominate $\Im(\mathbf{S}_\mathrm{c}\cdot\hat{\mathbf{n}})$ on a small dashed shell about the source; conserved outward power appears through $\Re(\mathbf{S}_\mathrm{c}\cdot\hat{\mathbf{n}})$ on a radiative shell $\Gamma_R$. Poynting and angular-momentum balances may be written on any admissible closed surface, yet signed radiative transport is least ambiguous where the field branch is predominantly outgoing \cite{chu1948,stratton1941}.

Source realizability and aperture closure are developed in Subsection~\ref{sec:ch01.0.1}; phase organization without net transport, and shell-dependent reactive fields, appear in Subsections~\ref{sec:ch01.6.2} and~\ref{sec:ch01.6.4} \cite{stratton1941,harrington2001}.
\phantomsection\label{thm:ch1.realizability-hierarchy}\label{prin:ch1.rep-physics}\label{def:ch1.hierarchy-classes}\label{hier:ch1.angular-status}\label{prin:ch1.finite-geometry-angular}\label{prin:ch1.accessibility}

\begin{figure}[t]
\centering
\VolIPhysFig{\VolIPhysThreeLevelHierarchy}
\caption{Impressed current on finite support $\Omega_s$, Maxwell fields $(\E,\H)$ on enclosing surface $\Gamma_R$, and angular variation of the tangential field in a harmonic chart on $S^2$. The chart lists angular organization of an already specified field pair; it is not a substitute for impressed sources and boundary closure \cite{stratton1941}.}
\label{fig:ch1.three-level-hierarchy}
\end{figure}

\subsection{Historical lineage of field energy, momentum, and angular structure}
\label{sec:ch01.0.historical-lineage}

Faraday's induction experiments made the field, rather than distant action, the primary agent of electromagnetic interaction \cite{maxwell1873}. Maxwell's 1865 memoir and the 1873 \emph{Treatise} supplied the pair $(\E,\H)$, completed Amp\`ere's law with displacement current, and predicted waves propagating at the speed of light \cite{maxwell1873}. Once the field became a dynamical system in its own right, two mechanical questions followed immediately: where does electromagnetic energy reside, and how does the field exchange linear and rotational action with matter?

Maxwell answered the mechanical side within the same field picture. His stress tensor
relates bilinear products of $(\E,\H)$ to force and momentum density in the
electromagnetic medium \cite{maxwell1873,stratton1941}. Poynting \cite{poynting1884}
then identified energy transport with the flux $\mathbf{s}=\E\times\H$, giving a local
account of power crossing a surface. Hertz's experiments (1887--1888) showed that
Maxwell's predicted waves could be generated and detected in the laboratory, converting
the wave equation from theory into measurable radiation \cite{stratton1941}. Radiation-pressure
measurements on absorbing and reflecting surfaces confirmed that the same field carries
directed momentum as well as energy; on a perfect absorber the pressure is $p_{\mathrm{abs}}=|\langle\mathbf{S}\rangle|/c$, on a
perfect reflector $p_{\mathrm{refl}}=2|\langle\mathbf{S}\rangle|/c$ in vacuum \cite{stratton1941}.

The early twentieth century did not \emph{introduce} electromagnetic stress---Maxwell
had already done so---but reopened the question of momentum \emph{inside} material
media. Abraham \cite{abraham1910} and Minkowski \cite{minkowski1908} proposed different
symmetric stress--energy forms for polarizable matter; their disagreement concerns how
much momentum is stored in a dielectric, not whether vacuum radiation exerts pressure.
In source-free exterior regions both reduce to the vacuum Maxwell stress used throughout
this chapter \cite{stratton1941}.

Stratton \cite{stratton1941}, Sommerfeld \cite{sommerfeld1949}, Harrington
\cite{harrington2001}, and Jackson \cite{jackson1999} consolidated the mathematical
apparatus of boundary-value radiation problems: impressed sources, outgoing conditions,
dyadic Green functions, and surface flux integrals for power and momentum. Chu
\cite{chu1948} emphasized that finite source size and stored reactive energy bound what
a compact antenna can radiate, so far-zone patterns cannot be separated from near-zone
storage. Tai \cite{tai1994} unified dyadic Green-function methods in monochromatic media. Gauge-invariant helicity and dual rotations on $(\E,\H)$ belong to the field-theoretic literature developed by Belinfante and successors \cite{belinfante1940,tai1994}.

Rotational action entered the laboratory with Beth's measurement of the
mechanical torque on a wave plate in circularly polarized light \cite{beth1936}. The result
linked a macroscopic suspension reading to field angular momentum carried by the beam and
deposited through Maxwell stress at the plate surface. Belinfante
\cite{belinfante1940} clarified that canonical and kinetic momentum densities differ by
the divergence of a spin current; correspondingly, local ``orbital'' and ``spin''
labels partition total angular momentum in a convention-dependent way, while closed-surface
flux integrals remain anchored to $(\E,\H)$ \cite{belinfante1940,stratton1941}.

Allen \emph{et al.} \cite{allen1992} supplied a countable azimuthal index $m$ for paraxial
beams through $\exp(\jj m\phi)$, which greatly eased experimental access to helical beams.
The observables retained here are nonetheless the surface fluxes already named---Poynting
power, Maxwell traction, and $\Mavg=\mathbf{r}\times\Tavg$ on a radiative shell---whose
intellectual line runs from Maxwell and Poynting through Beth to the developments cited
above \cite{poynting1884,beth1936,allen1992}.
\subsection{Source data, boundaries, and Maxwell consistency}
\label{sec:ch01.0.1}

Every flux integral in this chapter presupposes a Maxwell pair $(\E,\H)$ fixed by impressed
sources, constitutive laws, boundary data on $\Gamma$, and an outgoing condition whenever
radiation is asserted outward \cite{stratton1941,maxwell1873}. A tangential pattern sketched
on an aperture rim, without the impressed current and closure that complete Maxwell's problem,
does not define Poynting power or torque integrals tied to a realizable excitation
\cite{harrington2001}.

On a PEC-backed circular aperture, a charge-continuous impressed current $\Jimp$ closes Maxwell's equations and defines radiated flux on $\Gamma_R$. Prescribing a tangential field on the rim without a matching source may satisfy a local boundary relation yet leave Poynting and torque integrals detached from any synthesizable excitation \cite{stratton1941,harrington2001}. Figure~\ref{fig:ch1.wp-matched-aperture} contrasts the two configurations.

A radiation problem is specified by the spatial domain $\Omega$, its boundary $\Gamma$, constitutive data, impressed sources $(\Jimp,\rho)$, boundary conditions, and harmonic frequency under $\exp(+\jj\omega t)$ \cite{stratton1941}.

\begin{figure}[t]
\centering
\VolIPhysFig{\VolIPhysMatchedAperture}
\caption{Same circular aperture and observation shell $\Gamma_R$: left, a charge-continuous impressed current $\Jimp$ launches a Maxwell field whose outward flux on $\Gamma_R$ is defined; right, a tangential $\E_t$ pattern is prescribed on the rim without a matching source, so flux integrals need not correspond to any realizable excitation \cite{stratton1941,harrington2001}.}
\label{fig:ch1.wp-matched-aperture}
\end{figure}

Given impressed current and charge, constitutive data, boundary conditions on
$\Gamma=\partial\Omega$, and angular frequency $\omega>0$ under $\exp(+\jj\omega t)$,
Maxwell's equations determine $(\E,\H)$ once the radiation condition is imposed
\cite{maxwell1873,stratton1941}. Statements about angular organization inherit meaning only
from that closed boundary-value problem; an expanded spherical-harmonic chart cannot repair
incompatible rim traces \cite{stratton1941,harrington2001}.
\subsection{Rotation and Maxwell compatibility}
\label{sec:ch01.0.2}

Rotating a source is not the same as rotating the coordinate chart on which its components are listed. Under a proper rotation $R$ about axis $\hat{\mathbf{u}}$ through angle $\alpha$, a vector transforms as $\mathbf{v}'=R_\alpha\mathbf{v}$ and a second-order tensor as $T'_{ij}=R_{ik}T_{kl}R_{jl}$ ($\mathbf{T}'=R\mathbf{T}R^\top$) \cite{stratton1941}. A vector field transforms actively as $(T_R\mathbf{F})(\mathbf{r})=R\,\mathbf{F}(R^{-1}\mathbf{r})$: the excitation moves in space rather than being relabelled \cite{hansen1988}.

Rotate the impressed current on a feed through $\alpha$ about the horn axis while the
horn, ground plane, and surrounding medium remain axisymmetric. The radiated field rotates
with the source if and only if the full Maxwell problem---sources, constitutive data, and
boundaries---admits the same rotational symmetry. When that symmetry holds, solving
Maxwell's equations after rotating $\Jimp$ reproduces the field obtained by rotating the
unrotated solution \cite{stratton1941,hansen1988}. An anisotropic radome or an
azimuthally graded permittivity breaks the equivalence whenever the complete
boundary-value configuration lacks rotation symmetry about the axis under test
\cite{stratton1941}.

\subsection{Assumptions, finite energy, and domain of validity}
\label{sec:ch01.0.3}

A coax-fed pyramidal horn over a ground plane illustrates how an impressed current on a finite feed support, together with conducting boundaries and an outgoing radiation condition, defines a Maxwell boundary-value problem \cite{stratton1941}. Before Poynting, stress, or angular-flux integrals are interpreted, one must know whether the underlying field pair is regular enough on the chosen control volume for those quadratic flux densities to have finite surface integrals \cite{maxwell1873,poynting1884,stratton1941}.

Physical stored energy and mathematical field integrability are related but not identical. Under the harmonic convention \(\exp(+\jj\omega t)\), the cycle-averaged electromagnetic energy density in a \emph{linear, isotropic, lossless} medium with real, frequency-independent \(\varepsilon\) and \(\mu\) at the operating frequency is \(u=\tfrac14(\varepsilon|\E|^{2}+\mu|\H|^{2})\) \cite{stratton1941}. In dispersive or dissipative media the same phasor magnitudes do not, by themselves, define a unique local stored-energy density; dissipation appears in \(\Im(\varepsilon)\) and \(\Im(\mu)\) through the complex Poynting balance of Subsection~\ref{sec:ch01.1.4} \cite{stratton1941,harrington2001}. In vacuum, \(\varepsilon=\varepsilon_0\) and \(\mu=\mu_0\), so \(u\) has SI units J\,m\(^{-3}\). By contrast, \(|\E|^{2}\) alone carries units V\(^2\)m\(^{-2}\) and \(|\H|^{2}\) carries A\(^2\)m\(^{-2}\); adding them without constitutive weights is not a single physical density. The admissibility condition used here is therefore a \emph{weighted} field norm on a specified \emph{bounded} volume \(V\): Eq.~\eqref{eq:ch1.energy-class} certifies integrability of \(\varepsilon_0|\E|^{2}+\mu_0|\H|^{2}\) over \(V\). It does not, without further trace or smoothness hypotheses, certify classical boundary values on \(\partial V\) or the finiteness of every quadratic Poynting or stress flux on \(\partial V\) \cite{stratton1941}.

Let \(\Omega\subset\mathbb{R}^{3}\) be the analysis region with boundary \(\Gamma=\partial\Omega\). For any finite control volume \(V\subset\Omega\) that excludes the support of ideal point or filament singularities, require
\begin{equation}
\label{eq:ch1.energy-class}
\int_{V}\!\left(\varepsilon_0\,\E^*\!\cdot\E+\mu_0\,\H^*\!\cdot\H\right)\,dV<\infty.
\end{equation}
The integrand is the unweighted harmonic energy density \(4u\) in source-free vacuum; Eq.~\eqref{eq:ch1.energy-class} is therefore simultaneously a finite stored-energy test on \(V\) and an \(L^2\)-type integrability condition on the weighted field pair \cite{stratton1941}. It is imposed on the bounded volume \(V\) actually used in the calculation, not as a claim of global square-integrability over all space \cite{stratton1941}.

Outgoing \(1/r\) radiation fields are generally \emph{not} globally square-integrable on an unbounded exterior domain in the ordinary \(L^2\) sense, because their energy density decays only as \(r^{-2}\) \cite{stratton1941}. Nevertheless, time-averaged radiated power \(P_{\mathrm{rad}}\), integrated Maxwell traction \(\boldsymbol{\mathcal{T}}_J\), and angular flux \(\boldsymbol{\Phi}_{J,\mathrm{out}}\) on a finite enclosing shell \(S_R\) can remain well defined when the outgoing radiation condition selects the physical branch \cite{stratton1941}. Ideal point dipoles and filaments are retained only as explicit distributional limits of finite-width, charge-continuous sources; otherwise the regular impressed-current class of Eq.~\eqref{eq:ch1.source-charge-continuous} applies. Square-integrable \(\Jimp\) and charge continuity \(\divg\Jimp+\jj\omega\rho=0\) do not by themselves guarantee Eq.~\eqref{eq:ch1.energy-class} in the exterior region, nor do they guarantee classical normal traces of \(\Jimp\) on \(\Gamma\) or finiteness of every quadratic stress or Poynting surface integral invoked later---boundary data, outgoing selection, and the field regularity of the actual solution control those conclusions \cite{stratton1941}. Boundary data on \(\Gamma\) complete the physical problem together with the source specification \cite{maxwell1873,stratton1941}.
\phantomsection\label{eq:ch1.admissible-space}

\begin{figure}[t]
\centering
\VolIPhysFig{\VolIPhysGaussianWaist}
\caption{Schematic only: a Gaussian envelope illustrates why a finite analysis domain and enclosing shell $S$ must be chosen before interpreting flux integrals; it is not, by itself, a Maxwell solution with specified sources, boundaries, and outgoing selection \cite{allen1992,stratton1941}.}
\label{fig:ch1.wp-gaussian-waist}
\end{figure}

Before flux integrals acquire physical meaning, one must choose a domain that encloses the
current support and a shell in the far-enough exterior on which the field is to be read. A
Gaussian waist quoted in isolation does not specify power or angular momentum crossing that
shell \cite{allen1992,stratton1941}. Figure~\ref{fig:ch1.wp-gaussian-waist} states the
geometric requirement only; it is not, by itself, a closed Maxwell solution
\cite{stratton1941}.

\section{Maxwell Theory and Conservation Laws}
\label{sec:ch01.1}

The opening section separated three levels of description---angular organization, volume-associated angular momentum, and signed stress-moment flux---from the four surface observables used to read them. We now develop the Maxwell theory from which those balances and integrals follow. Divergence of Amp\`ere--Maxwell yields charge continuity; the curl structure of $(\E,\H)$ yields energy and linear momentum transport; angular momentum appears as the moment of electromagnetic momentum through the symmetric Maxwell stress tensor \cite{maxwell1873,poynting1884,beth1936}. Under $\exp(+\jj\omega t)$ the local dynamics take the compact form
\begin{equation}
\label{eq:ch1.maxwell-centre}
\curl \E=-\jj\omega\mu_0\H,\qquad
\curl \H=\Jimp+\jj\omega\varepsilon_0\E,
\end{equation}
together with the divergence constraints and boundary completion of section~\ref{sec:ch01.0} \cite{stratton1941}.

Later modal and helicity labels classify radiating fields once $(\E,\H)$ is already Maxwell-compatible; they are not extra dynamical degrees of freedom that can replace Faraday's and Amp\`ere--Maxwell's coupling \cite{maxwell1873,jackson1999}. A pattern that imprints azimuthal phase on $\E$ alone but leaves $\H$ unconstrained cannot represent a physical electromagnetic state, because $\curl\E=-\jj\omega\mu_0\H$ fixes the transverse magnetic response at every point \cite{stratton1941}.

\subsection{Maxwell equations: differential, integral, vacuum constitutive relations}
\label{sec:ch01.1.1}

A time-varying electric field cannot remain an isolated scalar pattern in space. Faraday's law ties its curl to $\H$; Amp\`ere--Maxwell ties the curl of $\H$ back to $\E$ and to the impressed sources that drive the system \cite{maxwell1873,stratton1941}. Prescribing only the spatial phase of one Cartesian component---for example a forward $\exp(-\jj k z)$ factor on $E_x$---does not yet specify an electromagnetic solution. The paired $\H$ must satisfy both curl equations and the divergence constraints on the same domain \cite{stratton1941,jackson1999}. Every later transport, flux, and angular-momentum statement in this chapter presupposes that completed Maxwell pair \cite{stratton1941}.

We work in a piecewise smooth region $\Omega\subset\mathbb{R}^{3}$ with boundary $\Gamma=\partial\Omega$, adopting the time factor $\exp(+\jj\omega t)$ throughout \cite{stratton1941}. Phasors $\E(\mathbf{r})$ and $\H(\mathbf{r})$ therefore satisfy $\partial/\partial t\to +\jj\omega$. Impressed current $\Jimp(\mathbf{r})$ and charge density $\rho(\mathbf{r})$ enter Amp\`ere--Maxwell and Gauss's electric law; vacuum constants $\varepsilon_0$ and $\mu_0$ link field intensities to flux densities in SI units \cite{stratton1941,harrington2001}. Maxwell's laws may be written at a point, around a finite loop or over a control volume, or completed by the vacuum relations between $(\E,\H)$ and $(\mathbf{D},\mathbf{B})$; the chapter uses whichever form makes the physical constraint visible \cite{stratton1941}.

\begin{figure}[t]
\centering
\VolIPhysFig{\VolIPhysVacuumClosureMap}
\caption{Vacuum constitutive relations linking $(\E,\D)$ and $(\H,\B)$ through $\varepsilon_0$ and $\mu_0$, and the outgoing wave whose kinematics those constants determine.}
\label{fig:ch1.vacuum-closure-map}
\end{figure}
Figure~\ref{fig:ch1.vacuum-closure-map} summarizes how $\varepsilon_0$ and $\mu_0$ connect $(\E,\H)$ to $(\mathbf{D},\mathbf{B})$ and how the same constants enter wave speed, impedance, and wavenumber in later formulas \cite{maxwell1873,stratton1941}. Any update to $\E$ that is not propagated to $\H$ through Faraday's law breaks that closure and cannot represent a Maxwell field.

The differential Maxwell system on $\Omega$ reads
\begin{equation}
\label{eq:ch1.maxwell-local}
\begin{aligned}
\curl \E&=-\jj\omega\mu_0\H, &
\curl \H&=\Jimp+\jj\omega\varepsilon_0\E,\\
\divg\!\left(\varepsilon_0\E\right)&=\rho, &
\divg\!\left(\mu_0\H\right)&=0.
\end{aligned}
\end{equation}
Faraday's law states that a circulating $\E$ requires a time-harmonic $\H$ proportional to $\curl\E$. Amp\`ere--Maxwell adds the impressed current $\Jimp$ and the displacement current $\jj\omega\varepsilon_0\E$ that makes the law consistent with time-varying charge \cite{maxwell1873}. Gauss's electric law, $\divg(\varepsilon_0\E)=\rho$, specifies how much electric flux originates from charge at a point; it is not the same statement as charge continuity, which constrains how $\rho$ and $\Jimp$ vary together \cite{stratton1941}. Gauss's magnetic law, $\divg(\mu_0\H)=0$, expresses the absence of magnetic monopoles in the classical model \cite{maxwell1873}.

Charge continuity in harmonic form follows from the divergence of Amp\`ere--Maxwell. Because $\divg(\curl\H)=0$ identically,
\begin{equation}
\label{eq:ch1.charge-continuity-harmonic}
\divg\Jimp+\jj\omega\rho=0,
\end{equation}
obtained by taking $\divg$ of $\curl\H=\Jimp+\jj\omega\varepsilon_0\E$ and substituting Gauss's law \cite{maxwell1873,stratton1941}. The impressed sources are therefore not independent data: any $\rho$ paired with $\Jimp$ must satisfy Eq.~\eqref{eq:ch1.charge-continuity-harmonic} or the four rows of Eq.~\eqref{eq:ch1.maxwell-local} cannot hold simultaneously \cite{hansen1988}.

Integrating Eq.~\eqref{eq:ch1.charge-continuity-harmonic} over a control volume $V$ with boundary $\partial V$ and outward unit normal $\hat{\mathbf{n}}$ gives
\begin{equation}
\label{eq:ch1.continuity-integral}
\int_{\partial V}\Jimp\cdot\hat{\mathbf{n}}\,dS
+\jj\omega\int_V\rho\,dV=0.
\end{equation}
The surface integral is net impressed current crossing $\partial V$ in the outward direction; the volume integral is the harmonic charge content inside $V$. Under $\exp(+\jj\omega t)$, a positive imaginary contribution from $\jj\omega\rho$ corresponds to a $90^\circ$ phase relation between charge and its rate of change in the time domain; the balance in Eq.~\eqref{eq:ch1.continuity-integral} states that current leaving the volume must be accompanied by the harmonic charge variation required by Eq.~\eqref{eq:ch1.charge-continuity-harmonic} \cite{maxwell1873,stratton1941}.

Cartesian components make the same coupling visible on a single row:
\begin{equation}
\label{eq:ch1.maxwell-component-example}
\begin{aligned}
\partial_y E_z-\partial_z E_y&=-\jj\omega\mu_0 H_x,\\
\partial_z H_x-\partial_x H_z&=J_{\mathrm{imp},y}+\jj\omega\varepsilon_0 E_y,\\
\partial_x(\varepsilon_0 E_x)+\partial_y(\varepsilon_0 E_y)+\partial_z(\varepsilon_0 E_z)&=\rho.
\end{aligned}
\end{equation}
Each line is one scalar component of Eq.~\eqref{eq:ch1.maxwell-local} \cite{stratton1941}. Subscripts $x,y,z$ label fixed Cartesian axes; $J_{\mathrm{imp},y}$ is the $y$-component of $\Jimp$ \cite{harrington2001}. A patch of $E_z$ alone still demands paired $H_x$ through Faraday's row and paired $\rho$ through Gauss's row, so field maps that display only one component can misrepresent the underlying Maxwell state \cite{maxwell1873}.

Laboratory and device questions are posed on finite volumes and surfaces: circulation of $\E$ or $\H$ around a wire loop, flux through an aperture, or enclosed charge. Stokes' and Gauss's theorems lift the differential rows to that scale on surfaces $S\subset\Omega$ with boundary $\partial S$ and volumes $V\subset\Omega$:
\begin{equation}
\label{eq:ch1.maxwell-integral}
\begin{aligned}
\oint_{\partial S}\E\cdot d\boldsymbol{\ell}
&=-\jj\omega\mu_0\int_S\H\cdot\hat{\mathbf{n}}\,dS,\\
\oint_{\partial S}\H\cdot d\boldsymbol{\ell}
&=\int_S\Jimp\cdot\hat{\mathbf{n}}\,dS
+\jj\omega\varepsilon_0\int_S\E\cdot\hat{\mathbf{n}}\,dS,\\
\int_{\partial V}\varepsilon_0\E\cdot\hat{\mathbf{n}}\,dS
&=\int_V\rho\,dV,\qquad
\int_{\partial V}\mu_0\H\cdot\hat{\mathbf{n}}\,dS=0.
\end{aligned}
\end{equation}
The circulations of $\E$ and $\H$ around $\partial S$ are tied to fluxes of $\H$ and $\E$ through $S$; the normal flux of $\varepsilon_0\E$ through $\partial V$ equals the charge in $V$ \cite{stratton1941}. These statements match Eq.~\eqref{eq:ch1.maxwell-local} when $\E$, $\H$, $\Jimp$, and $\rho$ are sufficiently smooth for the classical integral theorems; singular or distributional sources require the same laws in weak or limit form \cite{stratton1941}. \phantomsection\label{prop:ch1.local-integral-equivalence}
Shrinking $V$ or $S$ to an interior point recovers the differential rows.

Vacuum constitutive relations complete the model in free space:
\begin{equation}
\label{eq:ch1.vacuum-constitutive}
\mathbf{D}=\varepsilon_0\E,\qquad
\mathbf{B}=\mu_0\H,\qquad
c=\frac{1}{\sqrt{\varepsilon_0\mu_0}}.
\end{equation}
Here $\mathbf{D}$ is electric flux density (C\,m$^{-2}$), $\mathbf{B}$ is magnetic flux density (T), and $c$ is the speed of light in vacuum (m\,s$^{-1}$) \cite{stratton1941}. Without Eq.~\eqref{eq:ch1.vacuum-constitutive}, the curl laws alone leave $(\mathbf{D},\mathbf{B})$ undetermined.

Combining Eq.~\eqref{eq:ch1.maxwell-local} with Eq.~\eqref{eq:ch1.vacuum-constitutive} and eliminating $\H$ from the curl pair yields second-order operators whose source-free solutions propagate at speed $c$. The intrinsic impedance, wavenumber, and wavelength used throughout the chapter are
\begin{equation}
\label{eq:ch1.vacuum-impedance-wavenumber}
\eta_0=\sqrt{\frac{\mu_0}{\varepsilon_0}},\qquad
k_0=\omega\sqrt{\mu_0\varepsilon_0}=\frac{\omega}{c},\qquad
\lambda_0=\frac{2\pi}{k_0}.
\end{equation}
The quantity $\eta_0$ (ohms) relates transverse $\E$ and $\H$ amplitudes in a uniform plane wave; $k_0$ (rad\,m$^{-1}$) measures phase accumulation per unit distance at angular frequency $\omega$; $\lambda_0=2\pi/k_0$ is the corresponding vacuum wavelength \cite{stratton1941}. These definitions are not separate postulates: they follow from the same $\varepsilon_0$, $\mu_0$, and $\omega$ that appear in Eq.~\eqref{eq:ch1.maxwell-local}.

To obtain the inhomogeneous wave equation for $\E$, use the vector identity
\begin{equation}
\label{eq:ch1.curl-curl-identity}
\curl(\curl\E)=\grad(\divg\E)-\nabla^2\E,
\end{equation}
valid for continuously differentiable fields \cite{stratton1941}. With $\exp(+\jj\omega t)$ phasors, $\curl\E=-\jj\omega\mu_0\H$ and $\curl\H=\Jimp+\jj\omega\varepsilon_0\E$. Applying $\curl$ to Faraday's law and substituting Amp\`ere--Maxwell,
\begin{equation}
\label{eq:ch1.wave-derivation-intermediate}
\curl(\curl\E)=-\jj\omega\mu_0\Jimp+k_0^2\E,
\end{equation}
because $k_0^2=\omega^2\mu_0\varepsilon_0$. Gauss's law gives $\divg\E=\rho/\varepsilon_0$, so Eq.~\eqref{eq:ch1.curl-curl-identity} rearranges to
\begin{equation}
\label{eq:ch1.vacuum-wave-source-driven}
\nabla^2\E+k_0^2\E=\frac{1}{\varepsilon_0}\nabla\rho+\jj\omega\mu_0\Jimp,\qquad
\nabla^2\H+k_0^2\H=-\curl\Jimp.
\end{equation}
The electric equation carries both the charge-gradient forcing $\varepsilon_0^{-1}\nabla\rho$ and the impressed-current term $+\jj\omega\mu_0\Jimp$ under the $+\jj\omega t$ convention adopted here \cite{jackson1999,stratton1941}. The magnetic companion is obtained independently: $\curl\H=\Jimp+\jj\omega\varepsilon_0\E$ gives $\curl(\curl\H)=\curl\Jimp+\jj\omega\varepsilon_0\curl\E=\curl\Jimp-k_0^2\H$, and the same curl--curl identity yields $\nabla^2\H+k_0^2\H=-\curl\Jimp$. Homogeneous regions with $\Jimp=\mathbf{0}$ and $\rho=0$ satisfy the vector Helmholtz pair
\begin{equation}
\label{eq:ch1.vacuum-helmholtz-pair}
\nabla^2\E+k_0^2\E=\mathbf{0},\qquad
\nabla^2\H+k_0^2\H=\mathbf{0},
\end{equation}
but Eq.~\eqref{eq:ch1.vacuum-helmholtz-pair} alone does not select a unique physical solution: boundary data on $\Gamma$, radiation conditions at infinity, and the impressed source specification still determine which field belongs to a given antenna or scatterer problem \cite{stratton1941}.

For a uniform plane-wave ansatz $\E=\E_0 e^{-\jj\mathbf{k}\cdot\mathbf{r}}$ with $\mathbf{k}\cdot\E_0=0$, Eq.~\eqref{eq:ch1.vacuum-helmholtz-pair} requires $|\mathbf{k}|=k_0$. Faraday's law gives $\H=(\mathbf{k}\times\E)/(\omega\mu_0)=\hat{\mathbf{k}}\times\E/\eta_0$ with $\hat{\mathbf{k}}=\mathbf{k}/|\mathbf{k}|$, and Amp\`ere--Maxwell confirms $\mathbf{k}\cdot\H=0$ \cite{stratton1941,jackson1999}. Transverse vacuum plane waves therefore satisfy
\begin{equation}
\label{eq:ch1.local-plane-wave-relations}
\H=\frac{1}{\eta_0}\,\hat{\mathbf{k}}\times\E,\qquad
\hat{\mathbf{k}}\cdot\E=0,\qquad
\hat{\mathbf{k}}\cdot\H=0,
\end{equation}
where $\eta_0=\sqrt{\mu_0/\varepsilon_0}$ \cite{stratton1941,jackson1999}. This is a \emph{local} Maxwell-compatible plane-wave sector: it does not describe, in isolation, the full field of a bounded radiator, which requires outgoing selection and source integrals developed in the next subsection \cite{stratton1941}.

Take $\E=E_0 e^{-\jj k_0 z}\hat{\mathbf{x}}$ propagating in $+\hat{\mathbf{z}}$. Then $\curl\E=-\jj k_0 E_0 e^{-\jj k_0 z}\hat{\mathbf{y}}=-\jj\omega\mu_0\H$, so $\H=(E_0/\eta_0)e^{-\jj k_0 z}\hat{\mathbf{y}}$. Substituting into $\curl\H=\jj\omega\varepsilon_0\E$ (source-free) closes Amp\`ere--Maxwell with no remaining free parameters. Real-time phase surfaces satisfy $\omega t-k_0 z=\mathrm{const}$, hence $dz/dt=\omega/k_0=c>0$: the disturbance advances in $+\hat{\mathbf{z}}$, and the factor $e^{-\jj k_0 z}$ is the outgoing spatial branch for $\exp(+\jj\omega t)$ \cite{stratton1941,jackson1999}. The time-average Poynting direction is $+\hat{\mathbf{z}}$ with magnitude $|E_0|^2/(2\eta_0)$ for this polarization \cite{stratton1941}.

The same source-free physics may be written in curl-curl form:
\begin{equation}
\label{eq:ch1.vacuum-curl-second-order}
\curl\!\left(\mu_0^{-1}\curl \E\right)-\omega^2\varepsilon_0\E=\mathbf{0},\qquad
\curl\!\left(\varepsilon_0^{-1}\curl \H\right)-\omega^2\mu_0\H=\mathbf{0}.
\end{equation}
Eq.~\eqref{eq:ch1.vacuum-curl-second-order} is equivalent to Eq.~\eqref{eq:ch1.vacuum-helmholtz-pair} in vacuum and is often the preferred starting point for variational formulations with tangential boundary data \cite{stratton1941}.

Eqs.~\eqref{eq:ch1.maxwell-local}--\eqref{eq:ch1.local-plane-wave-relations} fix how $\E$ and $\H$ must balance locally for given $\Jimp$ and $\rho$, but they do not yet specify which source distributions, interface traces, and outer radiation conditions define a realizably fed structure on $\Omega$. That completion---Green-function reconstruction and the Maxwell boundary-value problem---is taken up immediately below \cite{stratton1941}.
% Ch.1 PASS K — Green-function and cross-chapter bridges (Ch. 2–4 seeds)
\subsection{From inhomogeneous Helmholtz equation to outgoing Green function}
\label{sec:ch01.1.1c}

The local Maxwell system of Subsection~\ref{sec:ch01.1.1} fixes how $\E$ and $\H$ must balance at each point for prescribed $\Jimp$ and $\rho$, but it does not yet supply a constructive map from a finite source distribution to the field at a distant observation point \cite{stratton1941,maxwell1873}. That map is the classical Green-function problem: given impressed data on compact support, determine the outgoing harmonic field that satisfies the radiation condition at infinity \cite{stratton1941,jackson1999}.

Let $\mathbf{r}$ denote the observation point and $\mathbf{r}'$ the source point (both in metres), with separation $R=|\mathbf{r}-\mathbf{r}'|$. In the time domain, causality requires the field at $\mathbf{r}$ to be driven by the source at retarded time $t-R/c$, where $c=1/\sqrt{\varepsilon_0\mu_0}$ \cite{stratton1941,jackson1999}. With the harmonic convention $\exp(+\jj\omega t)$ adopted in this chapter, the delay $R/c$ appears as the spatial phase factor $e^{-\jj k_0 R}$ in the frequency-domain response, because $e^{+\jj\omega(t-R/c)}=e^{+\jj\omega t}\,e^{-\jj k_0 R}$ and $k_0=\omega/c$ \cite{stratton1941}. The minus sign in $e^{-\jj k_0 R}$ selects the outgoing spherical wave that carries energy toward increasing $R$ under this time convention \cite{stratton1941,jackson1999}.

In a homogeneous, source-free region, each Cartesian component of $\E$ or $\H$ satisfies the vector Helmholtz equation of Eq.~\eqref{eq:ch1.vacuum-helmholtz-pair}. Restoring impressed sources leads to the scalar prototype
\begin{equation}
\label{eq:ch1.inhomogeneous-helmholtz}
(\nabla^2+k_0^2)\,U(\mathbf{r})=-S(\mathbf{r}),
\end{equation}
where $U$ is a scalar potential or field component under reconstruction and $S$ is the corresponding source density whose SI units match those of $(\nabla^2+k_0^2)U$ \cite{stratton1941,jackson1999}. Equation~\eqref{eq:ch1.inhomogeneous-helmholtz} has infinitely many solutions unless boundary data and, for open domains, an outgoing radiation condition are imposed \cite{stratton1941}.

The scalar Green function $G(\mathbf{r},\mathbf{r}')$ (units m\(^{-1}\)) is defined by
\begin{equation}
\label{eq:ch1.green-defining}
(\nabla^2+k_0^2)\,G(\mathbf{r},\mathbf{r}')=-\delta(\mathbf{r}-\mathbf{r}'),
\end{equation}
where $\delta$ is the three-dimensional Dirac distribution (units m\(^{-3}\)) and $\nabla^2$ acts on the observation coordinate $\mathbf{r}$ with $\mathbf{r}'$ fixed \cite{stratton1941,jackson1999}. For $R>0$, spherical symmetry about $\mathbf{r}'$ gives $G=G(R)$. With outgoing selection, write $G(R)=g\,e^{-\jj k_0 R}/R$. Substitution into the homogeneous Helmholtz equation for $R>0$ yields $g=\mathrm{const}$. Matching the strength of the singularity at $R=0$ through the identity $\nabla^2(1/R)=-4\pi\delta(\mathbf{r}-\mathbf{r}')$ fixes $g=1/(4\pi)$ and therefore
\begin{equation}
\label{eq:ch1.scalar-green-outgoing}
G(\mathbf{r},\mathbf{r}')
=
\frac{e^{-\jj k_0|\mathbf{r}-\mathbf{r}'|}}{4\pi|\mathbf{r}-\mathbf{r}'|}.
\end{equation}
Equation~\eqref{eq:ch1.scalar-green-outgoing} satisfies Eq.~\eqref{eq:ch1.green-defining} in the distributional sense. Without the outgoing radiation condition, the incoming branch proportional to $e^{+\jj k_0 R}/R$ would also satisfy the homogeneous equation in exterior regions and the representation would not be unique \cite{stratton1941,jackson1999}.

Because Eq.~\eqref{eq:ch1.scalar-green-outgoing} already carries the $1/(4\pi R)$ normalization, the retarded potentials use $G$ directly without an additional $1/(4\pi)$ prefactor (Stratton/Jackson convention A) \cite{stratton1941,jackson1999}. With $\phi$ the scalar potential (V), $\mathbf{A}$ the vector potential (T\,m), $\Jimp$ the impressed volume current density (A\,m$^{-2}$), and $\rho$ the charge density (C\,m$^{-3}$),
\begin{equation}
\label{eq:ch1.vector-potential-retarded}
\mathbf{A}(\mathbf{r})
=
\mu_0\int
\Jimp(\mathbf{r}')\,
G(\mathbf{r},\mathbf{r}')\,
d^3r',
\end{equation}
\begin{equation}
\label{eq:ch1.lorenz-scalar-potential-retarded}
\phi(\mathbf{r})
=
\frac{1}{\varepsilon_0}\int
\rho(\mathbf{r}')\,
G(\mathbf{r},\mathbf{r}')\,
d^3r',
\end{equation}
and, under $\exp(+\jj\omega t)$,
\begin{equation}
\label{eq:ch1.lorenz-field-reconstruction}
\E=-\nabla\phi-\jj\omega\mathbf{A},\qquad \B=\nabla\times\mathbf{A},\qquad \H=\B/\mu_0,
\end{equation}
with Lorenz gauge $\divg\mathbf{A}+\jj\omega\mu_0\varepsilon_0\phi=0$ \cite{stratton1941,jackson1999}. The gauge condition is a potential constraint; charge continuity $\divg\Jimp+\jj\omega\rho=0$ is the physical law that makes $(\phi,\mathbf{A})$ consistent with the Maxwell curl system of Subsection~\ref{sec:ch01.1.1} \cite{maxwell1873,stratton1941}.

\begin{example}[Finite line segment, end charges, and far-zone phase]
\label{ex:ch1.green-dipole-element}
Consider a straight conductor on the $z$-axis with phasor current $I_0$ (A) between $z'=-h/2$ and $z'=+h/2$. A minimal volume idealization uses uniform impressed current density
\[
\Jimp(\mathbf{r}')
=\frac{I_0}{A_c}\,\hat{\mathbf{z}}\,
\chi_{A_c}(x',y')\,\chi_{[-h/2,\,h/2]}(z'),
\]
where $A_c$ (m$^2$) is the conductor cross section, $\chi_{A_c}$ is unity over the transverse cross section and zero outside, and $\chi_{[-h/2,h/2]}$ restricts the axial support. Inside the open cylinder, $\partial J_z/\partial z'=0$, so $\divg\Jimp=0$ as an ordinary function; the idealization is not solenoidal on $\mathbb{R}^3$ because $\Jimp$ vanishes outside the segment and therefore has step discontinuities at $z'=\pm h/2$. In the distributional sense,
\begin{equation}
\label{eq:ch1.segment-div-J}
\divg\Jimp
=\frac{I_0}{A_c}\,\chi_{A_c}(x',y')\bigl[\delta(z'+h/2)-\delta(z'-h/2)\bigr],
\end{equation}
where the one-dimensional deltas are combined with the transverse indicator so that Eq.~\eqref{eq:ch1.segment-div-J} has units A\,m$^{-3}$ \cite{jackson1999,stratton1941}. Charge continuity $\divg\Jimp+\jj\omega\rho=0$ then requires
\begin{equation}
\label{eq:ch1.segment-end-charge}
\rho(\mathbf{r}')
=\frac{I_0}{\jj\omega A_c}\,\chi_{A_c}(x',y')\bigl[\delta(z'-h/2)-\delta(z'+h/2)\bigr],
\end{equation}
with units C\,m$^{-3}$ in the distributional sense. Integrating $\rho$ over a thin slab at $z'=+h/2$ gives end-face charge per unit area $\sigma_s=+I_0/(\jj\omega A_c)$ (C\,m$^{-2}$); at $z'=-h/2$, $\sigma_s=-I_0/(\jj\omega A_c)$ \cite{jackson1999,stratton1941}. Under $\exp(+\jj\omega t)$, the endpoint charge at $z'=+h/2$ is in quadrature with the axial impressed current (not a $180^\circ$ out-of-phase relation): current along $+\hat{\mathbf{z}}$ leaves the upper face at $z'=+h/2$ in this sign convention---the Hertzian dipole bookkeeping for a finite segment, not a universal model of a practical feed, which generally spreads charge and return current over a finite terminal region \cite{stratton1941,jackson1999}.

With Eqs.~\eqref{eq:ch1.vector-potential-retarded} and~\eqref{eq:ch1.lorenz-scalar-potential-retarded}, the vector potential follows from the axial current alone in the limit $A_c\to 0$ as
\begin{equation}
\label{eq:ch1.line-segment-potential}
A_z(\mathbf{r})=\frac{\mu_0 I_0}{4\pi}\int_{-h/2}^{h/2}
\frac{e^{-\jj k_0|\mathbf{r}-z'\hat{\mathbf{z}}|}}{|\mathbf{r}-z'\hat{\mathbf{z}}|}\,dz',
\end{equation}
which has units T\,m; $\phi$ receives equal contributions from the end charges in Eq.~\eqref{eq:ch1.segment-end-charge}. The pair $(\phi,\mathbf{A})$ satisfies Lorenz gauge and Eq.~\eqref{eq:ch1.charge-continuity-harmonic} for this idealization, but Eqs.~\eqref{eq:ch1.lorenz-field-reconstruction} still determine a field in all of space only together with outgoing selection at infinity (and any conducting boundaries if present). The line integral for $A_z$ alone does not, by itself, specify a complete boundary-value problem \cite{stratton1941}.

\emph{Geometrical and phase approximations.} Let $R=|\mathbf{r}-\mathbf{r}'|$ and $r=|\mathbf{r}|$. For $r\gg|\mathbf{r}'|$ the distance expands as $R\approx r-\hat{\mathbf{r}}\cdot\mathbf{r}'+\mathcal{O}(|\mathbf{r}'|^2/r)$; this is purely geometrical. Replacing $e^{-\jj k_0 R}$ by $e^{-\jj k_0 r}\,e^{+\jj k_0\hat{\mathbf{r}}\cdot\mathbf{r}'}$ additionally requires the neglected phase $k_0|\mathbf{r}'|^2/r$ to be small compared with $2\pi$ across the source---for a segment of length $h$, typically $k_0 h^2/r\ll 1$ (Fraunhofer-type) \cite{stratton1941,jackson1999}. Dominance of radiative $1/r$ fields over reactive $1/r^2$ and $1/r^3$ terms is a separate condition on $k_0 r\gg 1$ in the multipole sense \cite{stratton1941}. When $r=|\mathbf{r}|\gg h$ and both criteria hold,
\begin{equation}
\label{eq:ch1.green-far-zone-phase}
e^{-\jj k_0 R}\approx e^{-\jj k_0 r}\,e^{+\jj k_0\hat{\mathbf{r}}\cdot\mathbf{r}'},
\end{equation}
with $\mathbf{r}'=z'\hat{\mathbf{z}}$. The factor $e^{+\jj k_0\hat{\mathbf{r}}\cdot\mathbf{r}'}$ is angular \emph{phase} in the radiation integral; net angular-momentum transport requires the stress or flux integrals of later sections, not phase alone \cite{stratton1941,jackson1999}.
\end{example}

Figure~\ref{fig:ch1.green-function-seed} shows the geometry: source point $\mathbf{r}'$, observation point $\mathbf{r}$, separation $R$, and outgoing phase $e^{-\jj k_0 R}/(4\pi R)$ \cite{stratton1941,jackson1999,tai1994}. The same retarded phase underlies the multipole hierarchy of Subsection~\ref{sec:ch01.6.4} \cite{jackson1999,stratton1941}. Chapter~\ref{ch:02} develops the dyadic Green function under the same $\exp(+\jj\omega t)$ convention \cite{harrington2001,tai1994,stratton1941}.

\begin{figure}[t]
\centering
\VolIPhysFig{\VolIPhysGreenFunctionGeometry}
\caption{Localized impressed current at $\mathbf{r}'$, observation point $\mathbf{r}$, separation $R=|\mathbf{r}-\mathbf{r}'|$, observation direction $\hat{\mathbf{r}}$, and outgoing Green-function phase $e^{-\jj k_0 R}/(4\pi R)$; in the far zone $R\approx r-\hat{\mathbf{r}}\cdot\mathbf{r}'$ generates angular phase from source geometry \cite{stratton1941,jackson1999,tai1994}.}
\label{fig:ch1.green-function-seed}
\end{figure}


\begin{figure}[t]
\centering
\VolIPhysFig{\VolIPhysDipoleLoop}
\caption{Azimuthal filament loop of radius $b$ carrying $I_0\hat{\boldsymbol{\phi}}$: charge-continuous impressed current on compact support. Illustrative shells at $k_0 b=0.063$ (electrically small source) and $k_0 r=6.3$ ($r\approx\lambda_0$); the latter is an order-of-magnitude illustration, not by itself a proof of Fraunhofer or full far-zone asymptotics.}
\label{fig:ch1.wp-solenoidal-loop}
\end{figure}

Figure~\ref{fig:ch1.wp-solenoidal-loop} is the minimal charge-continuous source geometry: azimuthal
current on a filament loop produces dipole radiation because $\divg\Jimp=0$ on the support while
the impressed current closes the Amp\`ere row in Eq.~\eqref{eq:ch1.maxwell-local}. Replacing the
loop by a tangential $\E_t$ pattern on the same circle would satisfy a boundary display but not
the source row unless a matching $\Jimp$ is prescribed \cite{stratton1941,harrington2001}.

With solenoidal \(\Jimp\), Eq.~\eqref{eq:ch1.vacuum-wave-source-driven} yields a charge-continuous Maxwell field. Inserting a spurious \(\rho=\rho_0\delta(\mathbf{r})\) without adjusting \(\Jimp\) violates Eq.~\eqref{eq:ch1.charge-continuity-harmonic} and breaks the Amp\`ere--Gauss closure of Eqs.~\eqref{eq:ch1.maxwell-local} \cite{maxwell1873,stratton1941}.

\begin{example}[An electrically small current loop]
\label{ex:ch1.filament-loop}
A filament loop of radius \(b=3\,\mathrm{cm}\) carries phasor current \(I_0=1\,\mathrm{A}\) at \(f=100\,\mathrm{MHz}\). Then \(\lambda_0=c/f=3\,\mathrm{m}\), \(k_0=2\pi/\lambda_0\approx2.09\,\mathrm{rad\,m^{-1}}\), and \(k_0 b=0.063\ll1\): the loop is electrically small. The impressed azimuthal current is solenoidal on its support (\(\divg\Jimp=0\) in the distributional sense for a closed toroidal filament), so Eq.~\eqref{eq:ch1.vacuum-wave-source-driven} drives primarily magnetic-dipole radiation at asymptotic distances (Subsection~\ref{sec:ch01.2.4}) \cite{stratton1941,harrington2001}. At \(k_0 b=0.063\) reactive near fields dominate on shells comparable to \(b\). The value \(k_0 r=6.3\) corresponds to \(r=\lambda_0\) and is used \emph{illustratively} in Figure~\ref{fig:ch1.wp-solenoidal-loop}: it indicates order unity in \(k_0 r\) but does not, without the separate conditions \(k_0 b^2/r\ll1\) and established multipole dominance, replace a full far-zone proof. Fraunhofer-type phase accuracy for this loop would require \(k_0 b^2/r\ll1\), which holds easily at \(r=\lambda_0\) because \(b\ll\lambda_0\) \cite{stratton1941,jackson1999}.
\end{example}

Dispersive media require
frequency-dependent \(\varepsilon(\omega)\) and \(\mu(\omega)\), breaking the
algebraic completion of Eq.~\eqref{eq:ch1.vacuum-constitutive}. Relativistic or
quantum corrections lie outside the present harmonic class. Within the vacuum harmonic class,
Eqs.~\eqref{eq:ch1.maxwell-local}--\eqref{eq:ch1.local-plane-wave-relations}
form the complete local law set every later angular argument inherits
\cite{maxwell1873}.

\subsection{Sources, boundary conditions, and the Maxwell boundary-value problem}
\label{sec:ch01.1.2}

Free-space theory often begins with a localized Green-function representation that maps impressed sources to outgoing fields; an actual device or scatterer is specified, however, by the full Maxwell boundary-value problem on a domain $\Omega$ with boundary $\Gamma=\partial\Omega$. That statement names piecewise constitutive laws in subregions $\{\Omega_a\}$, impressed volume and surface sources, transmission conditions on internal interfaces, data on the outer boundary, and---when the exterior is open---an outgoing radiation selection \cite{stratton1941,maxwell1873}. Imposing a tangential field on an aperture without a charge-compatible $\Jimp$ and without compatible traces on $\Gamma$ need not correspond to any physically realizable Maxwell state \cite{harrington2001,chu1948}.

A fed antenna or inhomogeneous scatterer is specified by more than a volume current. At angular frequency $\omega>0$ under $\exp(+\jj\omega t)$, $(\Jimp,\rho)$ are impressed volume sources (A\,m$^{-2}$ and C\,m$^{-3}$); $\mathbf{J}_{s,ab}$ is impressed electric surface current on $\Sigma_{ab}=\partial\Omega_a\cap\partial\Omega_b$ (A\,m$^{-1}$); and $\rho_{s,ab}$ is free surface charge (C\,m$^{-2}$). In each subregion $\Omega_a$, $\mathbf{D}=\varepsilon_a\E$ and $\mathbf{B}=\mu_a\H$ with piecewise constant $\varepsilon_a,\mu_a$ unless stated otherwise \cite{stratton1941}. The finite-energy weighting of Subsection~\ref{sec:ch01.0.3} (Equation~\eqref{eq:ch1.energy-class}) tests admissibility on bounded domains; it is not the same statement as square-integrable energy on an unbounded exterior. Eq.~\eqref{eq:ch1.global-charge-compatibility} holds when $\Jimp$ admits a classical outward normal trace on $\Gamma$; for filamentary or distributional feeds the same balance is stated in weak form without a general trace theorem here \cite{stratton1941}.

\paragraph{Interface transmission.}
On each $\Sigma_{ab}$, let $\hat{\mathbf{n}}_{ab}$ be the unit normal directed from $\Omega_a$ to $\Omega_b$. Assume fields are smooth up to the interface except for the allowed jumps and that no impressed magnetic sheet current is present. Faraday's integral in Equation~\eqref{eq:ch1.maxwell-integral}, applied to a rectangular loop of area $\delta h$ straddling $\Sigma_{ab}$ with sides parallel and normal to the interface, yields $\oint \E\cdot d\boldsymbol{\ell}=-\jj\omega\mu_0\int_S \H\cdot\hat{\mathbf{n}}\,dS$. The right-hand side is $\mathcal{O}(\delta h)$, while the tangential contribution to the line integral is $[\hat{\mathbf{n}}_{ab}\times(\E_b-\E_a)]\cdot\Delta\boldsymbol{\ell}_t$ of order unity as the loop height $\delta\to 0$; and therefore $\hat{\mathbf{n}}_{ab}\times(\E_b-\E_a)=\mathbf{0}$ \cite{maxwell1873,jackson1999}. Amp\`ere--Maxwell on the same loop yields $\hat{\mathbf{n}}_{ab}\times(\H_b-\H_a)=\mathbf{J}_{s,ab}$ in the limit. Gauss's laws applied to a pillbox of height $\delta$ give $\hat{\mathbf{n}}_{ab}\cdot(\mathbf{D}_b-\mathbf{D}_a)=\rho_{s,ab}$ and, without magnetic monopoles, $\hat{\mathbf{n}}_{ab}\cdot(\mathbf{B}_b-\mathbf{B}_a)=0$ \cite{stratton1941}. Collecting,
\begin{equation}
\label{eq:ch1.interface-conditions}
\begin{aligned}
\hat{\mathbf{n}}_{ab}\times(\E_b-\E_a)&=\mathbf{0},\\
\hat{\mathbf{n}}_{ab}\times(\H_b-\H_a)&=\mathbf{J}_{s,ab},\\
\hat{\mathbf{n}}_{ab}\cdot(\mathbf{D}_b-\mathbf{D}_a)&=\rho_{s,ab},\\
\hat{\mathbf{n}}_{ab}\cdot(\mathbf{B}_b-\mathbf{B}_a)&=0.
\end{aligned}
\end{equation}
Equation~\eqref{eq:ch1.interface-conditions} is the complete transmission law used below; omitting $\rho_{s,ab}$ or $\mathbf{J}_{s,ab}$ when a coated or fed interface carries them misstates reflected power and source closure \cite{collin2001,jackson1999}.

\paragraph{Outer boundary traces.}
We partition $\Gamma=\Gamma_{\mathrm{PEC}}\cup\Gamma_{\mathrm{PMC}}\cup\Gamma_Z$. On every boundary-value problem, $\hat{\mathbf{n}}$ denotes the unit normal \emph{outward from the field domain} $\Omega$ on which Maxwell's equations are solved \cite{stratton1941}. For a bounded interior cavity, $\hat{\mathbf{n}}$ points from the cavity into the surrounding medium or exterior. For an exterior problem $\Omega_{\mathrm{ext}}$ (Subsection~\ref{sec:ch01.1.3}), the obstacle boundary is the inner surface of $\Omega_{\mathrm{ext}}$ and
\begin{equation}
\label{eq:ch1.exterior-normal-convention}
\hat{\mathbf{n}}_{\mathrm{ext}}=\hat{\mathbf{n}}\quad\text{on }\Gamma=\partial\mathcal{O},\qquad
\hat{\mathbf{n}}_{\mathrm{body}}=-\hat{\mathbf{n}}_{\mathrm{ext}},
\end{equation}
with $\hat{\mathbf{n}}_{\mathrm{body}}$ outward from the obstacle material into the exterior vacuum \cite{stratton1941,harrington2001}. A PEC boundary is an ideal limit of large tangential $\E$ screening: $\hat{\mathbf{n}}\times\E=\mathbf{0}$. A PMC boundary is the dual magnetic limit: $\hat{\mathbf{n}}\times\H=\mathbf{0}$ \cite{stratton1941}. Neither is universally interchangeable with a finite conductor model; both exclude penetration physics retained in richer models \cite{stratton1941}.

Passive surface impedance is stated on the \emph{obstacle} normal because the Leontovich relation ties tangential $\E$ and $\H$ on the material side of an absorbing sheet \cite{stratton1941,harrington2001}. Let $\mathbf{H}_t:=(\hat{\mathbf{n}}_{\mathrm{body}}\times\H)\times\hat{\mathbf{n}}_{\mathrm{body}}$ and $\mathbf{E}_t:=\hat{\mathbf{n}}_{\mathrm{body}}\times(\E\times\hat{\mathbf{n}}_{\mathrm{body}})$. The local scalar passive class used here is
\begin{equation}
\label{eq:ch1.outer-boundary-classes}
\begin{aligned}
\hat{\mathbf{n}}\times\E&=\mathbf{0}\quad &&\text{on }\Gamma_{\mathrm{PEC}},\\
\hat{\mathbf{n}}\times\H&=\mathbf{0}\quad &&\text{on }\Gamma_{\mathrm{PMC}},\\
\hat{\mathbf{n}}_{\mathrm{body}}\times\E&=-Z_s\,\mathbf{H}_t
\quad &&\text{on }\Gamma_Z,
\end{aligned}
\end{equation}
with scalar passive impedance $Z_s(\mathbf{r})$ specified on $\Gamma_Z$ (spatially constant on the sheet unless noted) and $\Re(Z_s)>0$ for absorption \cite{stratton1941,harrington2001}. Under the obstacle normal $\hat{\mathbf{n}}_{\mathrm{body}}$ (outward from the solid into the exterior vacuum), positive flux $\tfrac12\Re\{\hat{\mathbf{n}}_{\mathrm{body}}\cdot(\E\times\H^*)\}$ is power crossing the surface in the direction of $\hat{\mathbf{n}}_{\mathrm{body}}$, i.e.\ \emph{from the obstacle into the exterior field}, not into the obstacle \cite{stratton1941,harrington2001}. The minus sign in Equation~\eqref{eq:ch1.outer-boundary-classes} aligns the Leontovich relation with absorption: define absorbed power density
\begin{equation}
\label{eq:ch1.impedance-boundary-power}
P_{\mathrm{abs}}
=-\frac{1}{2}\Re\!\left\{
\hat{\mathbf{n}}_{\mathrm{body}}\cdot\left(\E\times\H^*\right)
\right\}
=\frac{1}{2}\Re\!\left\{
Z_s\left|\mathbf{H}_t\right|^2
\right\}\ge 0,
\end{equation}
by Equation~\eqref{eq:ch1.outer-boundary-classes} when $\Re(Z_s)>0$ \cite{stratton1941,harrington2001}. Because $\hat{\mathbf{n}}_{\mathrm{ext}}=-\hat{\mathbf{n}}_{\mathrm{body}}$,
\begin{equation}
\label{eq:ch1.impedance-exterior-normal-flux}
\Re\!\left\{
\hat{\mathbf{n}}_{\mathrm{ext}}\cdot\left(\E\times\H^*\right)
\right\}
=+\,P_{\mathrm{abs}}
=\frac{1}{2}\Re\!\left\{
Z_s\left|\mathbf{H}_t\right|^2
\right\}\ge 0
\end{equation}
on $\Gamma_Z$: power absorbed by the sheet is carried \emph{out of} the exterior field domain across the inner boundary \cite{stratton1941,harrington2001}. The passive sign in Equation~\eqref{eq:ch1.outer-boundary-classes} is fixed here by the local complex-Poynting identity and is checked independently for a uniform plane wave on $\Gamma_Z$. Existence and uniqueness for the full exterior impedance class are cited conditionally from Colton--Kress--Hsiao after convention conversion \cite{colton1981impedance}. All outer classes are collected by the boundary operator
\begin{equation}
\label{eq:ch1.boundary-data-operator}
\mathcal{B}_\Gamma[\E,\H]=\mathbf{g}_\Gamma,
\end{equation}
with prescribed data $\mathbf{g}_\Gamma$ on $\Gamma$ \cite{stratton1941}.
\phantomsection\label{eq:ch1.data-specification-tuple}

\paragraph{Outgoing selection.}
For exterior antenna and scattering problems, fields at large radius $R=|\mathbf{r}|$ must be asymptotically outgoing under $\exp(+\jj\omega t)$:
\begin{equation}
\label{eq:ch1.radiation-trace-condition}
\hat{\mathbf{n}}\times\H
+\eta_0^{-1}(\hat{\mathbf{n}}\times\E)\times\hat{\mathbf{n}}
=\mathbf{o}(R^{-1}),\qquad R\to\infty,
\end{equation}
with $\eta_0=\sqrt{\mu_0/\varepsilon_0}$ and $\hat{\mathbf{n}}=\hat{\mathbf{r}}$ \cite{stratton1941,sommerfeld1949}. Equation~\eqref{eq:ch1.radiation-trace-condition} excludes the incoming branch proportional to $e^{+\jj k_0 R}/R$ and matches $\E_{\mathrm{rad}}\sim \E_0(\theta,\phi)\,e^{-\jj k_0 R}/R$, $\H_{\mathrm{rad}}=\eta_0^{-1}\hat{\mathbf{r}}\times\E_{\mathrm{rad}}$. It is a selection rule on the exterior solution, not a proof of existence, uniqueness, finite global energy, or full trace regularity for later torque integrals \cite{harrington2001}. Absorbing layers at finite radius approximate the same physics on a truncated domain \cite{hansen1988}.

\paragraph{Source classes and global balance.}
Bounded excitation occupies $\Omega_s\subseteq\Omega$ and must satisfy Equation~\eqref{eq:ch1.charge-continuity-harmonic}. The regular class is
\begin{equation}
\label{eq:ch1.source-charge-continuous}
\Jimp\in L^2(\Omega_s)^3,\qquad
\operatorname{supp}(\Jimp)\subseteq\overline{\Omega_s},\qquad
\divg\Jimp+\jj\omega\rho=0
\ \text{in }\mathcal{D}'(\Omega).
\end{equation}
Square-integrable $\Jimp$ with distributional continuity does not automatically imply finite exterior field energy, classical normal traces of $\Jimp$ on $\Gamma$, or finite quadratic surface fluxes; those require separate hypotheses tied to the domain and outgoing selection \cite{stratton1941}. Closed filaments and segmented dipoles in examples are \emph{distributional} limits with $\divg\Jimp+\jj\omega\rho=0$ in $\mathcal{D}'(\Omega)$, not members of $L^2$ unless regularized by finite width \cite{stratton1941,jackson1999}. Integrating Equation~\eqref{eq:ch1.charge-continuity-harmonic} over $\Omega$ and applying the divergence theorem yields, \emph{when $\Jimp$ admits a classical outward normal trace on $\Gamma$},
\begin{equation}
\label{eq:ch1.global-charge-compatibility}
\int_{\Omega}\!\left(\divg\Jimp+\jj\omega\rho\right)\,dV
=\int_{\Gamma}\!\Jimp\cdot\hat{\mathbf{n}}\,dS+\jj\omega\int_{\Omega}\rho\,dV=0.
\end{equation}
The same argument shows that if $\Jimp$ is supported strictly inside $\Omega$, the surface term vanishes; otherwise impressed current crossing $\Gamma$ must be specified. A general distributional trace theorem for $\Jimp|_\Gamma$ is not established in this manuscript and remains an open extension \cite{maxwell1873,stratton1941}.

\phantomsection\label{prop:ch1.source-boundary-compatibility}
\begin{example}[Reflection at a lossless dielectric interface]
\label{ex:ch1.dielectric-interface}
Let $\Omega_1$ be vacuum ($z>0$) and $\Omega_2$ a semi-infinite dielectric with $\varepsilon_r=2.25$, $\mu_r=1$ occupying $z<0$, with interface at $z=0$ and $\hat{\mathbf{n}}_{12}=-\hat{\mathbf{z}}$ directed from $\Omega_1$ into $\Omega_2$, consistent with the definition of $\hat{\mathbf{n}}_{ab}$ in Subsection~\ref{sec:ch01.1.2}. A normally incident plane wave in $\Omega_1$ travels toward the interface from $z>0$ with $\E=E_0 e^{+\jj k_0 z}\hat{\mathbf{x}}$ under $\exp(+\jj\omega t)$. Eq.~\eqref{eq:ch1.interface-conditions} with $\mathbf{J}_{s,12}=\mathbf{0}$ gives continuous tangential $\E$ and $\H$. With intrinsic impedances $\eta_1=\eta_0$ and $\eta_2=\eta_0/\sqrt{\varepsilon_r}$, the reflection coefficient is $\Gamma=(\eta_2-\eta_1)/(\eta_2+\eta_1)=(1-\sqrt{\varepsilon_r})/(1+\sqrt{\varepsilon_r})=-0.2$, so $R=|\Gamma|^2=0.04$ and $T=1-R=0.96$ for a lossless interface; the transmitted field in $\Omega_2$ carries wavenumber $k_2=k_0\sqrt{\varepsilon_r}$ and amplitude $|E_t|=(1+\Gamma)|E_0|=0.8|E_0|$ \cite{stratton1941,jackson1999}.
\end{example}

\begin{figure}[t]
\centering
\VolIPhysFig{\csname VolIPhysWaveguideTE10\endcsname}
\caption{Rectangular waveguide cross section $a\times b$ with PEC sidewalls, TE$_{10}$ transverse $\E$ profile, and charge-compatible impressed-current feed at the aperture.}
\label{fig:ch1.wp-te10-waveguide}
\end{figure}
\begin{example}[TE$_{10}$ mode in a closed rectangular guide]
\label{ex:ch1.te10-waveguide}
A hollow PEC guide with cross section $a\times b$ supports $\E=E_0\sin(\pi x/a)\,e^{-\jj\beta z}\hat{\mathbf{y}}$ with angular cutoff $\omega_c=\pi c/a$ (rad\,s$^{-1}$) and propagation constant $\beta=\sqrt{k_0^2-(\pi/a)^2}$ when $k_0>\pi/a$. The corresponding cutoff \emph{frequency} is $f_c=c/(2a)$ (hertz), not $\omega_c/(2\pi)$ mislabelled as a separate symbol \cite{stratton1941,collin2001}. Tangential $\E$ vanishes on $x=0,a$ and $y=0,b$, implementing Eq.~\eqref{eq:ch1.outer-boundary-classes} on $\Gamma_{\mathrm{PEC}}$. This is a \emph{closed} boundary-value problem: Eq.~\eqref{eq:ch1.radiation-trace-condition} at infinity is not part of the closed guide eigenproblem, though an open-ended or fed guide coupled to free space requires a different exterior completion \cite{stratton1941,collin2001}. At $f=10\,\mathrm{GHz}$ with $a=2.3\,\mathrm{cm}$, $f_c\approx 6.52\,\mathrm{GHz}$ and $\beta/k_0\approx 0.759$; below cutoff the same geometry is evanescent and storage-dominated \cite{harrington2001}.
\end{example}

Compatible $(\Jimp,\rho)$, constitutive data, Equation~\eqref{eq:ch1.interface-conditions}, outer data Eq.~\eqref{eq:ch1.boundary-data-operator}, and outgoing selection Eq.~\eqref{eq:ch1.radiation-trace-condition} where applicable define the class of Maxwell fields under consideration. Existence of a solution in that class and uniqueness under additional hypotheses are treated in Subsection~\ref{sec:ch01.1.3}; flux, power, and angular-momentum integrals presuppose this closure \cite{maxwell1873,stratton1941}. Piecewise smooth interfaces and linear local constitutive laws are assumed; rough or nonlinear boundaries require enriched models before the same conservation interpretations apply \cite{stratton1941}.

\subsection{Boundary-Value Problems and Uniqueness}
\label{sec:ch01.1.3}

Surface integrals of power, momentum, and angular momentum name observables only when the underlying field is uniquely determined. Two distinct Maxwell solutions that share the same sources, interface data, boundary traces, and---on an open exterior---the same outgoing selection would make every such integral ambiguous \cite{stratton1941,sommerfeld1949}. Uniqueness is the condition that turns a completed boundary-value specification into one physical field.

Suppose $(\E_1,\H_1)$ and $(\E_2,\H_2)$ both satisfy the same data. The difference $(\Delta\E,\Delta\H)=(\E_1-\E_2,\H_1-\H_2)$ carries no impressed current ($\Delta\Jimp=\mathbf{0}$) and obeys the homogeneous boundary operator $\mathcal{B}_\Gamma[\Delta\E,\Delta\H]=\mathbf{0}$ on every prescribed trace \cite{stratton1941}. Under $\exp(+\jj\omega t)$ the difference satisfies the same curl system with the domain constitutive $\varepsilon$ and $\mu$. The central question is whether a nontrivial source-free difference can persist under the stated boundary and radiation hypotheses.

The energy route begins with a local quadratic identity. The vector relation $\divg(\Delta\E\times\Delta\H^*)=\Delta\H^*\cdot\curl\Delta\E-\Delta\E\cdot\curl\Delta\H^*$, together with $\curl\Delta\E=-\jj\omega\mu\Delta\H$, $\curl\Delta\H=\jj\omega\varepsilon\Delta\E$, and $\curl\Delta\H^*=-\jj\omega\varepsilon^*\Delta\E^*$ when $\varepsilon$ is complex, yields the local relation
\begin{equation}
\label{eq:ch1.uniqueness-local-div}
\divg\mathbf{S}_\Delta=-2\jj\omega\left(u_{\Delta H}-u_{\Delta E}\right)
\end{equation}
for $\mathbf{S}_\Delta=\tfrac12\Delta\E\times\Delta\H^*$, with $u_{\Delta E}=\tfrac14\varepsilon|\Delta\E|^{2}$ and $u_{\Delta H}=\tfrac14\mu|\Delta\H|^{2}$ \cite{stratton1941}. The factor $2$ matches Equation~\eqref{eq:ch1.uniqueness-energy-identity}: a half-sized local divergence would not reproduce the integrated identity. Integrating over $\Omega$ and applying the divergence theorem yields
\begin{equation}
\label{eq:ch1.uniqueness-energy-identity}
\int_{\Gamma}\hat{\mathbf{n}}\cdot\mathbf{S}_\Delta\,dS
=-2\jj\omega\!\int_{\Omega}\!\left(u_{\Delta H}-u_{\Delta E}\right)dV
\end{equation}
when $\Delta\E,\Delta\H$ are smooth enough for the classical divergence theorem on $\Omega$ \cite{stratton1941}. Equation~\eqref{eq:ch1.uniqueness-energy-identity} does not in isolation prove existence or global regularity; it constrains the difference field.

Three conditions govern the arguments below. The difference field is source-free. Dissipation, when present, is passive: isotropic loss with $\varepsilon'',\mu''\ge0$ and/or passive scalar impedance with $\Re(Z_s)>0$ on prescribed sheets. On open exteriors, the completion must carry outgoing selection Eq.~\eqref{eq:ch1.radiation-trace-condition} when flux is read at infinity. None of these items proves existence; each limits what a nontrivial $(\Delta\E,\Delta\H)$ may do under the difference Poynting balance.

For real $\varepsilon$ and $\mu$ on $\Omega$, Equation~\eqref{eq:ch1.uniqueness-energy-identity} has a purely imaginary right-hand side, so its real part gives $\Re\{\int_{\Gamma}\hat{\mathbf{n}}\cdot\mathbf{S}_\Delta\,dS\}=0$ for every difference field. A lossless interior therefore does not by itself forbid nontrivial source-free modes at cavity eigenfrequencies \cite{stratton1941}. When dissipation enters through complex isotropic $\varepsilon,\mu$ on the source-free difference problem, the same identity yields
\begin{equation}
\label{eq:ch1.uniqueness-local-div-complex}
\nabla\cdot\mathbf{S}_\Delta
=\frac{\jj\omega}{2}\left(\varepsilon^*|\Delta\E|^{2}-\mu|\Delta\H|^{2}\right),
\end{equation}
which reduces to Equation~\eqref{eq:ch1.uniqueness-local-div} when $\varepsilon$ and $\mu$ are real \cite{stratton1941,harrington2001}. When dissipation is included through $\varepsilon=\varepsilon'-\jj\varepsilon''$ and/or $\mu=\mu'-\jj\mu''$ with $\varepsilon'',\mu''\ge0$ (passive loss under $\exp(+\jj\omega t)$),
\[
\Re\{\nabla\cdot\mathbf{S}_\Delta\}
=-\frac{\omega}{2}\left(\varepsilon''|\Delta\E|^{2}+\mu''|\Delta\H|^{2}\right)
\]
for isotropic local $\varepsilon,\mu$ in the interior (volume divergence, not a surface flux density) \cite{stratton1941,harrington2001}. For later reference, integration over $\Omega$ and the divergence theorem yield
\[
\Re\!\left\{\int_{\Gamma}\hat{\mathbf{n}}\cdot\mathbf{S}_\Delta\,dS\right\}
=-\frac{\omega}{2}\int_{\Omega}\left(\varepsilon''|\Delta\E|^{2}+\mu''|\Delta\H|^{2}\right)dV
\le 0,
\]
with equality forcing $\Delta\E=\mathbf{0}$ on open subsets where $\varepsilon''>0$ and $\Delta\H=\mathbf{0}$ on open subsets where $\mu''>0$, subject to the integrated balance and boundary data \cite{stratton1941,harrington2001}. On an open region $\Omega_\varepsilon\subset\Omega$ with $\varepsilon''>0$ and finite nonsingular $\mu$, the volume integral forces $\Delta\E=\mathbf{0}$ almost everywhere on $\Omega_\varepsilon$; Faraday's law $\curl\Delta\E=-\jj\omega\mu\Delta\H$ then gives $\Delta\H=\mathbf{0}$ there in the distributional sense under the stated nondegeneracy of $\mu$ \cite{stratton1941,harrington2001}. The symmetric case with $\mu''>0$ on $\Omega_\mu$ uses the conjugated Amp\`ere--Maxwell row to annihilate $\Delta\H$ first and then $\Delta\E$ on $\Omega_\mu$. These are \emph{local} conclusions on dissipative subdomains; extending $\Delta\E=\Delta\H=\mathbf{0}$ from a proper dissipative subset to all of a connected domain still requires a separate unique-continuation argument unless passive dissipation covers the whole connected volume or the boundary data already exclude nontrivial modes \cite{stratton1941,harrington2001}.

On a bounded lossy domain with homogeneous PEC or PMC data on all of $\Gamma$, the real surface flux of $\mathbf{S}_\Delta$ vanishes. If $\varepsilon''>0$ throughout a connected $\Omega$ with finite nonsingular $\mu$, or $\mu''>0$ throughout with finite nonsingular $\varepsilon$, the local argument above yields $\Delta\E=\Delta\H=\mathbf{0}$ throughout that volume; when loss occupies only a proper subset of a connected domain, global vanishing on the whole domain is not automatic from the integral identity alone \cite{stratton1941}. Passive impedance sheets on $\Gamma_Z$ enter through Eqs.~\eqref{eq:ch1.outer-boundary-classes}--\eqref{eq:ch1.impedance-exterior-normal-flux}; for $\Re(Z_s)>0$, $P_{\mathrm{abs}}\ge0$ and $\Re\{\hat{\mathbf{n}}_{\mathrm{ext}}\cdot\mathbf{S}\}=P_{\mathrm{abs}}$ on $\Gamma_Z$ \cite{stratton1941,harrington2001}. The standard exception remains a \emph{lossless} closed cavity at an eigenfrequency, which may support a nontrivial homogeneous mode \cite{stratton1941,sommerfeld1949}.

Antennas and scatterers are posed on unbounded exteriors. Equation~\eqref{eq:ch1.uniqueness-energy-identity} must then be combined with outgoing selection Eq.~\eqref{eq:ch1.radiation-trace-condition}. That condition excludes incoming completion at infinity but does not, by itself, force a nontrivial difference field to vanish in the far zone \cite{sommerfeld1949,harrington2001}. The exterior class used below takes $\Omega_{\mathrm{ext}}$ as the complement of a bounded obstacle $\mathcal{O}$ with piecewise smooth boundary $\Gamma=\partial\mathcal{O}$, filled with homogeneous lossless vacuum ($\varepsilon_0,\mu_0$). Scalar passive impedance on $\Gamma$ is the third line of Eq.~\eqref{eq:ch1.outer-boundary-classes} with $\Re(Z_s)>0$, using $\hat{\mathbf{n}}_{\mathrm{body}}$ and Eq.~\eqref{eq:ch1.exterior-normal-convention}. For a source-free difference field $(\Delta\E,\Delta\H)$ in $\Omega_{\mathrm{ext}}$, Eqs.~\eqref{eq:ch1.impedance-boundary-power} and~\eqref{eq:ch1.impedance-exterior-normal-flux} give
\[
P_{\mathrm{abs},\Delta}
=\tfrac12\Re(Z_s)\,|\Delta\H_t|^{2}\ge 0,
\qquad
\Re\!\left\{\hat{\mathbf{n}}_{\mathrm{ext}}\cdot\mathbf{S}_\Delta\right\}
=+P_{\mathrm{abs},\Delta}\ge 0
\quad\text{on }\Gamma.
\]
Integrating the lossless form of Equation~\eqref{eq:ch1.uniqueness-local-div} over a truncated exterior domain $\Omega_{\mathrm{ext},R}$ bounded by $\Gamma$ and a large sphere $S_R$ and taking the real part of Eq.~\eqref{eq:ch1.uniqueness-energy-identity} with outward normal $\hat{\mathbf{n}}_{\mathrm{ext}}$ on $\Gamma$ yields, because the volume is source-free and lossless,
\begin{equation}
\label{eq:ch1.exterior-difference-energy-balance}
\Re\!\left\{\int_{S_R}\hat{\mathbf{r}}\cdot\mathbf{S}_\Delta\,dS\right\}
+\Re\!\left\{\int_{\Gamma}\hat{\mathbf{n}}_{\mathrm{ext}}\cdot\mathbf{S}_\Delta\,dS\right\}=0.
\end{equation}
On $\Gamma_Z$, Eqs.~\eqref{eq:ch1.impedance-boundary-power} and~\eqref{eq:ch1.impedance-exterior-normal-flux} applied to $(\Delta\E,\Delta\H)$ give
\[
\Re\!\left\{\hat{\mathbf{n}}_{\mathrm{ext}}\cdot\mathbf{S}_\Delta\right\}
=+P_{\mathrm{abs},\Delta}
=\tfrac12\Re(Z_s)\,|\Delta\H_t|^{2}\ge 0
\quad\text{on }\Gamma_Z,
\]
so the inner boundary contribution in Equation~\eqref{eq:ch1.exterior-difference-energy-balance} is \emph{nonnegative}. On a sphere at finite radius $R$, $\Re\{\hat{\mathbf{r}}\cdot\mathbf{S}_\Delta\}$ need not be pointwise nonnegative everywhere: reactive structure in the difference field can carry imaginary power through the shell. One therefore passes to the limit $R\to\infty$ before interpreting radiated power. With outgoing selection Eq.~\eqref{eq:ch1.radiation-trace-condition}, define the far-field pattern
\[
\E_\infty(\theta,\phi)=\lim_{R\to\infty}R\,e^{+\jj k_0 R}\,\E_\Delta(R,\theta,\phi)
\]
and the limiting real outward power
\[
\mathcal{P}_\infty=\lim_{R\to\infty}\Re\!\left\{\int_{S_R}\hat{\mathbf{r}}\cdot\mathbf{S}_\Delta\,dS\right\}
=\frac{1}{2\eta_0}\int_{4\pi}|\E_{\infty,t}|^{2}\,d\Omega\ge 0.
\]
Passing to the limit in Equation~\eqref{eq:ch1.exterior-difference-energy-balance} yields
\[
\mathcal{P}_\infty+\frac{1}{2}\Re(Z_s)\int_{\Gamma_Z}|\Delta\H_t|^{2}\,dS=0
\]
for scalar passive $Z_s$ with $\Re Z_s>0$ on $\Gamma_Z$. Hence $\E_\infty=\mathbf{0}$ and $\Delta\H_t=\mathbf{0}$ on $\Gamma_Z$, whence $\Delta\E_t=\mathbf{0}$ on $\Gamma_Z$ by Equation~\eqref{eq:ch1.outer-boundary-classes}. The continuation is stated in four explicit steps. \emph{(i)}~Vanishing far-field pattern: $\E_\infty=\mathbf{0}$. \emph{(ii)}~In source-free homogeneous vacuum, $\Delta\E$ admits the outgoing vector spherical-wave representation of Stratton and Harrington \cite{stratton1941,harrington2001}:
\[
\Delta\E
=\frac{k_0}{\eta_0}\sum_{\ell,m}\left(a_{\ell m}\,\mathbf M_{\ell m}^{+}+b_{\ell m}\,\mathbf N_{\ell m}^{+}\right),
\qquad
\Delta\H
=-\frac{\jj}{\eta_0}\sum_{\ell,m}\left(a_{\ell m}\,\mathbf N_{\ell m}^{+}-b_{\ell m}\,\mathbf M_{\ell m}^{+}\right),
\]
with $\mathbf M_{\ell m}^{+}=h_\ell^{(2)}(k_0 r)\,\vec{\mathcal{M}}_{\ell m}(\theta,\phi)$ and $\mathbf N_{\ell m}^{+}=(1/k_0)\nabla\times\mathbf M_{\ell m}^{+}$ \cite{stratton1941,harrington2001}. Each member of these Stratton--Harrington families is transverse: $\divg\mathbf M_{\ell m}^{+}=\mathbf{0}$ and $\divg\mathbf N_{\ell m}^{+}=\mathbf{0}$ in the source-free exterior, so $\divg\Delta\E=0$ is satisfied term-by-term and does \emph{not} impose an additional algebraic relation between the independent coefficients $(a_{\ell m},b_{\ell m})$ for fixed $(\ell,m)$ \cite{stratton1941,harrington2001}. Outgoing selection Eq.~\eqref{eq:ch1.radiation-trace-condition} singles out $h_\ell^{(2)}$ at infinity. The far-field pattern $\E_\infty$ extracts the tangential $r^{-1}$ coefficients of both families; $\E_\infty=\mathbf{0}$ therefore annihilates all outgoing $(a_{\ell m},b_{\ell m})$, so $\Delta\E_t=\mathbf{0}$ on any sphere $S_R$ that encloses $\mathcal{O}$ at radius $R$ large enough for the series to represent the difference field. \emph{(iii)}~By the Rellich unique-continuation property for outgoing source-free Maxwell fields in a connected exterior domain---with finite local energy on bounded subdomains, piecewise smooth obstacle regularity, and the Silver--M\"uller radiation condition consistent with Eq.~\eqref{eq:ch1.radiation-trace-condition}---vanishing tangential $\Delta\E$ on such an $S_R$ forces $\Delta\E=\mathbf{0}$ throughout the source-free shell exterior to $\mathcal{O}$ and up to the obstacle boundary $\Gamma$ \cite{sommerfeld1949,harrington2001}. \emph{(iv)}~The impedance traces enforced in step~(i) then yield $\Delta\H=\mathbf{0}$ on $\Gamma$ and, by Faraday's law with $\Delta\Jimp=\mathbf{0}$ in $\Omega_{\mathrm{ext}}$, $\Delta\H=\mathbf{0}$ everywhere in $\Omega_{\mathrm{ext}}$ \cite{stratton1941,harrington2001}. This argument closes \emph{homogeneous} uniqueness for the scalar passive class with $\hat{\mathbf{n}}_{\mathrm{ext}}=-\hat{\mathbf{n}}_{\mathrm{body}}$; it does not establish existence or continuous dependence for the inhomogeneous scattering problem, nor equation-for-equation agreement with Colton--Kress--Hsiao without inspection of that primary theorem \cite{colton1981impedance}. Nonlocal, anisotropic, or higher-order impedance operators lie outside the scalar passive class treated here \cite{stratton1941}.
\phantomsection\label{prop:ch1.uniqueness-passive-closure}
\phantomsection\label{eq:ch1.uniqueness-condition-index}

\subsection{Energy, linear momentum, and Maxwell stress tensor}
\label{sec:ch01.1.4}

Energy transport and mechanical force are not independent postulates bolted onto Maxwell's equations. On a declared control volume, the curl--divergence structure of $(\E,\H)$ yields the Poynting identity, the Lorentz force density, and the symmetric Maxwell stress tensor from which linear momentum and, in Subsection~\ref{sec:ch01.1.5}, angular momentum follow \cite{poynting1884,maxwell1873}. This subsection develops the instantaneous and harmonic energy balances; the momentum balances displayed here prepare the stress-moment integrals used later for torque. Let $V$ be a control volume with outward unit normal $\hat{\mathbf{n}}$ on $\partial V$; let $V_m$ enclose the matter whose force or torque is evaluated; let $\Gamma_R$ be a predominantly radiative shell when radiated power or flux is read \cite{stratton1941}. Finite surface integrals require the regularity hypotheses of Subsection~\ref{sec:ch01.0.3} on the same volume---Eq.~\eqref{eq:ch1.energy-class} alone does not supply every trace needed later.

Start with real-time fields in vacuum. Divergence of $\E\times\H$ splits into the curls of $\E$ and $\H$; substituting Maxwell's equations converts that product rule into a balance between impressed work, stored energy, and Poynting outflow:
\begin{equation}
\label{eq:ch1.poynting-instantaneous}
\divg(\E\times\H)
=-\Jimp\cdot\E
-\frac{\partial}{\partial t}\!\left(\frac{\varepsilon_0}{2}|\E|^2+\frac{\mu_0}{2}|\H|^2\right),
\end{equation}
where $\Jimp$ is impressed current density \cite{stratton1941}. We define instantaneous
Poynting vector $\mathbf{s}=\E\times\H$ and energy density
$u=\tfrac12(\varepsilon_0|\E|^2+\mu_0|\H|^2)$. It follows that
$\divg\mathbf{s}+\partial u/\partial t+\Jimp\cdot\E=0$. The term $\Jimp\cdot\E$ is
power density at which the electromagnetic field performs work on impressed current per unit volume; $\divg\mathbf{s}$ is net energy
outflow from an infinitesimal volume. Integrating over volume $V$ and applying the
divergence theorem yields the power balance through $\partial V$ \cite{poynting1884}.

Instantaneous transport uses real fields: $\mathbf{s}=\E\times\H$, $\mathbf{g}_{\mathrm{em}}=\varepsilon_0\,\E\times\B$, and $\Tinst$ in Eqs.~\eqref{eq:ch1.maxwell-stress-components}--\eqref{eq:ch1.momentum-balance-instantaneous} \cite{maxwell1873,poynting1884}. Harmonic phasors use $\mathbf{S}_\mathrm{c}=\tfrac12\E\times\H^*$, $\langle\mathbf{S}\rangle=\Re\mathbf{S}_\mathrm{c}$, and cycle-averaged stress $\Tavg$ of Equation~\eqref{eq:ch1.maxwell-stress} with $\Mavg=\mathbf{r}\times\Tavg$ \cite{stratton1941}. The two descriptions are linked by $\exp(+\jj\omega t)$ but must not be interchanged when reading reactive storage or radiative flux.

Harmonic steady state replaces time derivatives by $+\jj\omega$ on phasors. Cycle averages of products $\E(t)\cdot\H(t)$ then collapse to $\mathbf{S}_\mathrm{c}=\tfrac12\E\times\H^*$, $\langle\mathbf{S}\rangle=\Re\mathbf{S}_\mathrm{c}$, and the electric and magnetic storage terms $u_E$, $u_H$ already used above. Divergence of $\E\times\H^*$ is expanded with the product rule; substituting $\curl\E=-\jj\omega\mu_0\H$ and the conjugated Amp\`ere--Maxwell row yields Equation~\eqref{eq:ch1.poynting-complex} below \cite{stratton1941,harrington2001}.
With $\curl\E=-\jj\omega\mu_0\H$ and
$\curl\H=\Jimp+\jj\omega\varepsilon_0\E$, and $\curl\H^*=\Jimp^*-\jj\omega\varepsilon_0\E^*$,
\begin{equation}
\label{eq:ch1.poynting-intermediate}
\divg(\E\times\H^*)
=-\E\cdot\Jimp^*
+\jj\omega\bigl(\varepsilon_0\E\cdot\E^*-\mu_0\H^*\!\cdot\H\bigr).
\end{equation}
The storage difference satisfies
$\varepsilon_0\E\cdot\E^*-\mu_0\H^*\!\cdot\H=4(u_E-u_H)$ \cite{stratton1941}.
Dividing Equation~\eqref{eq:ch1.poynting-intermediate} by two then yields
\begin{equation}
\label{eq:ch1.poynting-average}\divg\mathbf{S}_\mathrm{c}
=-\frac{1}{2}\E\cdot\Jimp^*
-2\jj\omega\left(u_H-u_E\right).
\end{equation}
\addtocounter{equation}{-1}\refstepcounter{equation}\label{eq:ch1.dual-invariant-poynting}%
\addtocounter{equation}{-1}\refstepcounter{equation}\label{eq:ch1.poynting-complex}%
Physically, Equation~\eqref{eq:ch1.poynting-complex} states that net complex power leaving a differential volume equals electromagnetic work on impressed current plus harmonic exchange between electric and magnetic storage \cite{poynting1884}. The real part $\divg\Re(\mathbf{S}_\mathrm{c})$ balances impressed work; the imaginary part $\divg\Im(\mathbf{S}_\mathrm{c})$ tracks reactive exchange in the complex-power balance and is not, without further definition, a unique local stored-energy density \cite{stratton1941}. Momentum in material media involves additional stress--energy forms (Abraham, Minkowski); vacuum-side traction and angular flux use the Maxwell stress of Eq.~\eqref{eq:ch1.maxwell-stress} \cite{abraham1910,minkowski1908,stratton1941}.

\begin{equation}
\label{eq:ch1.poynting-real-imag-split}
\begin{aligned}
\divg\Re(\mathbf{S}_\mathrm{c})&=-\frac12\Re\!\left(\E\cdot\Jimp^*\right),\\
\divg\Im(\mathbf{S}_\mathrm{c})&=-\frac12\Im\!\left(\E\cdot\Jimp^*\right)-2\omega(u_H-u_E).
\end{aligned}
\end{equation}
\begin{equation}
\label{eq:ch1.dual-invariant-energy}\int_{\partial V}\mathbf{S}_\mathrm{c}\cdot\hat{\mathbf{n}}\,dS
=-\frac12\int_V\E\cdot\Jimp^*\,dV
-2\jj\omega\int_V(u_H-u_E)\,dV.
\end{equation}
\addtocounter{equation}{-1}\refstepcounter{equation}\label{eq:ch1.power-balance-integral}%
Let $V\subset\Omega$ be a control volume with outward normal $\hat{\mathbf{n}}$ on $\partial V$ \cite{stratton1941}. The surface integral is net complex power through $\partial V$; the volume terms are impressed work and reactive storage imbalance inside $V$ \cite{poynting1884}. Integrating the local balance over any finite region therefore relates power crossing the boundary to source work and stored reactive energy within the volume---the finite-volume form of Poynting's theorem \cite{stratton1941}.

Mechanical force on matter follows from the Lorentz law,
\begin{equation}
\label{eq:ch1.lorentz-force}
\mathbf{f}_{\mathrm{matter}}=\rho_{\mathrm{phys}}\E+\mathbf{J}_{\mathrm{phys}}\times\B,
\end{equation}
where $\rho_{\mathrm{phys}}$ and $\mathbf{J}_{\mathrm{phys}}$ are the charge and current densities of the carriers that experience the force (N\,m$^{-3}$) \cite{maxwell1873,stratton1941}.
In Eqs.~\eqref{eq:ch1.maxwell-local}, $\rho$ is the charge density in Gauss's law and $\Jimp$ is the impressed current in Amp\`ere's law---the mathematical source pair $(\rho,\Jimp)$ that closes Maxwell's equations \cite{maxwell1873,stratton1941}.
\textbf{Route~A.} When impressed sources model the actual charge carriers on their support, set $\rho_{\mathrm{phys}}=\rho$ and $\mathbf{J}_{\mathrm{phys}}=\Jimp$. The Lorentz force density $\mathbf{f}_{\mathrm{matter}}$ then coincides with the source-exchange density $\mathbf{f}_{\mathrm{src}}$ used in the momentum balance.

\textbf{Route~B.} When $(\rho,\Jimp)$ is a mathematical feed that does not equal the carrier densities on a load, $\mathbf{f}_{\mathrm{matter}}$ must be evaluated with the physical $(\rho_{\mathrm{phys}},\mathbf{J}_{\mathrm{phys}})$ of the material. Force and torque on a body in source-free vacuum are read from traction $\oint_{\partial V_b}\Tinst\cdot\hat{\mathbf{n}}\,dS$ on a closed surface in the exterior field. The term $-\tfrac12\E\cdot\Jimp^*$ in the complex Poynting identity records impressed-source exchange; it is not, by itself, mechanical work on an arbitrary body \cite{poynting1884,stratton1941,harrington2001}.
The symmetric vacuum stress tensor Eqs.~\eqref{eq:ch1.maxwell-stress-components}--\eqref{eq:ch1.maxwell-stress} and the local identity below are derived for regions where $(\E,\B)$ obey vacuum constitutive relations $\varepsilon_0,\mu_0$ and the displayed curl--divergence structure; they are not asserted pointwise inside arbitrary macroscopic media without an applicable material stress model \cite{stratton1941}. Using Maxwell's equations in such a vacuum/source region to eliminate the same $(\rho,\Jimp)$ as in Equation~\eqref{eq:ch1.maxwell-local}, in favour of $(\E,\B)$,
one forms the electromagnetic momentum density
\begin{equation}
\label{eq:ch1.momentum-density-instantaneous}
\mathbf{g}_{\mathrm{em}}=\varepsilon_0\,\E\times\B
=\mathbf{s}/c^2,
\end{equation}
where $\mathbf{s}=\E\times\H$ is the instantaneous Poynting vector \cite{maxwell1873,poynting1884}.
In this setting, radiation pressure on an absorber is momentum delivered through $\Tinst$ or, in harmonic form, through $\Tavg$ \cite{maxwell1873,stratton1941}. The instantaneous symmetric Maxwell stress in Cartesian components is
\begin{equation}
\label{eq:ch1.maxwell-stress-components}
T^{(i)}_{ij}=\varepsilon_0\!\left(E_iE_j-\tfrac12\delta_{ij}E^2\right)
+\frac{1}{\mu_0}\!\left(B_iB_j-\tfrac12\delta_{ij}B^2\right),
\end{equation}
with $E^2=\mathbf{E}\cdot\mathbf{E}$ and $B^2=\mathbf{B}\cdot\mathbf{B}$ \cite{maxwell1873,stratton1941}.
Index contraction with Gauss's and Amp\`ere's laws yields the field momentum balance
\begin{equation}
\label{eq:ch1.momentum-balance-instantaneous}
\frac{\partial g_{i}}{\partial t}+f_{\mathrm{src},i}=\partial_j T^{(i)}_{ij},
\end{equation}
equivalently $\partial_t\mathbf{g}_{\mathrm{em}}+\mathbf{f}_{\mathrm{src}}=\divg\Tinst$, where
\begin{equation}
\label{eq:ch1.maxwell-source-force}
\mathbf{f}_{\mathrm{src}}=\rho\E+\Jimp\times\B
\end{equation}
uses the same $\rho$ and $\Jimp$ as Eqs.~\eqref{eq:ch1.maxwell-local} \cite{maxwell1873,jackson1999}. With real fields $\E(\mathbf{r},t)=\Re\{\E(\mathbf{r})e^{+\jj\omega t}\}$ and similarly for $\H$, $\B$, $\Jimp$, and $\rho$, cycle averaging yields
$\langle\mathbf{f}_{\mathrm{src}}\rangle=\tfrac12\Re\{\rho\E^*+\Jimp\times\B^*\}$,
distinct from the instantaneous product $\mathbf{f}_{\mathrm{src}}=\rho\E+\Jimp\times\B$ \cite{stratton1941}. Equation~\eqref{eq:ch1.momentum-balance-instantaneous} is field--source bookkeeping; $\mathbf{f}_{\mathrm{matter}}$ equals $\mathbf{f}_{\mathrm{src}}$ only under Route~A on the relevant support.
Throughout the stress, Poynting, and torque balances of
Subsections~\ref{sec:ch01.1.4}--\ref{sec:ch01.1.5}, $V$ denotes a control volume
enclosing the matter whose mechanical force or torque is being evaluated, and
$\hat{\mathbf{n}}$ is the unit normal outward from $V$ \cite{stratton1941}.
Integrating over that $V$ with outward normal $\hat{\mathbf{n}}$,
\begin{equation}
\label{eq:ch1.momentum-integral-balance}
\mathbf{F}_{\mathrm{matter}}
=\int_V \mathbf{f}_{\mathrm{matter}}\,dV
=\int_{\partial V}\Tinst\cdot\hat{\mathbf{n}}\,dS
-\frac{d}{dt}\int_V \mathbf{g}_{\mathrm{em}}\,dV
\end{equation}
when $V$ lies in a vacuum/source region where Equation~\eqref{eq:ch1.momentum-balance-instantaneous} holds and Route~A applies throughout $V$, so $\mathbf{f}_{\mathrm{matter}}=\mathbf{f}_{\mathrm{src}}$ there \cite{stratton1941}. Under Route~B, $\int_V \mathbf{f}_{\mathrm{src}}\,dV$ is source-exchange bookkeeping on the chosen volume, while mechanical force on a body is $\mathbf{F}=\oint_{\partial V_b}\Tinst\cdot\hat{\mathbf{n}}\,dS$ (minus storage) for a closed $\partial V_b$ in source-free vacuum enclosing that body, with traction evaluated on the vacuum field side \cite{stratton1941,harrington2001}. Traction must not be identified with momentum density $\mathbf{g}_{\mathrm{em}}$ or with Poynting flux $\mathbf{S}_\mathrm{c}$.

\begin{figure}[t]
\centering
\VolIPhysFig{\VolIPhysStressTraction}
\caption{Differential surface element with outward normal $\hat{\mathbf{n}}$: cycle-averaged traction $\mathbf{t}=\Tavg\cdot\hat{\mathbf{n}}$ on a material-enclosing control volume $V_m$ \cite{maxwell1873,stratton1941}.}
\label{fig:ch1.stress-traction-element}
\end{figure}

On the same Maxwell pair, figure~\ref{fig:ch1.stress-traction-element} shows cycle-averaged traction $\mathbf{t}=\Tavg\cdot\hat{\mathbf{n}}$ on a material-enclosing volume $V_m$.

For harmonic fields with $\exp(+\jj\omega t)$, time averaging yields the cycle-averaged dyadic
\begin{equation}
\label{eq:ch1.maxwell-stress}
\Tavg
=\frac{1}{2}\Re\!\left[
\varepsilon_0\,\E\E^*
+\mu_0\,\H\H^*
-\frac{1}{2}\mathbf{I}
\left(\varepsilon_0\,\E\cdot\E^*+\mu_0\,\H\cdot\H^*\right)
\right],
\end{equation}
where \(\mathbf{I}\) is the identity dyad and juxtaposition such as \(\E\E^*\)
is dyadic product. Operand by operand, the first two dyads are directional field
intensity contributions, while the isotropic subtraction term
\(\tfrac{1}{2}\mathbf{I}(\cdots)\) removes trace pressure so that directional
tractions are represented correctly. The real-part operator keeps physically
observable time-averaged transfer \cite{stratton1941}.
Together, $\mathbf{S}_\mathrm{c}$ and $\langle\mathbf{S}\rangle$ give energetic flux densities and $\Tavg$ yields cycle-averaged mechanical traction; Subsection~\ref{sec:ch01.1.5} uses $\Mavg=\mathbf{r}\times\Tavg$ and stress-moment surface density $\Mavg\cdot\hat{\mathbf{n}}$ \cite{poynting1884}.

On a material-enclosing volume $V_m$, traction $\mathbf{t}=\Tavg\cdot\hat{\mathbf{n}}$ is cycle-averaged force per unit area on matter, and $\mathbf{F}=\int_{\partial V_m}\Tavg\cdot\hat{\mathbf{n}}\,dS$ \cite{stratton1941}.
For a plane wave in $+\hat{\mathbf{z}}$ incident on a perfectly absorbing sheet at
$z=0$ with vacuum control volume $z<0$, the outward unit normal on the interface is
$\hat{\mathbf{n}}=+\hat{\mathbf{z}}$. Radiation pressure on the absorber is
$p_{\mathrm{abs}}=|\langle\mathbf{S}\rangle|/c$. A perfect reflector at normal incidence receives
$p_{\mathrm{refl}}=2|\langle\mathbf{S}\rangle|/c$ \cite{stratton1941,maxwell1873}.

\begin{example}[Radiation pressure on an absorbing surface]
\label{ex:ch1.radiation-pressure}
With $\langle\mathbf{S}\rangle=(|E_0|^2/2\eta_0)\hat{\mathbf{z}}$ from Eq.~\eqref{eq:ch1.plane-wave-transverse}, $p_{\mathrm{abs}}=|E_0|^2/(2\eta_0 c)$. If $|E_0|=100\,\mathrm{V/m}$, then $\langle S_z\rangle=13.2\,\mathrm{W/m^2}$ and $p_{\mathrm{abs}}\approx4.4\times10^{-8}\,\mathrm{N/m^2}$; a perfect reflector doubles the pressure \cite{stratton1941}. The same result follows from the momentum balance: in vacuum $\langle\mathbf{g}_{\mathrm{em}}\rangle=\Re(\mathbf{S}_\mathrm{c})/c^2$ (Eq.~\eqref{eq:ch1.rotation-generator-split}), so the cycle-averaged momentum flux through the interface is $\langle\mathbf{g}_{\mathrm{em}}\rangle\cdot\hat{\mathbf{n}}$ times $c$ when the wave is absorbed; specular reversal reverses $\langle\mathbf{g}_{\mathrm{em}}\rangle$ and doubles the magnitude of the momentum delivered to the reflector \cite{stratton1941}.
\end{example}

\begin{figure}[t]
\centering
\VolIPhysFig{\VolIPhysPlaneWaveRadiationPressure}
\caption{Uniform plane wave in $+\hat{\mathbf{z}}$ on field control volume $V_f$ ($z<0$): outward interface normal $\hat{\mathbf{n}}=+\hat{\mathbf{z}}$ at $z=0$; perfect absorber receives radiation pressure $|\langle\mathbf{S}\rangle|/c$, perfect reflector at normal incidence $2|\langle\mathbf{S}\rangle|/c$; transverse $\E$, $\H$, and $\langle\mathbf{S}\rangle$ remain right-handed \cite{stratton1941,maxwell1873}.}
\label{fig:ch1.plane-wave-handedness}
\end{figure}
Integrating the momentum flux over a control volume yields the harmonic steady-state balance
\begin{equation}
\label{eq:ch1.linear-momentum-balance}
\langle\mathbf{f}_{\mathrm{src}}\rangle=\divg\Tavg,
\end{equation}
where \(\langle\mathbf{f}_{\mathrm{src}}\rangle=\tfrac12\Re\{\rho\E^*+\Jimp\times\B^*\}\) is the cycle-averaged source-exchange density (distinct from the instantaneous \(\mathbf{f}_{\mathrm{src}}\) of Equation~\eqref{eq:ch1.maxwell-source-force}) and \(\partial_t\langle\mathbf{g}_{\mathrm{em}}\rangle=\mathbf{0}\) on the harmonic branch \cite{stratton1941}. Eq.~\eqref{eq:ch1.linear-momentum-balance} is the steady phasor limit of Eq.~\eqref{eq:ch1.momentum-balance-instantaneous}; mechanical force on carriers remains \(\mathbf{f}_{\mathrm{matter}}\) of Eq.~\eqref{eq:ch1.lorentz-force}, equal to \(\mathbf{f}_{\mathrm{src}}\) only under Route~A.

\phantomsection\label{eq:ch1.energy-momentum-directional-ratio}
The plane wave provides the simplest standard example in which energy flow and momentum flow become visibly parallel \cite{poynting1884,stratton1941}. Classically, for an outgoing wave with $\mathbf{k}=k\hat{\mathbf{z}}$ and transverse electric field $\E=E_0 e^{-\jj k z}\hat{\mathbf{x}}$,
\begin{equation}
\label{eq:ch1.plane-wave-transverse}
\H=\frac{1}{\eta_0}\,\hat{\mathbf{k}}\times\E,
\qquad
\langle\mathbf{S}\rangle=\frac{|\E|^{2}}{2\eta_0}\,\hat{\mathbf{k}},
\end{equation}
where $\eta_0=\sqrt{\mu_0/\varepsilon_0}$ is the vacuum intrinsic impedance \cite{stratton1941}. The time-averaged Poynting vector and normal Maxwell traction then point along $\hat{\mathbf{k}}$, with radiation pressure $|\mathbf{t}|=|\langle\mathbf{S}\rangle|/c$ on the same Maxwell branch. In structured near fields energetic and mechanical transport does not necessarily be co-directed---a common situation in reactive storage zones \cite{stratton1941}.

\phantomsection\label{prop:ch1.energy-momentum-consistency}
For a Maxwell solution branch fixed by impressed sources, interface data, boundary completion, and vacuum constitutive law on the integration region, Eqs.~\eqref{eq:ch1.poynting-complex} and~\eqref{eq:ch1.maxwell-stress} are quadratic functionals of the same \((\E,\H)\). Poynting and stress surface integrals on \(\partial V\) are defined only when the corresponding integrands are integrable on that surface---typically when \((\E,\H)\) has sufficient regularity on \(V\) (Subsection~\ref{sec:ch01.0.3}); Equation~\eqref{eq:ch1.energy-class} on \(V\) is helpful but not sufficient in isolation, and finite \(P_{\mathrm{rad}}\) on a shell \(S_R\) alone does not replace those hypotheses \cite{stratton1941}. Both balances then describe transport generated by one field state, which is the common basis for the angular-flux discussion later in the chapter \cite{stratton1941}.

Eqs.~\eqref{eq:ch1.poynting-complex} and
\eqref{eq:ch1.maxwell-stress} presuppose harmonic averaging in passive media with
local \(\varepsilon,\mu\). Dispersive, anisotropic, or strongly time-varying
constitutive laws modify both storage terms and traction operators. Active boundaries
can inject net gain that invalidates the passive-balance route used in
Subsection~\ref{sec:ch01.1.3}
\cite{stratton1941}.

% R10.2 — table layout (slightly widened definition columns)
\begin{table}[t]
\centering
\caption{Fundamental harmonic electromagnetic quantities in vacuum (cycle-averaged where noted).}
\label{tab:ch1.fundamental-quantities}
\small
\begin{tabular}{@{}p{0.13\linewidth}p{0.27\linewidth}p{0.23\linewidth}p{0.19\linewidth}p{0.11\linewidth}@{}}
\toprule
Quantity & Definition & Physical meaning & Conservation law & Units \\
\midrule
$\mathbf{S}_\mathrm{c}$, $\langle\mathbf{S}\rangle$
& $\mathbf{S}_\mathrm{c}=\tfrac12\E\times\H^*$, $\langle\mathbf{S}\rangle=\Re\mathbf{S}_\mathrm{c}$
& Energy flux density & Poynting theorem & W\,m$^{-2}$ \\
$\mathbf{g}_{\mathrm{em}}$, $\langle\mathbf{g}_{\mathrm{em}}\rangle$
& $\mathbf{g}_{\mathrm{em}}=\varepsilon_0\,\E\times\B$; $\langle\mathbf{g}_{\mathrm{em}}\rangle=\varepsilon_0\mu_0\langle\mathbf{S}\rangle$
& Instantaneous vs cycle-averaged momentum density
& Momentum balance with $\Tinst$ / $\Tavg$
& kg\,m$^{-2}$\,s$^{-1}$ \\
$\Tinst$, $\Tavg$
& Eqs.~\eqref{eq:ch1.maxwell-stress-components}, \eqref{eq:ch1.maxwell-stress}
& Instantaneous and cycle-averaged Maxwell stress tensors
& $\langle\mathbf{f}_{\mathrm{src}}\rangle=\divg\Tavg$ (steady; Eq.~\eqref{eq:ch1.linear-momentum-balance})
& N\,m$^{-2}$ \\
$\mathbf{j}_{\mathrm{em}}$, $\langle\mathbf{j}_{\mathrm{em}}\rangle$
& $\mathbf{j}_{\mathrm{em}}=\mathbf{r}\times\mathbf{g}_{\mathrm{em}}$; $\langle\mathbf{j}_{\mathrm{em}}\rangle=\mathbf{r}\times\langle\mathbf{g}_{\mathrm{em}}\rangle$
& Instantaneous vs cycle-averaged angular-momentum density
& Angular-momentum balance
& kg\,m$^{-1}$\,s$^{-1}$ \\
$\Minst$, $\Mavg$
& $\Minst=\mathbf{r}\times\Tinst$; $\Mavg=\mathbf{r}\times\Tavg$
& Stress-moment density; surface density $\Mavg\cdot\hat{\mathbf{n}}$
& $\boldsymbol{\mathcal{T}}_J=\oint\Mavg\cdot\hat{\mathbf{n}}\,dS$
& N\,m$^{-1}$ ($\Minst$, $\Mavg$); N\,m ($\boldsymbol{\mathcal{T}}_J$) \\
\bottomrule
\end{tabular}
\end{table}

\begin{table}[t]
\centering
\caption{Four angular concepts that answer different physical questions.}
\label{tab:ch1.four-not-confused}
\small
\begin{tabular}{@{}p{0.24\linewidth}p{0.70\linewidth}@{}}
\toprule
Concept & Physical question answered \\
\midrule
Azimuthal phase $e^{\jj m\phi}$
& How is the transverse field organized in angle on a plane or chart? \\
Polarization / helicity
& How does the electric vector rotate in the transverse plane? \\
Angular-momentum density $\mathbf{j}_{\mathrm{em}}$
& Where is rotational action stored in the harmonic field at a point? \\
Angular-momentum flux $\boldsymbol{\Phi}_{J,\mathrm{out}}$
& What signed angular-momentum transport crosses a closed radiative shell? \\
\bottomrule
\end{tabular}
\end{table}

\begin{table}[t]
\centering
\caption{Near, transition, and far zones for an electric dipole ($r>0$, source-free exterior).}
\label{tab:ch1.near-far-zones}
\footnotesize
\begin{tabular}{@{}p{0.11\linewidth}p{0.17\linewidth}p{0.22\linewidth}p{0.22\linewidth}p{0.22\linewidth}@{}}
\toprule
Region & Dominant $r$-dependence & Interpretation & Poynting character & Angular-momentum note \\
\midrule
Near ($kr\ll1$) & $r^{-3}$, $r^{-2}$ & Static / induction structure & Large reactive $\Im(\mathbf{S}_\mathrm{c}\cdot\hat{\mathbf{r}})$ & Near-zone stress structure; not identical to stored energy \\
Transition ($kr\sim1$) & Mixed & Handover to radiation & $|P_{\mathrm{reac}}(r)|/P_{\mathrm{rad}}$ large & Shell choice matters for reading flux \\
Far ($kr\gg1$) & $r^{-1}$ & Outgoing wave & $P_{\mathrm{rad}}$ fixed on nested shells & $\E$, $\H$ transverse; $U(\theta)\propto\sin^2\theta$ \\
\bottomrule
\end{tabular}
\end{table}


Figure~\ref{fig:ch1.wp-reactive-vs-radiative} illustrates how the reactive--radiative distinction enters these balances.
\begin{figure}[!htbp]
\centering
\VolIPhysFig{\VolIPhysReactiveVsRadiative}
\caption{Concentric evaluation shells at $kr=0.5$, $kr=3$, and $kr=10$ for a small dipole: reactive power exchange dominates on the inner shell, while outward $\langle\mathbf{S}\rangle$ and steady traction emerge on the radiative shell \cite{stratton1941,chu1948}.}
\label{fig:ch1.wp-reactive-vs-radiative}
\end{figure}

Figure~\ref{fig:ch1.wp-reactive-vs-radiative} marks three concentric evaluation shells at $kr=0.5$, $kr=3$, and $kr=10$ for an electrically small dipole ($k\ell\ll1$). At $kr=0.5$ reactive
storage dominates: Equation~\eqref{eq:ch1.poynting-real-imag-split} yields
\(\int_{4\pi}\Re(\mathbf{S}_\mathrm{c}\cdot\hat{\mathbf{r}})\,dS=P_{\mathrm{rad}}\ll
\int_{4\pi}|\Im(\mathbf{S}_\mathrm{c}\cdot\hat{\mathbf{r}})|\,dS=|P_{\mathrm{reac}}|\), and traction on a
small patch is oscillatory rather than proportional to \(P_{\mathrm{rad}}/c\). At
$kr=3$ the transition regime begins; at
$kr=10$ outward $\Re(\mathbf{S}_\mathrm{c})$ dominates and the shell supports the radiation-pressure and torque integrals of section~\ref{sec:ch01.6.4} \cite{stratton1941,chu1948}.
The conservation law holds on any admissible closed surface; a predominantly radiative shell is preferred when the integral is to be read as outward radiated transport rather than local reactive exchange \cite{chu1948,stratton1941}.

Using the uniform plane-wave example of Subsection~\ref{sec:ch01.1.1}, take $\mathbf{S}_\mathrm{c}=(|E_0|^2/2\eta_0)\hat{\mathbf{z}}$
and evaluate power through a disk of radius $a$ normal to $\hat{\mathbf{z}}$:
\begin{equation}
\label{eq:ch1.worked-aperture-power}
P(a)=\int_0^{2\pi}\!\!\int_0^a \Re(\mathbf{S}_\mathrm{c}\cdot\hat{\mathbf{z}})\,\rho\,d\rho\,d\phi
=\frac{\pi a^2 |E_0|^2}{2\eta_0}.
\end{equation}
For a uniform plane wave, $P(a)$ in Equation~\eqref{eq:ch1.worked-aperture-power} is the active power intercepted by a disk of radius $a$ whose normal aligns with $\langle\mathbf{S}\rangle$ \cite{poynting1884}.

Within the harmonic vacuum class treated here, Eqs.~\eqref{eq:ch1.poynting-complex} and~\eqref{eq:ch1.maxwell-stress} fix how power and linear momentum cross a surface. The next subsection takes the moment of that momentum balance about an origin and introduces $\Mavg=\mathbf{r}\times\Tavg$ as the stress-moment density from which torque integrals are built \cite{stratton1941}.

\subsection{Angular momentum conservation and torque balance}
\label{sec:ch01.1.5}

Linear momentum is already balanced by impressed force, stored electromagnetic momentum, and stress traction. Rotational balance asks how that linear transport appears about a chosen origin. Multiplying the index form of Equation~\eqref{eq:ch1.momentum-balance-instantaneous} by $\varepsilon_{ijk}r_j$ and applying the Levi-Civita contraction gives, in index notation,
\begin{equation}
\label{eq:ch1.angular-momentum-index-start}
\varepsilon_{ijk}\,r_j\!\left(\partial_t g_k+f_{\mathrm{src},k}-\partial_l T_{kl}\right)
=0.
\end{equation}
The product rule expands the stress term as
\begin{equation}
\label{eq:ch1.angular-momentum-product-rule}
\varepsilon_{ijk}\,r_j\,\partial_l T_{kl}
=\partial_l\!\left(\varepsilon_{ijk}\,r_j T_{kl}\right)
-\varepsilon_{ijk}\,T_{kj},
\end{equation}
The remaining term $\varepsilon_{ijk}T_{kj}$ would represent torque density carried by an antisymmetric stress. The Maxwell tensor is symmetric ($T_{kj}=T_{jk}$), so $\varepsilon_{ijk}T_{kj}=0$ and that contribution vanishes \cite{maxwell1873,stratton1941}. Define the instantaneous stress-moment density $\Minst=\mathbf{r}\times\Tinst$ and the angular-momentum density $\mathbf{j}_{\mathrm{em}}=\mathbf{r}\times\mathbf{g}_{\mathrm{em}}$. Equations~\eqref{eq:ch1.angular-momentum-index-start} and~\eqref{eq:ch1.angular-momentum-product-rule} then combine to
\begin{equation}
\label{eq:ch1.angular-momentum-balance}
\frac{\partial\mathbf{j}_{\mathrm{em}}}{\partial t}+\boldsymbol{\tau}_{\mathrm{src}}
=\divg\Minst,
\end{equation}
with source-exchange torque density $\boldsymbol{\tau}_{\mathrm{src}}=\mathbf{r}\times\mathbf{f}_{\mathrm{src}}$ (N\,m\,m$^{-3}$) obtained from the same $\rho$ and $\Jimp$ as Equation~\eqref{eq:ch1.momentum-balance-instantaneous} \cite{maxwell1873,stratton1941}. Mechanical torque on carriers is $\boldsymbol{\tau}_{\mathrm{matter}}=\mathbf{r}\times\mathbf{f}_{\mathrm{matter}}$ of Eq.~\eqref{eq:ch1.lorentz-force}; $\boldsymbol{\tau}_{\mathrm{src}}=\boldsymbol{\tau}_{\mathrm{matter}}$ only under Route~A on the support where $\mathbf{f}_{\mathrm{src}}=\mathbf{f}_{\mathrm{matter}}$. The resultant mechanical torque $\mathbf{T}_{\mathrm{matter}}=\int_{V_m}\boldsymbol{\tau}_{\mathrm{matter}}\,dV$ (N\,m) follows from Eq.~\eqref{eq:ch1.angular-balance-integral} under Route~A throughout $V_m$, or from the closed stress integral on $\partial V_m$ in source-free vacuum under Route~B (Eq.~\eqref{eq:ch1.torque-stress-surface}) \cite{stratton1941,beth1936}. In harmonic steady state with vanishing storage rate, $\mathbf{T}_{\mathrm{matter}}=\boldsymbol{\mathcal{T}}_J=-\boldsymbol{\Phi}_{J,\mathrm{out}}$ when the integration domain and route assumptions support that identification \cite{beth1936}. Eq.~\eqref{eq:ch1.angular-momentum-balance} is the instantaneous rotational counterpart of Eq.~\eqref{eq:ch1.momentum-balance-instantaneous}: stored field angular momentum changes when source-exchange torque density acts or when stress-moment flux diverges.

For time-harmonic fields under $\exp(+\jj\omega t)$, the same conservation laws may be cycle-averaged. Averaging Equation~\eqref{eq:ch1.momentum-balance-instantaneous} yields $\partial_t\langle\mathbf{g}_{\mathrm{em}}\rangle+\langle\mathbf{f}_{\mathrm{src}}\rangle=\divg\Tavg$, and averaging Eq.~\eqref{eq:ch1.angular-momentum-balance} yields
\begin{equation}
\label{eq:ch1.angular-momentum-balance-averaged}
\frac{\partial\langle\mathbf{j}_{\mathrm{em}}\rangle}{\partial t}+\langle\boldsymbol{\tau}_{\mathrm{src}}\rangle
=\divg\Mavg,
\end{equation}
where $\Mavg=\mathbf{r}\times\Tavg$, $\langle\mathbf{j}_{\mathrm{em}}\rangle=\mathbf{r}\times\langle\mathbf{g}_{\mathrm{em}}\rangle$, and $\langle\boldsymbol{\tau}_{\mathrm{src}}\rangle=\mathbf{r}\times\langle\mathbf{f}_{\mathrm{src}}\rangle$ \cite{stratton1941}. Under Route~A, $\langle\boldsymbol{\tau}_{\mathrm{matter}}\rangle=\langle\boldsymbol{\tau}_{\mathrm{src}}\rangle$ on the relevant support; under Route~B, $\langle\boldsymbol{\tau}_{\mathrm{src}}\rangle$ remains field--source bookkeeping and mechanical torque on a body is read from the stress surface integral (Equation~\eqref{eq:ch1.torque-stress-surface}). Angle brackets mark cycle-averaged harmonic quantities; $\Minst$ and $\Mavg$ are not interchangeable in the same displayed balance.

In linear isotropic vacuum the cycle-averaged momentum density is
\begin{equation}
\label{eq:ch1.rotation-generator-split}
\langle\mathbf{g}_{\mathrm{em}}\rangle=\varepsilon_0\mu_0\,\langle\mathbf{S}\rangle
=\varepsilon_0\mu_0\,\Re(\mathbf{S}_\mathrm{c}),
\end{equation}
and the corresponding angular-momentum density is
\begin{equation}
\label{eq:ch1.angular-momentum-density}
\langle\mathbf{j}_{\mathrm{em}}\rangle=\mathbf{r}\times\langle\mathbf{g}_{\mathrm{em}}\rangle.
\end{equation}
\addtocounter{equation}{-1}\refstepcounter{equation}\label{eq:ch1.angular-density-split}%

Integrating Equation~\eqref{eq:ch1.angular-momentum-balance-averaged} over a control volume $V_m$ with outward normal $\hat{\mathbf{n}}$ yields a valid volume identity only where that local balance holds pointwise (vacuum $\varepsilon_0,\mu_0$ region with the same $(\rho,\Jimp)$ as Maxwell's equations). When $V_m$ encloses matter, the surface balance Eqs.~\eqref{eq:ch1.torque-stress-surface} and~\eqref{eq:ch1.outward-am-flux} on $\partial V_m$ in source-free vacuum supply the mechanical resultant under Route~B; the volume integral of $\langle\boldsymbol{\tau}_{\mathrm{src}}\rangle$ is then source-exchange bookkeeping, not automatic load torque \cite{stratton1941}. With that scope,
\begin{equation}
\label{eq:ch1.angular-balance-integral}
\frac{d}{dt}\int_{V_m} \langle\mathbf{j}_{\mathrm{em}}\rangle\,dV
+\int_{V_m} \langle\boldsymbol{\tau}_{\mathrm{src}}\rangle\,dV
=\oint_{\partial V_m}\Mavg\cdot\hat{\mathbf{n}}\,dS.
\end{equation}
The volume term $\int_{V_m}\langle\boldsymbol{\tau}_{\mathrm{src}}\rangle\,dV$ is the integrated source-exchange torque associated with $(\rho,\Jimp)$; it equals $\int_{V_m}\langle\boldsymbol{\tau}_{\mathrm{matter}}\rangle\,dV$ only under Route~A throughout $V_m$ \cite{stratton1941}. The surface term is the stress moment $\boldsymbol{\mathcal{T}}_J=\oint\Mavg\cdot\hat{\mathbf{n}}\,dS$ \cite{stratton1941}. At harmonic steady state with $\partial_t\langle\mathbf{j}_{\mathrm{em}}\rangle=\mathbf{0}$, Equation~\eqref{eq:ch1.torque-stress-surface} and~\eqref{eq:ch1.outward-am-flux} give the \emph{resultant} mechanical torque $\mathbf{T}_{\mathrm{matter}}=\boldsymbol{\mathcal{T}}_J-d\mathbf{J}_{\mathrm{em,store}}/dt=-\boldsymbol{\Phi}_{J,\mathrm{out}}-d\mathbf{J}_{\mathrm{em,store}}/dt$ when the enclosing surface lies in source-free vacuum and the route assumptions are satisfied (N\,m), not the local density $\boldsymbol{\tau}_{\mathrm{src}}$ or $\boldsymbol{\tau}_{\mathrm{matter}}$ \cite{beth1936}. Beth \cite{beth1936} measured load torque from net flux through the plate's control surface (incident minus transmitted), not from equating an impressed-source volume integral to load mechanics without that surface balance \cite{stratton1941}.

Figure~\ref{fig:ch1.momentum-angular-geometry} traces the cycle-averaged chain at a field point and on a surface element: $\langle\mathbf{g}_{\mathrm{em}}\rangle$, $\langle\mathbf{j}_{\mathrm{em}}\rangle=\mathbf{r}\times\langle\mathbf{g}_{\mathrm{em}}\rangle$, traction $\Tavg\cdot\hat{\mathbf{n}}$, and stress moment $\mathbf{r}\times(\Tavg\cdot\hat{\mathbf{n}})$ \cite{maxwell1873,beth1936}.

\begin{figure}[t]
\centering
\VolIPhysFigWide{\VolIPhysMomentumAngularGeometry}
\caption{Cycle-averaged angular-momentum geometry: at observation point $\mathbf{r}$, $\langle\mathbf{g}_{\mathrm{em}}\rangle$ and $\langle\mathbf{j}_{\mathrm{em}}\rangle=\mathbf{r}\times\langle\mathbf{g}_{\mathrm{em}}\rangle$; on enclosing surface element with outward $\hat{\mathbf{n}}$, traction $\Tavg\cdot\hat{\mathbf{n}}$ and stress moment $\mathbf{r}\times(\Tavg\cdot\hat{\mathbf{n}})$ that enters torque integrals \cite{maxwell1873,beth1936,stratton1941}.}
\label{fig:ch1.momentum-angular-geometry}
\end{figure}

The cycle-averaged stress-moment density $\Mavg=\mathbf{r}\times\Tavg$ satisfies
$\Mavg\cdot\hat{\mathbf{n}}=\mathbf{r}\times(\Tavg\cdot\hat{\mathbf{n}})$ on $\partial V$ (stress-moment surface density).
At harmonic steady state the cycle-averaged \emph{resultant} mechanical torque on matter enclosed in $V_m$ is
\begin{equation}
\label{eq:ch1.torque-stress-surface}
\mathbf{T}_{\mathrm{matter}}
\equiv
\int_{V_m}\langle\boldsymbol{\tau}_{\mathrm{matter}}\rangle\,dV
=\oint_{\partial V_m}\Mavg\cdot\hat{\mathbf{n}}\,dS
-\frac{d}{dt}\int_{V_m} \langle\mathbf{j}_{\mathrm{em}}\rangle\,dV
=\oint_{\partial V_m}\mathbf{r}\times(\Tavg\cdot\hat{\mathbf{n}})\,dS
-\frac{d}{dt}\int_{V_m} \langle\mathbf{j}_{\mathrm{em}}\rangle\,dV,
\end{equation}
with outward $\hat{\mathbf{n}}$ on $\partial V_m$ \cite{beth1936,stratton1941}. The volume identity $\int_{V_m}\langle\boldsymbol{\tau}_{\mathrm{matter}}\rangle\,dV=\int_{V_m}\langle\boldsymbol{\tau}_{\mathrm{src}}\rangle\,dV$ holds under Route~A throughout $V_m$. Under Route~B, $\mathbf{T}_{\mathrm{matter}}$ is obtained from the stress surface integral when $\partial V_m$ lies in source-free vacuum. The stress moment on the same closed surface is
\begin{equation}
\label{eq:ch1.stress-moment-integral}
\boldsymbol{\mathcal{T}}_J
\equiv
\oint_{\partial V_m}\mathbf{r}\times(\Tavg\cdot\hat{\mathbf{n}})\,dS
=\oint_{\partial V_m}\Mavg\cdot\hat{\mathbf{n}}\,dS.
\end{equation}
Equation~\eqref{eq:ch1.torque-stress-surface} yields $\mathbf{T}_{\mathrm{matter}}=\boldsymbol{\mathcal{T}}_J-d\mathbf{J}_{\mathrm{em,store}}/dt$ with stored field angular momentum $\mathbf{J}_{\mathrm{em,store}}=\int_{V_m}\langle\mathbf{j}_{\mathrm{em}}\rangle\,dV$. Outward electromagnetic angular-momentum flux through $\partial V_m$ is
\begin{equation}
\label{eq:ch1.outward-am-flux}
\boldsymbol{\Phi}_{J,\mathrm{out}}
\equiv
\dot{\mathbf{J}}_{\mathrm{em,out}}
=-\boldsymbol{\mathcal{T}}_J
=-\oint_{\partial V_m}\mathbf{r}\times(\Tavg\cdot\hat{\mathbf{n}})\,dS,
\end{equation}
so $\mathbf{T}_{\mathrm{matter}}=-\boldsymbol{\Phi}_{J,\mathrm{out}}-d\mathbf{J}_{\mathrm{em,store}}/dt$. The stress moment $\boldsymbol{\mathcal{T}}_J$ is not the outward flux: outward transport is $\boldsymbol{\Phi}_{J,\mathrm{out}}=-\boldsymbol{\mathcal{T}}_J$ \cite{stratton1941,beth1936}. In harmonic steady state with $d\mathbf{J}_{\mathrm{em,store}}/dt=\mathbf{0}$, $\mathbf{T}_{\mathrm{matter}}=\boldsymbol{\mathcal{T}}_J=-\boldsymbol{\Phi}_{J,\mathrm{out}}$; source and load resultants $\boldsymbol{\tau}_{\mathrm{source}}=-\boldsymbol{\Phi}_{J,\mathrm{out}}$ and $\boldsymbol{\tau}_{\mathrm{load}}=+\boldsymbol{\Phi}_{J,\mathrm{out}}$ apply only when the declared control surface encloses the support and the load intercepts the net flux \cite{beth1936}.

Belinfante's 1940 field-theoretic analysis \cite{belinfante1940} showed that canonical
energy--momentum and symmetric (kinetic) tensors differ by the divergence of a spin
current. Correspondingly, total angular momentum
$\int_V \mathbf{j}_{\mathrm{em}}\,dV$ and closed-surface flux
$\oint \Mavg\cdot\hat{\mathbf{n}}\,dS$ are gauge invariant, whereas local spin--orbit splits are convention-dependent \cite{belinfante1940}. Classically, primary quantities are $\mathbf{j}_{\mathrm{em}}=\mathbf{r}\times\mathbf{g}_{\mathrm{em}}$ and stress-moment flux through $\Mavg\cdot\hat{\mathbf{n}}$. A separate
``residual spin'' obtained by subtracting these identical definitions is not used.
Detailed spin--orbit and canonical decompositions are representation-dependent and are
developed where the modal basis is fixed (Chapter~\ref{ch:03}) \cite{belinfante1940,tai1994}.
Beth's experiment on a wave plate under circular polarization
\cite{beth1936} showed that the spin-associated part carries mechanical consequences; later orbital-indexed beam experiments
\cite{allen1992} extended the vocabulary to spatial phase gradients; the measurable
balance remains Equation~\eqref{eq:ch1.angular-balance-integral} \cite{belinfante1940}.

For torque about axis $\hat{\mathbf{u}}$, project Equation~\eqref{eq:ch1.angular-balance-integral}:
\begin{equation}
\label{eq:ch1.torque-observable}
\tau_{\hat{\mathbf{u}}}(V_m)
=\hat{\mathbf{u}}\cdot
\left(
\oint_{\partial V_m}\Mavg\cdot\hat{\mathbf{n}}\,dS
-\frac{d}{dt}\int_{V_m} \langle\mathbf{j}_{\mathrm{em}}\rangle\,dV
\right),
\end{equation}
which is the axis-projected form of Equation~\eqref{eq:ch1.angular-balance-integral}.
Torque is therefore a projected surface integral of the stress moment, not a scalar read from a phase portrait \cite{beth1936,stratton1941}.

To compare directional concentration of stress moment on a sphere of radius $r$, note that $\avgStressMoment\cdot\hat{\mathbf{r}}$ states how the local stress moment projects onto each solid-angle sector. A main-lobe concentration in $\|\avgStressMoment\cdot\hat{\mathbf{r}}\|$ does not in isolation establish net torque about an axis until Equation~\eqref{eq:ch1.torque-observable} is evaluated with the declared origin and surface \cite{stratton1941,beth1936}.
\phantomsection\label{eq:ch1.angular-flux-directional-index}

Under the conventions adopted here, torque about axis $\hat{\mathbf{u}}$ is fixed once origin $\mathbf{r}_0$, axis $\hat{\mathbf{u}}$, and control surface $\partial V$ entering $\Mavg=\mathbf{r}\times\Tavg$ are declared \cite{stratton1941}. Changing any of these geometric choices rotates flux projections without altering the underlying Maxwell pair; published torque values therefore state $(\mathbf{r}_0,\hat{\mathbf{u}},\partial V)$ alongside the field solution \cite{beth1936}.

Consider two Maxwell field states \(\mathcal{F}_1,\mathcal{F}_2\) at common radius
\(r\), with equal total radiated power but different directional stress-moment
densities \(\mathbf{M}_1\cdot\hat{\mathbf{r}}\neq
\mathbf{M}_2\cdot\hat{\mathbf{r}}\) on a subset of $4\pi$ of nonzero measure. Equal $P_{\mathrm{rad}}$ constrains only the energetic integral; it does not fix the stress-moment distribution direction by direction. Projected torque through Eq.~\eqref{eq:ch1.torque-observable} can therefore differ between the two states even when their radiation patterns share the same total power \cite{stratton1941}.

\phantomsection\label{prop:ch1.torque-observable-criterion}
\begin{figure}[t]
\centering
\VolIPhysFig{\VolIPhysTorqueBalanceSurface}
\caption{Spherical control shell $S_R$ with axis $\hat{\mathbf{z}}$: axisymmetric radiation with azimuthal order $m=0$ may carry zero $\dot J_{z,\mathrm{out}}$ from the stress-moment integral even when $P_{\mathrm{rad}}$ is finite; circular polarization fixes local handedness but does not by itself guarantee nonzero $\dot J_{z,\mathrm{out}}$ on $S_R$ under Eq.~\eqref{eq:ch1.outward-am-flux}. Azimuthal modes with $m\neq 0$ can break symmetry and yield nonzero net flux when the Maxwell-valid field and surface integral support it \cite{beth1936,stratton1941}.}
\label{fig:ch1.wp-torque-balance-surface}
\end{figure}

Figure~\ref{fig:ch1.wp-torque-balance-surface} contrasts symmetric radiation with azimuthally phased
radiation on the same spherical shell: when the stress flux integrand possesses definite azimuthal
symmetry about $\hat{\mathbf{z}}$, the solid-angle integral of
$[\mathbf{r}\times(\Tavg\cdot\hat{\mathbf{n}})]\cdot\hat{\mathbf{z}}$ may vanish even though
local phase structure is visible on a near-field display---the symmetry argument Beth exploited for
circular polarization \cite{beth1936,stratton1941}.

For a matter-enclosing control volume $V_m$ whose stress-evaluation surface lies in source-free vacuum,
Equation~\eqref{eq:ch1.angular-balance-integral} reduces at harmonic steady state to
\(T_{z,\mathrm{matter}}=\oint_{\partial V_m}[\mathbf{r}\times(\Tavg\cdot\hat{\mathbf{n}})]\cdot\hat{\mathbf{z}}\,dS\)
and \(\dot J_{z,\mathrm{em,out}}=-T_{z,\mathrm{matter}}\) by Equation~\eqref{eq:ch1.outward-am-flux} (here $T_{z,\mathrm{matter}}$ is the resultant torque in N\,m, not the local density $\boldsymbol{\tau}_{\mathrm{src}}$ of Eq.~\eqref{eq:ch1.angular-momentum-balance}).
On sphere \(S_R\), use \(\mathbf{r}=R\hat{\mathbf{r}}\) and
\(dS=R^2\sin\theta\,d\theta\,d\phi\). For axisymmetric radiation with azimuthal index $m=0$, the far-zone pattern carries no azimuthal phase circulation about $\hat{\mathbf{z}}$, so $\dot J_{z,\mathrm{out}}$ from orbital winding vanishes while $P_{\mathrm{rad}}$ remains finite. For a uniform circularly polarized plane wave, the symmetric Maxwell stress tensor on a centered transverse disk gives $\dot J_{z,\mathrm{em,out}}=0$ while power crossing the disk is nonzero (Subsection~\ref{sec:ch01.6.1}); $(m+\sigma)P/\omega$ and $\sigma P/\omega$ symbols elsewhere are modal-class evaluators unless derived from a Maxwell-valid field on the chosen surface \cite{beth1936,stratton1941}.
Torque is an integral consequence of the Maxwell stress tensor through Equation~\eqref{eq:ch1.torque-observable}, not a local phase label
\cite{stratton1941}.

Eqs.~\eqref{eq:ch1.angular-momentum-balance}
and \eqref{eq:ch1.angular-balance-integral} assume a fixed origin, a symmetric
stress-tensor completion, and local linear constitutive behavior in the present
regime. Strongly nonlocal or dissipative media may require additional closure before
torque partition is interpreted without ambiguity \cite{stratton1941}.

\section{Angular Structure of Electromagnetic Fields}
\label{sec:ch01.2}

Once energy, linear momentum, and stress-moment flux are defined on a control surface, one can ask how the field's angular organization---its phase winding, polarization handedness, and direction-resolved intensity---connects to signed transport and torque on that same surface \cite{allen1992,stratton1941,beth1936}. The subsections below follow that connection from rotating feeds and phased apertures to helical phase, stress integrals, and the reactive versus radiative shells on which those integrals are evaluated.
\subsection{Rotating source versus rotating pattern}
\label{sec:ch01.2.0}

Far-zone isophase charts on a sphere $S_R$ often show spiral contours in $\arg E$. That display records how transverse phase winds with azimuth; it does not by itself fix signed angular-momentum transport or tell a mechanically rotating feed from a stationary array whose elements impose $\exp(\jj m\phi)$ \cite{beth1936,harrington2001}. Only a Maxwell-valid $(\E,\H)$ together with the stress--torque observables of Eqs.~\eqref{eq:ch1.directional-am-flux} and~\eqref{eq:ch1.torque-observable} answers whether rotational action is radiated \cite{stratton1941,beth1936}.

Figure~\ref{fig:ch1.signature-rotation-duality} contrasts those two excitations at the level of the source geometry; the following paragraphs compare their balances when the patterns on $S_R$ look alike.

\begin{figure}[t]
\centering
\VolIPhysFig{\VolIPhysRotationDuality}
\caption{Physical source rotation (a) versus stationary phased feed (b): both can yield identical spiral isophase contours on far-zone shell $S_R$, yet signed flux and torque observables distinguish radiated angular action from appearance alone.}
\label{fig:ch1.signature-rotation-duality}
\end{figure}

In panel~(a) the compact dipole \(\mathbf{p}\) of Subsection~\ref{sec:ch01.0.2} rotates physically about \(\hat{\mathbf{z}}\). Real mechanical rotation is generally non-sinusoidal and can generate sidebands; the harmonic analysis at fixed \(\omega\) therefore uses a stated single-frequency approximation---frozen orientation, selected sideband, or slow-rotation limit---rather than the full time history \cite{stratton1941}. The motorized constraint that supplies \(\Omega_s\) enters the angular-momentum budget separately from the field stress integral.

In panel~(b) the hardware is fixed in the laboratory. A ring of radius \(a\ll\lambda\) carries aperture phasing \(\E_a(\phi)=\E_0\exp(\jj m\phi)\hat{\boldsymbol{\rho}}\), so the same azimuthal index \(m\) may appear on $S_R$ without any rotating conductor \cite{harrington2001}. Passive feed networks can impose that phase law; they do not duplicate the mechanical power channel of panel~(a).

On sphere \(S_R\) in the radiative zone, define the directional observables of
Eqs.~\eqref{eq:ch1.radiation-intensity-Ur}--\eqref{eq:ch1.directional-am-flux}.
For configuration~(b) with single azimuthal index \(m\) and axisymmetric magnitude,
\begin{equation}
\label{eq:ch1.signature-phase-ramp}
\E(R,\theta,\phi)\propto \E_\theta(\theta)\,e^{\jj m\phi},
\end{equation}
labels the azimuthal winding on $S_R$ for the phased feed. Signed transport $\dot J_{z,\mathrm{em,out}}$ is obtained only from Equation~\eqref{eq:ch1.ex61-flux-from-stress} or Eqs.~\eqref{eq:ch1.signature-flux-torque-pair}--\eqref{eq:ch1.signature-source-torque-z} for a Maxwell-valid $(\E,\H)$; Eq.~\eqref{eq:ch1.signature-phase-ramp} does not imply $\dot J_{z,\mathrm{em,out}}=(m/\omega)P_{\mathrm{rad}}$ \cite{stratton1941,allen1992}.

For panel~(a) with slow rotation \(\Omega_s\ll\omega\), the radiated spectrum is a sideband superposition. A snapshot on \(S_R\) may then resemble Eq.~\eqref{eq:ch1.signature-phase-ramp}, yet torque on a Beth-type load still includes mechanical work on the feed as well as field exchange. Beth's half-wave plate on a suspension measures the torque from reversing circular polarization on the transmitted beam \cite{beth1936}. A quarter-wave plate that converts circular to linear polarization is a different arrangement: it yields \(\tau_{z,\mathrm{load}}\) through Equation~\eqref{eq:ch1.torque-observable} without reproducing Beth's suspended-plate geometry. Panel~(a) can exchange angular momentum through mechanics and through flux; panel~(b) with a passive feed exchanges torque through the field when the impressed currents are fixed in space \cite{stratton1941}.

Closed-surface integrals, not the phase portrait alone, fix transport:
\begin{equation}
\label{eq:ch1.signature-flux-torque-pair}
\Phi_{J,z,\mathrm{out}}(R)
=-\hat{\mathbf{z}}\cdot\int_{4\pi}\mathbf{j}_J(\theta,\phi;R)\,d\Omega
=-\hat{\mathbf{z}}\cdot\boldsymbol{\Phi}_{J,\mathrm{out}}(R),
\qquad
\dot J_{z,\mathrm{em,out}}(R)=\Phi_{J,z,\mathrm{out}}(R),
\end{equation}
\begin{equation}
\label{eq:ch1.signature-source-torque-z}
\tau_{z,\mathrm{source}}(R)
=\oint_{S_R}\big[\mathbf{r}\times(\Tavg\cdot\hat{\mathbf{n}})\big]\cdot\hat{\mathbf{z}}\,dS
=-\Phi_{J,z,\mathrm{out}}(R)
\end{equation}
at harmonic steady state when storage in the enclosing volume averages to zero \cite{beth1936,stratton1941}.
On \(S_R\), configuration~(b) with \(|m|\ge 1\) and \(|P_{\mathrm{reac}}(R)|/P_{\mathrm{rad}}\ll 1\) is a predominantly radiative shell; net signed \(\Phi_{J,z,\mathrm{out}}(R)\) and \(\tau_{z,\mathrm{source}}(R)\) follow only from evaluating Eqs.~\eqref{eq:ch1.signature-flux-torque-pair} and~\eqref{eq:ch1.signature-source-torque-z} for the Maxwell-valid \((\E,\H)\) on that surface. Azimuthal index and a small reactive-power ratio do not by themselves guarantee nonzero transport \cite{stratton1941,allen1992}.
Configuration~(a) matched to the same snapshot pattern can share \(U_r(\theta,\phi;R)\)
yet carry additional mechanical boundary terms on the feed support that do not appear
in a pure phase-synthesis model; omitting those terms violates
Equation~\eqref{eq:ch1.angular-balance-integral}. Conversely, a balanced pair of
counter-rotating phase ramps \(\exp(\pm\jj m\phi)\) with equal power superpose to a standing azimuthal pattern \(\sqrt{2}\cos(m\phi)\); under the fixed linear-polarization and modal assumptions of section~\ref{sec:ch01.6.2}, the signed stress integral in Equation~\eqref{eq:ch1.directional-am-total} yields \(\boldsymbol{\Phi}_{J,\mathrm{out}}(R)=\mathbf{0}\) for that template---an orbital-bookkeeping counterexample, not a universal cancellation of all spin-plus-orbital contributions for arbitrary polarization \cite{allen1992,stratton1941}.

Azimuthal index $m$ organizes phase on a chart; Eqs.~\eqref{eq:ch1.signature-flux-torque-pair} measure transport. Panels~(a) and~(b) cannot be separated by phase on $S_R$ alone, but $U_r$, the stress integrand, and $\tau_{z,\mathrm{source}}$ on the same surface can \cite{stratton1941,beth1936}.
\subsection{Directional dependence and angular momentum}
\label{sec:ch01.2.1}


Power through a sphere is read from the normal component of $\mathbf{S}_\mathrm{c}$; signed angular-momentum transport uses the stress-moment integrand built from $\Tavg$. Both observables are evaluated on the same $S_R$, but they are different functionals of $(\E,\H)$ \cite{poynting1884,belinfante1940,beth1936,stratton1941}. An axisymmetric dipole may carry finite $P_{\mathrm{rad}}$ yet vanishing net $z$-flux when opposing sectors cancel \cite{stratton1941}. Let \(\hat{\mathbf{r}}=\hat{\mathbf{r}}(\theta,\phi)\) be radial unit vector on a
radius-\(r\) sphere and \(d\Omega=\sin\theta\,d\theta\,d\phi\) the solid-angle
measure. With \(\mathbf{S}_\mathrm{c}\) from Equation~\eqref{eq:ch1.intro-poynting}, the radiation intensity is
\begin{equation}
\label{eq:ch1.radiation-intensity-Ur}
U_r(r,\theta,\phi)=
r^2\,\Re\!\left\{\mathbf{S}_\mathrm{c}(r,\theta,\phi)\cdot\hat{\mathbf{r}}\right\},
\end{equation}
in watts per steradian \cite{poynting1884,harrington2001}. The far-zone limit is $U(\theta,\phi)=\lim_{r\to\infty}U_r$, total radiated power is
\begin{equation}
\label{eq:ch1.directional-power-integral}
P_{\mathrm{rad}}(r)=\int_{4\pi}U_r(r,\theta,\phi)\,d\Omega,
\end{equation}
and directivity is
\begin{equation}
\label{eq:ch1.directivity-definition}
D(\theta,\phi)=\frac{4\pi\,U(\theta,\phi)}{P_{\mathrm{rad}}}.
\end{equation}
Equal integrated power does not imply equal directional angular-momentum structure \cite{stratton1941}.

For rotational transport, define from
Equation~\eqref{eq:ch1.angular-momentum-balance}
\begin{equation}
\label{eq:ch1.directional-am-flux}
\mathbf{j}_J(\theta,\phi;r)=
r^2\,\big[\mathbf{r}\times(\Tavg\cdot\hat{\mathbf{r}})\big].
\end{equation}
Here $\mathbf{j}_J$ is the directional stress-moment density (N\,m/sr): solid-angle resolution of $\Mavg=\mathbf{r}\times\Tavg$ on a sphere of radius $r$ \cite{beth1936,stratton1941}. Radiation intensity $U_r$ from Equation~\eqref{eq:ch1.radiation-intensity-Ur} describes energetic transport; $\mathbf{j}_J$ describes the corresponding rotational transport density.

The shell integrals are
\begin{equation}
\label{eq:ch1.directional-am-total}
\boldsymbol{\mathcal{T}}_J(r)=
\int_{4\pi}\mathbf{j}_J(\theta,\phi;r)\,d\Omega,
\qquad
\boldsymbol{\Phi}_{J,\mathrm{out}}(r)=-\boldsymbol{\mathcal{T}}_J(r).
\end{equation}
Integrated totals are a final summary of direction-resolved structure. Opposing solid-angle
sectors can carry large \(|\mathbf{j}_J|\) yet cancel in
\(\boldsymbol{\mathcal{T}}_J\); outward \(U_r>0\) along a ray does not by itself imply
signed angular-momentum transport along that ray \cite{stratton1941,belinfante1940}.

\phantomsection\label{prop:ch1.directional-nonequivalence}
Let two radiative states \(\mathcal{F}_1,\mathcal{F}_2\) at radius \(r\) satisfy
\(P_{\mathrm{rad}}^{(1)}(r)=P_{\mathrm{rad}}^{(2)}(r)\). If \(\mathbf{j}_J^{(1)}\neq\mathbf{j}_J^{(2)}\) on a subset \(\mathcal{U}\subset4\pi\) of nonzero measure while \(U_r^{(1)}=U_r^{(2)}\) on \(\mathcal{U}\), the two states differ in direction-resolved stress-moment transport even though they radiate the same power. Axis-projected torque about \(\hat{\mathbf{z}}\) can then disagree while \(P_{\mathrm{rad}}\) is shared, because \(\hat{\mathbf{z}}\cdot\mathbf{j}_J\) is redistributed over solid angle \cite{stratton1941}.

Origin declaration remains essential because angular transport is moment-based.
Changing origin re-weights directional moments even when global balances remain
unchanged. Directional flux and torque statements therefore require the same origin convention established in section~\ref{sec:ch01.1.5}.

Directional observables represent physical transport
only when outward flux is well-defined on the chosen shell and Equation~\eqref{eq:ch1.am-net-transport-condition} holds for the field under study.
In strongly reactive near zones, strongly nonlocal media, or rapidly modulated
boundaries, extra decomposition terms may be necessary before \(U_r\) and
\(\mathbf{j}_J\) are treated as direct rotational-transfer observables
\cite{stratton1941}.

For fully polarized harmonic transverse fields, the two independent electric
components \(E_\theta\) and \(E_\phi\) on a chosen surface fix the oscillation
ellipse. Left- and right-circular polarization correspond to \(\E\) and \(\H\)
tracing the same rotation sense about the propagation direction
\cite{stratton1941}. Beth's quartz-plate experiment \cite{beth1936} is the
mechanical consequence when such a field transfers angular momentum to matter
through Maxwell stress at the interface.
\phantomsection\label{eq:ch1.polarization-ellipse-angles}
\phantomsection\label{eq:ch1.jones-rotation}
\phantomsection\label{eq:ch1.coherency-matrix}
\phantomsection\label{eq:ch1.polarization-degree}

Directional dependence asks which portion of $S^{2}$ carries Poynting flux and whether that same directional distribution yields nonzero signed angular-momentum flux through the stress integral, independently of any modal or OAM label \cite{belinfante1940}. Circular polarization fixes the handedness of the transverse \(\E\)--\(\H\) triad locally; outward angular flux still requires a signed stress integral on an enclosing surface \cite{stratton1941,beth1936}. Helical phase topology is separated from radiated action below \cite{allen1992}.

\subsection{Helical phase structure and transport}
\label{sec:ch01.2.2}

Helical isophases on a transverse plane record the winding of $\arg E$. Outward transport is a different statement: the signed integral of the stress-moment density over a closed shell, sketched in Figure~\ref{fig:ch1.phase-vs-transport-sketch}. A polarization converter can alter $\oint \Mavg\cdot\hat{\mathbf{n}}\,dS$ on $\Gamma_R$ while a near-zone phase sketch looks unchanged \cite{allen1992,beth1936}; Subsection~\ref{sec:ch01.6.2} gives counter-rotating examples where the phase portrait stays structured but the net flux cancels.

\begin{figure}[t]
\centering
\ChOneFigBlock{%
\begin{tikzpicture}[ch1geo]
 \begin{scope}[shift={(-2.1,0)}]
 \draw[ch1axis] (-0.95,0) -- (0.95,0) node[right, font=\scriptsize] {$x$};
 \draw[ch1axis] (0,-0.95) -- (0,0.95) node[above, font=\scriptsize] {$y$};
 \draw[ch1shellO, dotted] (0,0) circle (0.12);
 \foreach \r in {0.28,0.48,0.68} {
   \draw[gray!55!black, line width=0.55pt] (0,0) circle (\r);
 }
 \draw[-{Latex[length=1.5mm,width=1mm]}, gray!45!black, line width=0.5pt]
   (0.75,0) arc [start angle=0, end angle=35, radius=0.75];
 \draw[-{Latex[length=1.5mm,width=1mm]}, gray!50!black, line width=0.5pt]
   ({0.75*cos(60)},{0.75*sin(60)}) arc [start angle=60, end angle=95, radius=0.75];
 \draw[-{Latex[length=1.5mm,width=1mm]}, gray!50!black, line width=0.5pt]
   ({0.75*cos(120)},{0.75*sin(120)}) arc [start angle=120, end angle=155, radius=0.75];
 \draw[-{Latex[length=1.5mm,width=1mm]}, gray!50!black, line width=0.5pt]
   ({0.75*cos(180)},{0.75*sin(180)}) arc [start angle=180, end angle=215, radius=0.75];
 \draw[-{Latex[length=1.5mm,width=1mm]}, gray!50!black, line width=0.5pt]
   ({0.75*cos(240)},{0.75*sin(240)}) arc [start angle=240, end angle=275, radius=0.75];
 \draw[-{Latex[length=1.5mm,width=1mm]}, gray!50!black, line width=0.5pt]
   ({0.75*cos(300)},{0.75*sin(300)}) arc [start angle=300, end angle=335, radius=0.75];
 \node[ch1oncanvas, font=\scriptsize, align=center] at (0,-1.15)
 {winding of $\arg E$};
 \end{scope}
 \begin{scope}[shift={(2.1,0)}]
 \draw[ch1shell] (0,0) circle (0.72);
 \fill (0,0) circle (1.5pt);
 \draw[-{Latex[length=2mm,width=1.3mm]}, gray!45!black, line width=0.65pt]
   (0.72,0) -- (1.05,0);
 \draw[-{Latex[length=2mm,width=1.3mm]}, gray!65!black, line width=0.65pt]
   (0,0.72) -- (0,1.05);
 \node[ch1oncanvas, font=\scriptsize] at (1.15,0) {$+\mathbf{j}_J$};
 \node[ch1oncanvas, font=\scriptsize] at (0,1.18) {$-\mathbf{j}_J$};
 \node[ch1oncanvas, font=\scriptsize] at (0.95,-0.55) {$\Gamma_R$};
 \node[ch1oncanvas, font=\scriptsize, align=center] at (0,-1.15)
 {$\boldsymbol{\Phi}_{J,\mathrm{out}}$};
 \end{scope}
\end{tikzpicture}%
}
\caption{Phase topology on a transverse plane (left) and angular-flux transport on a control surface $\Gamma_R$ (right). Helical isophases show the winding of $\arg E$; angular-momentum transport is established only when direction-resolved $\mathbf{j}_J$ integrates to nonzero $\boldsymbol{\Phi}_{J,\mathrm{out}}$ without cancellation between counter-rotating sectors \cite{allen1992,beth1936}.}
\label{fig:ch1.phase-vs-transport-sketch}
\end{figure}

Returning to the boundary-value data, Allen \emph{et al.} \cite{allen1992} identified paraxial modes; the following is a waist-plane paraxial modal factor with the outgoing axial phase retained for the present discussion, not the full propagated Laguerre--Gaussian solution with $w(z)$, wavefront curvature, and Gouy phase \cite{allen1992}.
\begin{equation}
\label{eq:ch1.phase-winding-number}E_\phi(\rho,\phi,z)
=
E_0
\left(\frac{\rho}{w_0}\right)^{|m|}
L_p^{|m|}\!\left(\frac{2\rho^2}{w_0^2}\right)
e^{-\rho^2/w_0^2}\,
e^{\jj(m\phi-k_0 z)},
\end{equation}
\addtocounter{equation}{-1}\refstepcounter{equation}\label{eq:ch1.lg-vortex-field}%
where \(L_p^{|m|}\) is an associated Laguerre polynomial, \(w_0\) is beam waist,
and integer \(m\) is azimuthal index. For \(p=0\) and \(|m|=1\), the intensity vanishes
on the axis (\(\rho=0\)), while the phase advances by \(2\pi m\) in one circuit about the axis.
Phase topology is the winding structure of the transverse phase; transport is the outward real-power and angular-flux integral on \(\Gamma\), and topology can persist while transport cancels on closed surfaces. Equation~\eqref{eq:ch1.lg-vortex-field} supplies only a transverse electric template in the modal class of Allen \emph{et al.}; it is not, by itself, a Maxwell-valid vector field because $\H$, longitudinal components, and curl constraints are not displayed \cite{allen1992,stratton1941}. Net structured angular-momentum transport requires a nonzero signed solid-angle integral built from a completed $(\E,\H)$ pair: superposing the same outgoing modal factor of Equation~\eqref{eq:ch1.lg-vortex-field} with the opposite winding $m\to -m$ produces the same standing pattern as Equation~\eqref{eq:ch1.ex62-counter-rotating-superposition}, developed in Subsection~\ref{sec:ch01.6.2}. Any statement about azimuthal Poynting flow from Eq.~\eqref{eq:ch1.lg-vortex-field} alone is therefore deferred until first-order paraxial Maxwell completion is carried out (Chapter~\ref{ch:02}) \cite{allen1992,harrington2001}.

For a uniform transverse plane wave, $\H^{(0)}=(\hat{\mathbf{k}}\times\E^{(0)})/\eta_0$ fixes axial power transport through $\Re(\mathbf{S}_\mathrm{c})\cdot\hat{\mathbf{k}}$. That leading-order relation does not extend pointwise to the spatially varying template of Eq.~\eqref{eq:ch1.lg-vortex-field} without deriving the corresponding $\H^{(0)}$, $E_z$, $H_z$, and verifying $\curl\E$ and $\curl\H$ to the stated paraxial order \cite{stratton1941,allen1992}. Phase winding can accompany reactive $\Im(\mathbf{S}_\mathrm{c})$ exchange in near zones, but signed outward angular-momentum transport is fixed only by the stress integral on the stated enclosing surface when opposing solid-angle sectors need not cancel \cite{allen1992,poynting1884,stratton1941}.

Poynting circulation can be locally strong while the outward real-power integral on a surrounding shell remains small or unchanged: the complex Poynting vector separates reactive exchange through $\Im\{\mathbf{S}_\mathrm{c}\cdot\hat{\mathbf{n}}\}$ from time-averaged radiation through $\Re\{\mathbf{S}_\mathrm{c}\cdot\hat{\mathbf{n}}\}$ \cite{poynting1884}. Poynting's 1884 memoir introduced the energy-flux vector $\mathbf{s}=\E\times\H$, the energy density in the field, and the mechanical account of energy transfer to and from conductors through the surrounding electromagnetic medium \cite{poynting1884}. Its flow-line pictures visualize \emph{energy} transport and storage exchange---not the Maxwell-stress angular-momentum flux $\oint_{\partial V}\Mavg\cdot\hat{\mathbf{n}}\,dS$, which is derived later from symmetric stress and momentum balance \cite{maxwell1873,stratton1941,beth1936}. Locally, circulation in $\mathbf{P}=\langle\mathbf{S}\rangle=\Re\{\mathbf{S}_\mathrm{c}\}$ can accompany azimuthal phase structure and reactive coupling through $\Im(\mathbf{S}_\mathrm{c})$ without fixing signed angular-momentum transport; the latter requires the stress integral on the stated enclosing surface (subsection~\ref{sec:ch01.2.3}) \cite{stratton1941}.

In this setting, directional dependence tracks radiation geometry on $\Gamma_R$: the far-zone pattern may exhibit strong angular variation while source and field angular-momentum balances remain distinct objects tied to impressed currents and flux integrals \cite{belinfante1940}.
\subsection{Physical relations for angular-momentum transport}
\label{sec:ch01.2.3}

Azimuthal phase structure and radiation intensity $U_r$ answer different questions: net outward angular-momentum transport follows from the Maxwell-stress balance and
\begin{equation}
\label{eq:ch1.am-net-transport-condition}
\boldsymbol{\Phi}_{J,\mathrm{out}}(r)=-\oint_{S_r}\mathbf{r}\times(\Tavg\cdot\hat{\mathbf{n}})\,dS
=-\int_{4\pi}\mathbf{j}_J(\theta,\phi;r)\,d\Omega,
\end{equation}
and the axis-projected scalar transport is
\begin{equation}
\label{eq:ch1.am-net-transport-z}
\Phi_{J,z,\mathrm{out}}(r)=\hat{\mathbf{z}}\cdot\boldsymbol{\Phi}_{J,\mathrm{out}}(r)
=-\int_{4\pi}\hat{\mathbf{z}}\cdot\mathbf{j}_J(\theta,\phi;r)\,d\Omega
=\dot J_{z,\mathrm{em,out}}(r).
\end{equation}
with $\mathbf{j}_J=r^2[\mathbf{r}\times(\Tavg\cdot\hat{\mathbf{r}})]$ from Equation~\eqref{eq:ch1.directional-am-flux} \cite{stratton1941,beth1936}. Equal counter-rotating components with $P_+=P_-$ give $\dot J_{z,\mathrm{orb,out}}=(m/\omega)(P_+-P_-)=0$ for the declared orbital modal bookkeeping scale while $|E_t|\propto\sqrt{2}\,|A(\rho)|\,|\cos(m\phi)|$ may remain large (Subsection~\ref{sec:ch01.6.2}); total Maxwell transport $\dot J_{z,\mathrm{em,out}}$ from Eq.~\eqref{eq:ch1.am-net-transport-z} must still be evaluated on the stated field and surface \cite{allen1992,stratton1941}. Outward power in a sector requires $\int_{\Delta\Omega}U_r\,d\Omega>0$; outward angular momentum requires $\boldsymbol{\Phi}_{J,\mathrm{out}}\neq\mathbf{0}$---neither implies the other \cite{poynting1884}.
\phantomsection\label{eq:ch1.flux-qualified-decision}%
\phantomsection\label{eq:ch1.transport-charge-continuous-index}%
\phantomsection\label{tab:ch1.triad-numeric-dipole}

\subsection{Reactive and radiative field regimes}
\label{sec:ch01.2.4}

The field of a finite electric dipole exhibits the classical $1/r^3$, $1/r^2$, and $1/r$ radial hierarchy in its exact exterior solution \cite{stratton1941}. In particular, higher multipoles possess their own higher-order near-field radial structure, while every radiating spherical multipole carries an outgoing $1/r$ branch at infinity \cite{stratton1941,hansen1988}. Near the source the dipole $1/r^3$ and $1/r^2$ terms may be comparable with, or larger than, the $1/r$ radiation field. The near zone can therefore contain substantial electric--magnetic energy exchange and complicated phase structure while the outward real power through an enclosing sphere already equals the total radiated power of the source.

\begin{figure}[t]
\centering
\VolIPhysFig{\VolIPhysRadialFieldZones}
\caption{Radial organization of a localized radiator: $r^{-3}$ electrostatic storage, $r^{-2}$ induction, and $r^{-1}$ radiation dominate at increasing distance. The boundaries are qualitative; all three terms coexist at every radius, but their relative weight shifts with $kr$ \cite{stratton1941,poynting1884}.}
\label{fig:ch1.radial-field-zones}
\end{figure}

On the same Maxwell pair, figure~\ref{fig:ch1.radial-field-zones} matches the multipole counting developed in Subsection~\ref{sec:ch01.6.4}: moving outward does not create new radiated power, but changes which radial term carries the local Poynting and stress flux.

At increasing distance the $1/r$ terms dominate. The electric and magnetic fields then approach the transverse relation characteristic of an outgoing wave, and the real Poynting flux becomes the principal carrier of energy away from the source \cite{poynting1884}. Figure~\ref{fig:ch1.wp-reactive-vs-radiative} illustrates this shell dependence: reactive circulation in $\Im\{\mathbf{S}_\mathrm{c}\cdot\hat{\mathbf{n}}\}$ can be large on a small-$kr$ surface even when the net outward real power is unchanged, whereas on a large-$kr$ surface the same source supports a field whose local phase relation is predominantly radiative.

The distinction is evaluated on a spherical observation surface $S_r$ of radius $r$ with outward unit normal $\hat{\mathbf{n}}$. Define
\begin{equation}
\label{eq:ch1.ex64-power-functionals}
P_{\mathrm{rad}}=\int_{S_r}\Re\!\left\{\mathbf{S}_\mathrm{c}\cdot\hat{\mathbf{n}}\right\}dS,\qquad
P_{\mathrm{reac}}(r)=\int_{S_r}\Im\!\left\{\mathbf{S}_\mathrm{c}\cdot\hat{\mathbf{n}}\right\}dS,
\end{equation}
where \(\mathbf{S}_\mathrm{c}=\tfrac{1}{2}\E\times\H^*\) is the complex Poynting vector
\cite{poynting1884}. In a lossless source-free region between enclosing surfaces, $P_{\mathrm{rad}}$ is independent of $r$ \cite{stratton1941}: it is the total time-averaged power radiated by the source. The integral $P_{\mathrm{reac}}(r)$ measures reactive power associated with local oscillatory exchange between electric and magnetic storage on the chosen shell; $|P_{\mathrm{reac}}(r)|$ generally varies with $r$ even when $P_{\mathrm{rad}}$ does not.

The real and imaginary parts of the complex Poynting vector distinguish active power transport from reactive harmonic exchange because they weight different phase relationships between \(\E\) and \(\H\). On a reactive shell, \(\Im\{\mathbf{S}_\mathrm{c}\cdot\hat{\mathbf{n}}\}\) can integrate to a large \(|P_{\mathrm{reac}}(r)|\) while \(P_{\mathrm{rad}}\) remains fixed; on a radiative shell the ratio \(|P_{\mathrm{reac}}(r)|/P_{\mathrm{rad}}\) is typically small even though the same $P_{\mathrm{rad}}$ crosses every enclosing sphere \cite{poynting1884,stratton1941}.

Within the harmonic formulation, near the source, with \(kr\ll 1\), the \(r^{-3}\) and \(r^{-2}\) multipole terms dominate and \(|P_{\mathrm{reac}}(r)|/P_{\mathrm{rad}}\) is typically much greater than unity: reactive circulation dominates the normal flux pattern on that shell even though \(P_{\mathrm{rad}}\) is unchanged \cite{stratton1941}. The scale \(kr\sim 1\) marks a transition region in which no single radial power law fixes the shell integrals: the ratio \(|P_{\mathrm{reac}}(r)|/P_{\mathrm{rad}}\) must be evaluated for the actual field on \(S_r\), not inferred from \(kr\) alone \cite{stratton1941}. When \(kr\gg 1\) and the \(r^{-1}\) radiation field dominates, \(|P_{\mathrm{reac}}(r)|/P_{\mathrm{rad}}\ll 1\) typically holds and the field on \(S_r\) is predominantly radiative \cite{poynting1884,stratton1941}. Any angular-momentum interpretation that ignores shell radius therefore risks
attributing storage-dominated structure to propagative torque flux
\cite{stratton1941}. Reactive shells carry standing-wave structure whose phase rotates
locally while time-averaged Poynting flux through the same surface may vanish
\cite{stratton1941}. The same small electric dipole at $kr=0.5$ versus $kr=5$ illustrates why
shell placement matters: the inner shell is storage-dominated with \(|P_{\mathrm{reac}}(r)|/P_{\mathrm{rad}}\gg 1\); the outer shell carries \(|P_{\mathrm{reac}}(r)|/P_{\mathrm{rad}}\ll 1\) and exhibits a predominantly radiative field \cite{harrington2001}. Subsection~\ref{sec:ch01.6.4} develops the two-shell comparison on the same dipole using these integrals.

\section{Rotational covariance and reference frames}
\label{sec:ch01.3}

Section~\ref{sec:ch01.2} evaluated fluxes in a fixed laboratory frame about a declared origin. One must now separate two operations: rotating the physical source, currents, and boundaries while the axes stay fixed, versus holding the excitation fixed and rotating only the coordinate labels used to describe it \cite{stratton1941,hansen1988}. The first changes the Maxwell problem when material and impressed data rotate together; the second relists components of the same field. Under proper rotations, $\E$, $\H$, and the Maxwell stress transform as vectors and tensors so that traction and stress-moment integrals rotate consistently. Gauge transformations of the potentials are an additional non-uniqueness, treated in section~\ref{sec:ch01.4} \cite{belinfante1940,maxwell1873}.

\phantomsection\label{sec:ch01.3.1}

\subsubsection*{Active and passive rotations; transformation of vectors, tensors, and dyads}

\begin{figure}[t]
\centering
\ChOneFigBlock{%
\begin{tikzpicture}[ch1geo]
 \node[ch1oncanvas, font=\scriptsize] at (-1.55,1.45) {(a) source rotation};
 \draw[-{Latex[length=1.8mm,width=1.1mm]}, gray!60!black, line width=0.55pt]
   (-0.15,-1.1) -- (-0.15,1.1);
 \draw[-{Latex[length=1.8mm,width=1.1mm]}, gray!60!black, line width=0.55pt]
   (-1.1,-0.15) -- (1.1,-0.15);
 \node[ch1oncanvas, font=\scriptsize] at (0.05,1.05) {$\hat{\mathbf{z}}$};
 \node[ch1oncanvas, font=\scriptsize] at (1.0,-0.05) {$\hat{\mathbf{x}}$};
 \draw[-{Latex[length=2mm,width=1.3mm]}, gray!40!black, line width=0.65pt]
   (0,0) -- (0.55,0.75);
 \draw[gray!40!black, line width=0.55pt] (0,0) circle (3pt);
 \node[ch1oncanvas, font=\scriptsize] at (0.35,0.45) {$\Jimp$};
 \draw[gray!50!black, dashed, line width=0.5pt] (0,0) circle (0.85);
 \draw[-{Latex[length=1.5mm,width=1mm]}, gray!60!black, line width=0.5pt]
   (0.72,0.35) -- (1.0,0.55);
 \node[ch1oncanvas, font=\scriptsize] at (0.95,0.7) {$\E$};
 \node[ch1oncanvas, font=\scriptsize, align=center] at (-0.15,-1.5) {fixed reference\\axes};

 \node[ch1oncanvas, font=\scriptsize] at (2.45,1.45) {(b) chart rotation};
 \draw[-{Latex[length=1.8mm,width=1.1mm]}, gray!60!black, line width=0.55pt]
   (2.35,-1.1) -- (2.35,1.1);
 \draw[-{Latex[length=1.8mm,width=1.1mm]}, gray!55!black, line width=0.55pt]
   (1.5,0.35) -- (3.3,1.15);
 \node[ch1oncanvas, font=\scriptsize] at (3.05,1.25) {$\hat{\mathbf{x}}'$};
 \draw[-{Latex[length=1.8mm,width=1.1mm]}, gray!55!black, line width=0.55pt]
   (2.05,-0.35) -- (2.85,-0.95);
 \node[ch1oncanvas, font=\scriptsize] at (2.55,-1.05) {$\hat{\mathbf{y}}'$};
 \draw[-{Latex[length=2mm,width=1.3mm]}, gray!40!black, line width=0.65pt]
   (2.35,0) -- (2.9,0);
 \draw[gray!40!black, line width=0.55pt] (2.35,0) circle (3pt);
 \node[ch1oncanvas, font=\scriptsize] at (2.75,0.2) {$\Jimp$};
 \draw[gray!50!black, dashed, line width=0.5pt] (2.35,0) circle (0.85);
 \draw[-{Latex[length=1.5mm,width=1mm]}, gray!60!black, line width=0.5pt]
   (3.07,0) -- (3.35,0.2);
 \node[ch1oncanvas, font=\scriptsize] at (3.45,0.25) {$\E$};
 \node[ch1oncanvas, font=\scriptsize, align=center] at (2.35,-1.5) {same source,\\relabeled axes};
\end{tikzpicture}%
}
\caption{Physical source rotation in a fixed reference frame (a) versus passive relabeling of coordinate axes while the impressed current remains fixed (b). Maxwell radiation changes in (a) when source, medium, and boundary do not share rotational symmetry; passive relabeling in (b) reorders field components without moving the source \cite{hansen1988,stratton1941}.}
\label{fig:ch1.source-vs-chart-rotation}
\end{figure}

Figure~\ref{fig:ch1.source-vs-chart-rotation} keeps apart three transformations that the literature often treats as interchangeable. Panel~(a) rotates the impressed current while the reference triad stays fixed; panel~(b) holds the source fixed and relabels the component chart. A third circumstance---allowing azimuthal phase to advance past a stationary line of sight without moving either source or axes---can produce angular \emph{appearance} in a single frame; that circumstance need not coincide with (a) or (b) \cite{hansen1988}.

For panel~(a), whenever source, medium, and boundaries share rotational symmetry about the chosen axis, one may rotate the impressed current and thereafter solve Maxwell's equations, or solve first and rotate the field---the two procedures agree on the physical field \cite{stratton1941,hansen1988}. An anisotropic or inhomogeneous medium violates that commutativity only when constitutive and boundary data fail to transform covariantly under the selected rotation \cite{stratton1941}.

Passive relabeling in~(b) touches only the component chart. Traction $\mathbf{t}=\Tavg\cdot\hat{\mathbf{n}}$ and its moment about a declared origin are assembled from $\Tavg$ (Eq.~\eqref{eq:ch1.maxwell-stress}). Vectors carry one factor of the rotation matrix; the stress tensor carries two, $\Tavg'=R\Tavg R^\top$ \cite{beth1936,stratton1941}.

Let $R\in\mathrm{SO}(3)$ denote a proper rotation, so that $R^\top R=I$ and $\det R=1$. Write a vector as $\mathbf{v}=v_j\hat{\mathbf{e}}_j$ in an orthonormal basis $\hat{\mathbf{e}}_j$. Under the active convention adopted here, a physical rotation gives $\mathbf{v}'_{\mathrm{active}}=R\mathbf{v}$; under a passive rotation of the basis by the same $R$, components obey $\mathbf{v}'_{\mathrm{passive}}=R^\top\mathbf{v}$. The two records describe one geometric object when active and passive maps are inverse. Eq.~\eqref{eq:ch1.vector-rotation-components} records the active component law
\begin{equation}
\label{eq:ch1.infinitesimal-active-rotation}v_i' = R_{ij}v_j.
\end{equation}
\addtocounter{equation}{-1}\refstepcounter{equation}\label{eq:ch1.vector-rotation-components}%
\addtocounter{equation}{-1}\refstepcounter{equation}\label{eq:ch1.active-rotation-action}%
A finite rotation through angle \(\varphi\) about the unit axis \(\hat{\mathbf{u}}\) is expressed by the active Rodrigues formula
\begin{equation}
\label{eq:ch1.rodrigues-rotation}
R(\varphi,\hat{\mathbf{u}})\,\mathbf{v}
=
\mathbf{v}\cos\varphi
+(\hat{\mathbf{u}}\times\mathbf{v})\sin\varphi
+\hat{\mathbf{u}}(\hat{\mathbf{u}}\cdot\mathbf{v})(1-\cos\varphi),
\end{equation}
which agrees, on vector fields, with the generator exponential introduced below and carried forward in Chapter~\ref{ch:03} \cite{jackson1999,stratton1941}.

One must not identify rotation of a field with rotation of its component chart alone. For $R\in\mathrm{SO}(3)$, the active map on a scalar $F$ and on a vector field $\mathbf{F}$ reads
\begin{equation}
\label{eq:ch1.field-rotation-operator}
(T_R F)(\mathbf{r})=F(R^{-1}\mathbf{r}),
\qquad
(T_R\mathbf{F})(\mathbf{r})=R\,\mathbf{F}(R^{-1}\mathbf{r}),
\end{equation}
so that sample point and vector direction rotate coherently \cite{jackson1999}. Evaluating $\E(\mathbf{r})$ after an active rotation is not the same as multiplying $\E(\mathbf{r})$ by $R$ at fixed $\mathbf{r}$, except when only the chart has been relabeled \cite{stratton1941}. Source-free vacuum Maxwell equations are covariant under $T_R$ because $\curl$ and $\divg$ commute with the spatial part of the action whenever $\varepsilon$ and $\mu$ are scalars \cite{maxwell1873}; impressed sources generally break rotate-then-solve equivalence, and an anisotropic medium does so whenever constitutive and boundary data lack the stated symmetry---even though the transformation law itself stays valid \cite{stratton1941}.

Following an active source rotation, the same field strength $|E|$ on $S_R$ appears in relabeled $\hat{\mathbf{x}},\hat{\mathbf{y}}$ components \cite{hansen1988}. An active rotation of the \emph{source} (Subsection~\ref{sec:ch01.0.2}) modifies $\Jimp$ and can change the radiated field when symmetry is lacking; a passive rotation of the \emph{chart}, with the source fixed, permutes components without displacing the impressed current \cite{stratton1941}.

The Maxwell stress tensor couples $\mathbf{T}$, $\hat{\mathbf{n}}$, and $\mathbf{t}$: each column of $\mathbf{T}$ gives the traction on a face whose outward normal aligns with a coordinate direction. As a geometric object that maps normals to force densities, $\mathbf{T}$ must inherit \emph{two} factors of $R$, one per tensor index:
\begin{equation}
\label{eq:ch1.passive-rotation-action}\Tavg' = R\,\Tavg\,R^\top,
\qquad
T'_{ij}=R_{ip}R_{jq}T_{pq}.
\end{equation}
\addtocounter{equation}{-1}\refstepcounter{equation}\label{eq:ch1.tensor-rotation-rule}%
\addtocounter{equation}{-1}\refstepcounter{equation}\label{eq:ch1.vector-observable-covariance}%
The factors $R_{ip}$ and $R_{jq}$ rotate the two tensor slots in concert. Eq.~\eqref{eq:ch1.tensor-rotation-rule} guarantees that the stress tensor of Eq.~\eqref{eq:ch1.maxwell-stress} transforms as one object acting on normals, rather than as nine detached scalars.
Treating $T_{ij}$ as independently rotating entries---a common slip in stress-integral work---introduces spurious torque when results are compared with mechanical reaction on a polarizable body \cite{beth1936}.

When the reference triad rotates, the oriented surface patch rotates with it; traction on the patch remains $\mathbf{t}=\Tavg\cdot\hat{\mathbf{n}}$, and
\begin{equation}
\label{eq:ch1.traction-rotation-covariance}
\mathbf{t}'=\Tavg'\cdot\hat{\mathbf{n}}' = R\mathbf{t}.
\end{equation}
Eq.~\eqref{eq:ch1.traction-rotation-covariance} expresses the mechanical content: the force density on the rotated element is the same physical traction, now written in the rotated basis \cite{beth1936}. In vacuum the cycle-averaged Maxwell stress is symmetric, $\Tavg=\Tavg^\top$, for every $R\in\mathrm{SO}(3)$ \cite{maxwell1873}.

When the field is axially symmetric about \(\hat{\mathbf{z}}\), azimuthal dependence factorizes as \(\psi(\phi)=e^{\jj m\phi}\), and the orbital angular index obeys
\begin{equation}
\label{eq:ch1.azimuthal-eigenvalue}
-\jj\,\frac{\partial\psi}{\partial\phi}=m\,\psi.
\end{equation}
Eq.~\eqref{eq:ch1.azimuthal-eigenvalue} labels scalar azimuthal phase; it is not, without further structure, the vector angular-momentum operator on $(\E,\H)$ \cite{stratton1941}. Meaningful torque comparisons on matter demand an explicit record of whether $R$ acted on the source, on the chart, or on both, together with a fixed reference origin \cite{beth1936,stratton1941}.

\phantomsection\label{eq:ch1.active-passive-duality}
\phantomsection\label{eq:ch1.tensor-invariant-trace}
\phantomsection\label{eq:ch1.stress-eigen-covariance}
\phantomsection\label{eq:ch1.rotation-group-action}
\phantomsection\label{eq:ch1.dyadic-covariance}
\phantomsection\label{eq:ch1.rank1-dyad-rotation}
\phantomsection\label{sec:ch01.3.2b}
\phantomsection\label{eq:ch1.total-angular-generator}
\phantomsection\label{eq:ch1.orbital-angular-generator}
\phantomsection\label{eq:ch1.spin-angular-generator}
\phantomsection\label{eq:ch1.rotation-from-generators}

\phantomsection\label{sec:ch01.3.3}
\subsubsection*{Origin dependence}

Torque is invariably torque \emph{about a point}. Shifting the reference origin from $\mathbf{r}_0$ to $\mathbf{r}_0'=\mathbf{r}_0+\Delta\mathbf{r}$ leaves net force unchanged yet adds the transport couple $\boldsymbol{\tau}_{\mathbf{r}_0'}=\boldsymbol{\tau}_{\mathbf{r}_0}-\Delta\mathbf{r}\times\mathbf{F}$, with $\mathbf{F}=\oint_{S_m}\Tavg\cdot\hat{\mathbf{n}}\,dS$, even when $(\E,\H)$ and $\mathbf{T}$ on $S_m$ are held fixed \cite{beth1936,stratton1941}. Figure~\ref{fig:ch1.origin-shift-torque} separates force, which follows from traction alone, from torque, which inherits the moment arm through the origin. The angular-momentum flux $\Mavg=\mathbf{r}\times\Tavg$ obeys the same rule \cite{stratton1941}.

\begin{figure}[t]
\centering
\ChOneFigBlock{%
\begin{tikzpicture}[ch1geo]
 \coordinate (O) at (0,0);
 \coordinate (Op) at (0.55,0.35);
 \coordinate (P) at (1.35,0.55);
 \fill (O) circle (2pt);
 \fill (Op) circle (2pt);
 \fill (P) circle (2pt);
 \node[ch1oncanvas, font=\scriptsize, anchor=east] at (O) {$O$};
 \node[ch1oncanvas, font=\scriptsize, anchor=west] at (Op) {$O'$};
 \node[ch1oncanvas, font=\scriptsize, anchor=west] at (P) {$P$};
 \draw[gray!55!black, line width=0.45pt] (-0.35,-0.55) -- (1.65,-0.55);
 \node[ch1oncanvas, font=\scriptsize] at (0.65,-0.75) {material surface $S_m$};
 \draw[-{Latex[length=2mm,width=1.3mm]}, gray!50!black, line width=0.6pt]
   (P) -- (1.75,0.55);
 \node[ch1oncanvas, font=\scriptsize] at (1.95,0.55) {$\mathbf{t}$};
 \draw[-{Latex[length=1.8mm,width=1.1mm]}, gray!60!black, line width=0.55pt]
   (O) -- (P);
 \node[ch1oncanvas, font=\scriptsize] at (0.75,0.15) {$\mathbf{r}$};
 \draw[-{Latex[length=1.8mm,width=1.1mm]}, gray!55!black, line width=0.55pt]
   (Op) -- (P);
 \node[ch1oncanvas, font=\scriptsize] at (1.05,0.05) {$\mathbf{r}'$};
 \draw[-{Latex[length=1.6mm,width=1mm]}, gray!40!black, line width=0.55pt, dashed]
   (O) -- (Op);
 \node[ch1oncanvas, font=\scriptsize] at (0.15,0.25) {$\Delta\mathbf{r}$};
 \draw[-{Latex[length=1.8mm,width=1.1mm]}, gray!45!black, line width=0.55pt]
   (O) -- (-0.15,0.55);
 \node[ch1oncanvas, font=\scriptsize] at (-0.35,0.65) {$\boldsymbol{\tau}_O$};
 \draw[-{Latex[length=1.8mm,width=1.1mm]}, gray!45!black, line width=0.55pt]
   (Op) -- (0.95,0.95);
 \node[ch1oncanvas, font=\scriptsize] at (1.15,1.05) {$\boldsymbol{\tau}_{O'}$};
\end{tikzpicture}%
}
\caption{Force density $\mathbf{t}$ at point $P$ on $S_m$ contributes to resultant force independent of origin, but torque $\boldsymbol{\tau}_O$ about $O$ differs from $\boldsymbol{\tau}_{O'}$ about $O'=O+\Delta\mathbf{r}$ because the moment arm changes \cite{beth1936}.}
\label{fig:ch1.origin-shift-torque}
\end{figure}

Active rotation of the impressed current, constitutive data, and boundaries defines a new Maxwell problem when those data lack full rotational symmetry; passive relabeling of the coordinate chart records the same physical field in new components. Stress traction and stress-moment flux inherit vector and tensor transformation laws, and both mechanical torque and $\Mavg=\mathbf{r}\times\Tavg$ demand an explicit reference origin. These covariance statements close the kinematic thread opened in section~\ref{sec:ch01.2}; the remaining non-uniqueness of the potential description is taken up in section~\ref{sec:ch01.4}.

\phantomsection\label{sec:ch01.3.4}

\FloatBarrier
\section{Gauge structure, helicity, and spin--orbit limits}
\label{sec:ch01.4}

Rotations fix how $(\E,\H)$ and the stress tensor appear in a new frame, but they do not exhaust the redundancies in the potential description. This section treats gauge freedom, the field observables that survive it, helicity as a property of the Maxwell pair, dual symmetry where it applies, and the limits of assigning unique local spin and orbital parts to angular momentum \cite{belinfante1940,jackson1999,stratton1941}.

\phantomsection\label{sec:ch01.4.2}
\phantomsection\label{prin:ch1.angular-em-structure}

\subsubsection*{Gauge freedom and angular structure}

Once rotation and origin have been declared, the mechanical reports built from $(\E,\H)$ are objective. Potentials raise a distinct issue that rotation does not resolve: may $(\mathbf{A},\phi)$ be altered while $(\E,\H)$ stay fixed? On a compact range two reconstructions can disagree on $\mathbf{A}(\mathbf{r})$ at a sample point yet yield identical tangential $\E$ and $\H$ on the scan surface \cite{stratton1941,hansen1988}. That freedom is not a defect of Maxwell theory; it belongs to the potential representation. Wattmeters, torsion balances, and stress integrals respond to $(\E,\H)$, $\mathbf{S}_\mathrm{c}$, and $\Tavg$, not to $\mathbf{A}$ in isolation \cite{maxwell1873,beth1936}.

\begin{figure}[t]
\centering
\ChOneFigBlock{%
\begin{tikzpicture}[ch1geo]
 \node[ch1oncanvas, font=\scriptsize] at (-1.35,1.35) {(a) $\mathbf{A}_1$};
 \draw[-{Latex[length=1.8mm,width=1.1mm]}, gray!60!black, line width=0.55pt]
   (-0.15,-0.95) -- (-0.15,0.95);
 \draw[-{Latex[length=1.8mm,width=1.1mm]}, gray!60!black, line width=0.55pt]
   (-0.95,-0.15) -- (0.95,-0.15);
 \draw[-{Latex[length=2mm,width=1.3mm]}, gray!60!black, line width=0.65pt]
   (-0.55,0.35) arc (145:35:0.65);
 \draw[-{Latex[length=2mm,width=1.3mm]}, gray!40!black, line width=0.65pt]
   (0.15,0.15) -- (0.75,0.75);
 \node[ch1oncanvas, font=\scriptsize] at (0.55,0.85) {$\E$};
 \node[ch1oncanvas, font=\scriptsize] at (-0.55,0.55) {$\mathbf{A}_1$};

 \node[ch1oncanvas, font=\scriptsize] at (1.65,1.35) {(b) $\mathbf{A}_2=\mathbf{A}_1+\nabla\chi$};
 \draw[-{Latex[length=1.8mm,width=1.1mm]}, gray!60!black, line width=0.55pt]
   (1.35,-0.95) -- (1.35,0.95);
 \draw[-{Latex[length=1.8mm,width=1.1mm]}, gray!60!black, line width=0.55pt]
   (0.55,-0.15) -- (2.45,-0.15);
 \draw[-{Latex[length=2mm,width=1.3mm]}, gray!60!black, line width=0.65pt, dashed]
   (1.05,0.55) .. controls (1.55,0.15) .. (2.05,0.45);
 \draw[-{Latex[length=2mm,width=1.3mm]}, gray!40!black, line width=0.65pt]
   (1.65,0.15) -- (2.25,0.75);
 \node[ch1oncanvas, font=\scriptsize] at (2.05,0.85) {$\E$};
 \node[ch1oncanvas, font=\scriptsize] at (1.15,0.55) {$\mathbf{A}_2$};

 \node[ch1oncanvas, font=\scriptsize, align=center] at (0.55,-1.25)
 {same $\E$, $\H$; different $\mathbf{A}$, $\phi$};
\end{tikzpicture}%
}
\caption{Two potential pairs can describe the same physical field: $\mathbf{A}_2=\mathbf{A}_1+\nabla\chi$ relabels the vector potential while $\E=-\nabla\phi-\jj\omega\mathbf{A}$ and $\H=\curl\mathbf{A}/\mu$ remain unchanged \cite{maxwell1873}.}
\label{fig:ch1.gauge-same-field}
\end{figure}

Because $\mathbf{B}=\mu\H=\curl\mathbf{A}$, adding $\nabla\chi$ to $\mathbf{A}$ and co-adjusting $\phi$ leaves $(\E,\H)$ invariant (Figure~\ref{fig:ch1.gauge-same-field}). The subsection below states the transformation under $\exp(+\jj\omega t)$, lists the field observables that survive it, and separates gauge maps from electromagnetic duality.

\subsection{Gauge-invariant observables, helicity, and dual symmetry}
\label{sec:ch01.4.3}

Take the phasor state $\mathcal{F}=(\E,\H)$ in a linear isotropic region with scalar $\varepsilon$, scalar $\mu$, angular frequency $\omega$, and intrinsic impedance $\eta=\sqrt{\mu/\varepsilon}$. Let $(\mathbf{A},\phi)$ denote a potential pair tied to the physical fields by
\begin{equation}
\label{eq:ch1.potential-field-relations}
\E=-\nabla\phi-\jj\omega\mathbf{A},\qquad \mathbf{B}=\curl\mathbf{A}=\mu\H,
\end{equation}
The vector potential $\mathbf{A}$ is a calculational device, not a field observable: its curl fixes $\mathbf{B}$, and its coupling with $\phi$ fixes $\E$. Since $\curl\nabla\chi=\mathbf{0}$, a gradient shift in $\mathbf{A}$ leaves $\mathbf{B}$ unchanged; compensating $\phi$ by $-\jj\omega\chi$ restores $\E$:
\begin{equation}
\label{eq:ch1.gauge-transform}
\mathbf{A}'=\mathbf{A}+\nabla\chi,\qquad \phi'=\phi-\jj\omega\chi,\qquad \E'=\E,\ \H'=\H.
\end{equation}
Direct substitution into $\E=-\nabla\phi-\jj\omega\mathbf{A}$ yields $\E'=\E$ and $\mathbf{B}'=\mathbf{B}$, as in Figure~\ref{fig:ch1.gauge-same-field}. Eq.~\eqref{eq:ch1.gauge-transform} encodes that relabeling. Observables are $\E$, $\B$, $\mathbf{S}_\mathrm{c}$, $\Tavg$, and closed-surface integrals free of explicit potential dependence \cite{maxwell1873,belinfante1940}. \phantomsection\label{eq:ch1.gauge-observable-invariance}\phantomsection\label{eq:ch1.observable-admissibility-triad}\phantomsection\label{eq:ch1.observable-charge-continuous-triad}
Any local map tagged to potentials must state its gauge \cite{belinfante1940}.

Coulomb and Lorenz gauges agree on $(\E,\H)$ but differ on $\phi$ and $\nabla\phi$ in the interior, so a potential-based local map may disagree pointwise between gauges even while traction and angular-flux integrals built from $\Tavg$ and $\Mavg=\mathbf{r}\times\Tavg$ stay fixed on chosen surfaces \cite{belinfante1940,maxwell1873}. For illustration alone, introduce the gauge-linked map below; it exhibits non-uniqueness of potential-based local constructions and is not an electromagnetic angular-momentum density.


\begin{equation}
\label{eq:ch1.representation-linked-jloc}
\mathbf{q}_{\mathrm{loc}}^{(\mathrm{pot})}
=
\frac{\varepsilon_0}{\jj\omega}\,\mathbf{r}\times\nabla\phi,
\end{equation}
Here $\phi$ is the scalar potential in the gauge of Eq.~\eqref{eq:ch1.potential-field-relations}. Under $\phi'=\phi-\jj\omega\chi$, one has $\nabla\phi'=\nabla\phi-\jj\omega\nabla\chi$, so $\mathbf{q}_{\mathrm{loc}}^{(\mathrm{pot})}$ acquires the shift $-\varepsilon_0\,\mathbf{r}\times\nabla\chi$. Eq.~\eqref{eq:ch1.representation-linked-jloc} is thus an orbital map tied to $\phi$, not the gauge-invariant density $\mathbf{j}_{\mathrm{em}}$ and not a Belinfante-class spin density \cite{belinfante1940,stratton1941}.

Inside a control volume $V_g$ of radius $R=2\lambda$ about a dipole, Coulomb and Lorenz gauges coincide on $(\E,\H)$ over $\partial V_g$ yet disagree on $\nabla\phi$ and hence on Eq.~\eqref{eq:ch1.representation-linked-jloc} in the interior; $\oint_{\partial V_g}\Mavg\cdot\hat{\mathbf{n}}\,dS$ is unaffected \cite{belinfante1940,stratton1941}. Helicity records the handedness with which $\E$ and $\H$ wind about $\mathbf{k}$. In a monochromatic circularly polarized plane wave, $\E$ and $\H$ rotate in lockstep while staying transverse to $\mathbf{k}$; right- and left-handed states differ by that winding sense, not by radiated power \cite{stratton1941}. Figure~\ref{fig:ch1.helicity-cp-wave} fixes the geometry before the density is stated.


\begin{figure}[t]
\centering
\ChOneFigBlock{%
\begin{tikzpicture}[ch1geo]
 \draw[-{Latex[length=2mm,width=1.3mm]}, gray!60!black, line width=0.65pt]
   (-1.1,0) -- (1.1,0);
 \node[ch1oncanvas, font=\scriptsize] at (1.25,0) {$\mathbf{k}$};
 \draw[-{Latex[length=1.8mm,width=1.1mm]}, gray!35!black, line width=0.6pt]
   (0,0.75) -- (0,1.05);
 \draw[-{Latex[length=1.8mm,width=1.1mm]}, gray!55!black, line width=0.6pt]
   (0.75,0) -- (1.05,0);
 \draw[gray!35!black, line width=0.55pt] (0,0) circle (0.75);
 \draw[gray!55!black, line width=0.55pt] (0,0) circle (0.55);
 \node[ch1oncanvas, font=\scriptsize] at (-0.95,0.95) {$\E$};
 \node[ch1oncanvas, font=\scriptsize] at (0.95,0.55) {$\H$};
 \node[ch1oncanvas, font=\scriptsize, align=center] at (0,-1.05)
 {same rotation sense $\Rightarrow$ definite helicity};
\end{tikzpicture}%
}
\caption{Circularly polarized plane wave: $\E$ and $\H$ rotate in the same handedness about propagation axis $\mathbf{k}$. Helicity is the field-intrinsic sign of that handedness \cite{stratton1941,belinfante1940}.}
\label{fig:ch1.helicity-cp-wave}
\end{figure}

The vacuum helicity density adopted in this chapter, at a single angular frequency, is
\begin{equation}
\label{eq:ch1.helicity-density}
h(\mathbf{r})=
\frac{1}{2\omega c}\,
\Im\!\left\{\E\cdot\H^*\right\},
\end{equation}
under the $\exp(+\jj\omega t)$ phasor convention \cite{stratton1941,belinfante1940,tai1994}.
Dimensionally, $[\E\cdot\H^*]=\mathrm{W\,m^{-2}}$ while $[h]=\mathrm{J\,s\,m^{-3}}$---angular momentum per unit volume---because $\omega$ enters with $\mathrm{rad\,s^{-1}}$ and $c$ with $\mathrm{m\,s^{-1}}$.
Eq.~\eqref{eq:ch1.helicity-density} must be distinguished from optical chirality (Lipkin zilch) and from Stokes $S_3$: each is a separate pseudoscalar with its own normalization and conservation law \cite{stratton1941,tai1994}. The scalar $h$ reports the handedness of the transverse $\E$--$\H$ triad; local spin--orbit splits become definite only under an explicit convention (section~\ref{sec:ch01.4.5}) \cite{belinfante1940,stratton1941}.

Consider an outgoing circularly polarized plane wave with \emph{real} amplitude $E_0>0$,
$\E=(E_0/\sqrt{2})(\hat{\mathbf{x}}+\jj\sigma\hat{\mathbf{y}})\,e^{-\jj kz}$, $\sigma=\pm1$, $\mathbf{k}=k\hat{\mathbf{z}}$, and right-handed $(\hat{\mathbf{x}},\hat{\mathbf{y}},\hat{\mathbf{z}})$. Faraday's law gives
$\H=(1/\eta_0)\,\hat{\mathbf{z}}\times\E=(E_0/\eta_0\sqrt{2})(\hat{\mathbf{y}}-\jj\sigma\hat{\mathbf{x}})\,e^{-\jj kz}$.
Hence $\E\cdot\H^*=\jj\sigma E_0^2/\eta_0$ and
\begin{equation}
\label{eq:ch1.helicity-plane-wave}
h=\frac{\sigma E_0^2}{2\omega c\eta_0}
=\frac{\sigma\varepsilon_0 E_0^2}{2\omega},
\end{equation}
so that $h$ reverses sign between opposite circular polarizations and vanishes for linear polarization \cite{stratton1941,belinfante1940}. Labels $\sigma=\pm1$ are meaningful only after the spatial triad and time convention are fixed; they cannot be identified with Stokes $S_3$ without an explicit sign map \cite{stratton1941}.

Duality acts on field pairs through
\begin{equation}
\label{eq:ch1.duality-rotation}
\E'=\cos\xi\,\E+\sin\xi\,\eta\H,\qquad
\H'=-\sin\xi\,\E/\eta+\cos\xi\,\H,\qquad \xi\in\mathbb{R}.
\end{equation}
Eq.~\eqref{eq:ch1.duality-rotation} is not a gauge map. In a source-free region with spatially uniform scalar $\varepsilon$ and $\mu$, direct substitution into $\curl\E=-\jj\omega\mu\H$ and $\curl\H=\jj\omega\varepsilon\E$ shows that $(\E',\H')$ obeys the same curl pair provided $\eta=\sqrt{\mu/\varepsilon}$ is the impedance used in Eq.~\eqref{eq:ch1.duality-rotation}. Complex $\varepsilon$ and $\mu$ at fixed $\omega$ are permitted in this \emph{frequency-domain algebraic} sense; a separate time-domain or energy-hermitian statement is not developed here \cite{stratton1941}. Because Maxwell's equations are linear, the map sends any one source-free solution $(\E,\H)$ to another on the same homogeneous constitutive branch; a sum of such solutions is mapped componentwise. The map does not extend to a complete problem with impressed $\rho$ and $\Jimp$ unless sources are transformed consistently, nor across interfaces where $\eta$ jumps unless tangential boundary data are dual-covariant. Spatially varying $\varepsilon(\mathbf r)$, $\mu(\mathbf r)$ with non-uniform $\eta(\mathbf r)$ likewise require a position-dependent impedance in the rotation rule and are outside the uniform-$\eta$ closure of Eq.~\eqref{eq:ch1.duality-rotation} \cite{stratton1941,fernandezcorbaton2013}. At a plane interface between two uniform half-spaces, tangential $\E$ and $\H$ remain continuous, but a global rotation mixing $(\E,\H)$ with a single $\eta$ on both sides does not in general map one physical boundary-value problem into another unless the impedance data are dual-covariant. To first order in $\xi=\delta\xi$, $\delta\E=\delta\xi\,\eta\H$ and $\delta\H=-\delta\xi\,\E/\eta$. For transverse states in a fixed polarization sector, a rotation through $\xi=\pi/2$ exchanges predominantly electric and magnetic roles while Eq.~\eqref{eq:ch1.helicity-density}, evaluated on the physical $(\E,\H)$ pair, continues to report handedness relative to the declared triad and $\exp(+\jj\omega t)$ convention \cite{stratton1941,belinfante1940}.

Because Eq.~\eqref{eq:ch1.helicity-density} depends only on $(\E,\H)$, $h$ is gauge invariant, whereas canonical helicity definitions tied to potentials may move under $\mathbf{A}'=\mathbf{A}+\nabla\chi$ and must not be conflated with Eq.~\eqref{eq:ch1.helicity-density} without naming the formalism \cite{belinfante1940,tai1994,jackson1999}. For circular polarization, $\langle\mathbf{S}\rangle$ has equal magnitude for $\sigma=\pm1$ while $h$ tracks the handedness; linear polarization gives $h=0$ \cite{stratton1941}. This chapter treats helicity as a field-scalar record of handedness and does not develop a separate helicity continuity law, since energy, momentum, and angular momentum already supply the transport structure required for radiation \cite{stratton1941}.


Helicity and duality complete the gauge thread of section~\ref{sec:ch01.4}. Maxwell's equations fix $(\E,\H)$ once sources, constitutive data, boundaries, and outgoing selection are given \cite{maxwell1873,stratton1941}. On any chosen surface, cycle-averaged power follows from $\langle\mathbf{S}\rangle=\Re\mathbf{S}_\mathrm{c}$, mechanical loading from $\Tavg\cdot\hat{\mathbf{n}}$, and angular transport from $\Mavg=\mathbf{r}\times\Tavg$ \cite{poynting1884,beth1936}.

\phantomsection\label{sec:ch01.4.4}
\subsection{Limits of local spin-orbit decompositions}
\label{sec:ch01.4.5}

Total electromagnetic angular momentum is conserved and tied to $(\E,\H)$ through Eq.~\eqref{eq:ch1.angular-momentum-balance}. Helicity, as developed above, supplies a gauge-invariant handedness sign for the pair. Local ``spin'' and ``orbit'' names split $\mathbf{j}_{\mathrm{em}}$ into pointwise pieces that aid interpretation yet are not fixed uniquely by Maxwell's equations alone \cite{belinfante1940,stratton1941}. The total $\mathbf{j}_{\mathrm{em}}$ follows from $(\E,\H)$; convention-dependent $\mathbf{s}_{\mathrm{loc}}$ and $\mathbf{l}_{\mathrm{loc}}$ may reshuffle under gauge while closed-surface flux stays put \cite{belinfante1940,stratton1941}. What is measurable remains $\oint \Mavg\cdot\hat{\mathbf{n}}\,dS$ and the balance of Eq.~\eqref{eq:ch1.angular-momentum-balance-averaged}; the paraxial index $m$ is taken up in Chapter~\ref{ch:02} \cite{allen1992}.

\begin{figure}[!htbp]
\centering
\adjustbox{max width=0.78\linewidth, max height=0.42\textheight, center}{%
\begin{tikzpicture}[ch1geo]
 \draw[-{Latex[length=2mm,width=1.3mm]}, gray!60!black, line width=0.65pt]
   (-1.0,0) -- (1.0,0);
 \node[ch1oncanvas, font=\scriptsize] at (1.15,0) {$\mathbf{k}$};
 \draw[gray!55!black, line width=0.55pt] (0,0) circle (0.55);
 \node[ch1oncanvas, font=\scriptsize] at (0,0.75) {helicity};
 \draw[gray!40!black, dashed, line width=0.5pt]
   (-0.45,-0.35) arc (-140:40:0.55);
 \node[ch1oncanvas, font=\scriptsize] at (-0.55,-0.55) {phase winding};
 \draw[-{Latex[length=2mm,width=1.3mm]}, gray!30!black, line width=0.65pt]
   (0,0) -- (0,0.95);
 \node[ch1oncanvas, font=\scriptsize] at (0.25,1.05) {$\mathbf{j}_{\mathrm{em}}$};
 \draw[-{Latex[length=1.6mm,width=1mm]}, gray!45!black, line width=0.55pt]
   (0.35,0.15) -- (0.75,0.35);
 \node[ch1oncanvas, font=\scriptsize] at (0.95,0.45) {$\mathbf{s}_{\mathrm{loc}}$};
 \draw[-{Latex[length=1.6mm,width=1mm]}, gray!65!black, line width=0.55pt]
   (-0.35,0.15) -- (-0.75,0.35);
 \node[ch1oncanvas, font=\scriptsize] at (-1.05,0.45) {$\mathbf{l}_{\mathrm{loc}}$};
 \node[ch1oncanvas, font=\scriptsize, align=center] at (0,-1.05)
 {chosen split of one total};
\end{tikzpicture}%
}
\caption{Total angular-momentum density $\mathbf{j}_{\mathrm{em}}$ may be partitioned into spin-like $\mathbf{s}_{\mathrm{loc}}$ and orbital-like $\mathbf{l}_{\mathrm{loc}}$ contributions for a structured beam; helicity (polarization handedness) and azimuthal phase winding are distinct physical features, and the local split is a convention, not two separate substances \cite{belinfante1940,allen1992}.}
\label{fig:ch1.spin-orbit-decomposition}
\end{figure}

Formal spin--orbit partitions are useful when their conventions are stated. In a declared paraxial beam they separate polarization-linked circulation from azimuthal phase gradients; in a general vector solution they redistribute the same total $\mathbf{j}_{\mathrm{em}}$ without altering $\oint_{\partial V}\Mavg\cdot\hat{\mathbf{n}}\,dS$ on a chosen shell \cite{belinfante1940,allen1992,stratton1941}. A collimated mode with azimuthal index $m$ is an angular \emph{label} on the chart; signed transport through $\Gamma$ still follows from the Maxwell stress integral with the full $(\E,\H)$ pair \cite{allen1992,stratton1941}.

\FloatBarrier
\section{Surface flux, torque, and physical observation}
\label{sec:ch01.5}

The conservation laws of section~\ref{sec:ch01.1} define fluxes and torques as surface integrals of $\mathbf{S}_\mathrm{c}$, $\Tavg$, and $\Mavg$. Experiment and device modelling require the further step: which tangential samples on $\Gamma$ enter a wattmeter reading, which traction patch loads a mechanical support, and which partial field charts are used only for synthesis \cite{poynting1884,beth1936,stratton1941}. Harmonic coefficients of $\E_t$ aid reconstruction in Chapter~\ref{ch:02}; they are not substitutes for the closed-surface balances \cite{hansen1988,harrington2001}.

\begin{figure}[t]
\centering
\ChOneFigBlock{%
\begin{tikzpicture}[ch1geo]
 \draw[gray!50!black, dashed, line width=0.45pt] (0,0) circle (0.95);
 \node[ch1oncanvas, font=\scriptsize] at (0,1.15) {$\Gamma$};
 \fill[gray!40] (0,0) circle (0.12);
 \node[ch1oncanvas, font=\scriptsize] at (0,-0.28) {source};
 \draw[-{Latex[length=1.8mm,width=1.1mm]}, gray!45!black, line width=0.55pt]
   (0.55,0.35) -- (0.95,0.55);
 \draw[-{Latex[length=1.8mm,width=1.1mm]}, gray!45!black, line width=0.55pt]
   (-0.35,0.65) -- (-0.55,0.95);
 \node[ch1oncanvas, font=\scriptsize] at (1.05,0.65) {$\langle\mathbf{S}\rangle$};
 \draw[gray!55!black, line width=0.5pt] (-0.95,0) -- (-0.55,0);
 \draw[-{Latex[length=1.6mm,width=1mm]}, gray!55!black, line width=0.55pt]
   (-0.75,0) -- (-0.75,0.45);
 \draw[-{Latex[length=1.6mm,width=1mm]}, gray!55!black, line width=0.55pt]
   (-0.75,0) -- (-0.95,0.18);
 \node[ch1oncanvas, font=\scriptsize, align=center] at (-1.05,0.55)
 {$\mathbf{t}$};
 \draw[-{Latex[length=1.4mm,width=0.9mm]}, gray!40!black, line width=0.5pt]
   (-0.75,0) -- (-0.95,-0.35);
 \node[ch1oncanvas, font=\scriptsize] at (-1.05,-0.45) {$\mathbf{r}_0$};
 \draw[-{Latex[length=1.6mm,width=1mm]}, gray!65!black, line width=0.5pt]
   (0.65,-0.55) arc [start angle=-40, end angle=20, radius=0.75];
 \node[ch1oncanvas, font=\scriptsize] at (0.95,-0.35) {$\E_t$};
 \node[ch1oncanvas, font=\scriptsize, align=center] at (0,-1.15)
 {flux, traction, and tangential field on one shell};
\end{tikzpicture}%
}
\caption{Maxwell field content on a closed surface $\Gamma$: outward cycle-averaged power flux $\langle\mathbf{S}\rangle$ enters power balance; stress traction $\mathbf{t}=\Tavg\cdot\hat{\mathbf{n}}$ on a material patch enters force and torque about origin $\mathbf{r}_0$; tangential $\E_t$ is part of the same field pair \cite{poynting1884,beth1936,stratton1941}.}
\label{fig:ch1.measurement-surface-readout}
\end{figure}

Field--matter coupling is the stage on which Maxwell predictions become mechanical fact: traction $\mathbf{t}=\Tavg\cdot\hat{\mathbf{n}}$ operates across conductors, dielectrics, and polarizable bodies, and the attendant force and torque report how absorption or polarization conversion redistributes angular action \cite{beth1936}. Beth's quartz wave plate under circular polarization is the standard demonstration---the plate face supplies $S_m$ for the stress integral, and the measured twist follows stress at the interface, not an orbital label read off a harmonic chart \cite{beth1936,allen1992}.

\subsection{Stress traction, force, and torque on matter}
\label{sec:ch01.5.1}

The cycle-averaged Maxwell stress $\Tavg$ is the local flux of linear momentum \cite{maxwell1873,stratton1941}. At a vacuum--matter interface the traction
\begin{equation}
\label{eq:ch1.stress-traction-definition}
\mathbf{t}=\Tavg\cdot\hat{\mathbf{n}}
\end{equation}
gives the force per unit area that the field exerts on the material face with outward normal $\hat{\mathbf{n}}$. Let $\partial V_b$ denote the \emph{closed} surface that bounds the body, all faces included. The net electromagnetic force is
\begin{equation}
\label{eq:ch1.stress-force-integral}
\mathbf{F}=\oint_{\partial V_b}\Tavg\cdot\hat{\mathbf{n}}\,dS.
\end{equation}
A finite patch $S_m\subset\partial V_b$ contributes only a portion of that integral unless the omitted faces and any edge terms are restored \cite{stratton1941}. Angular momentum is the moment of linear momentum about a chosen origin $\mathbf{r}_0$. The electromagnetic torque is
\begin{equation}
\label{eq:ch1.stress-torque-integral}
\mathbf{T}_{\mathbf{r}_0}=\oint_{\partial V_b}(\mathbf{r}-\mathbf{r}_0)\times(\Tavg\cdot\hat{\mathbf{n}})\,dS.
\end{equation}
The radius vector $\mathbf{r}$ is taken from the same origin as in $\Mavg=\mathbf{r}\times\Tavg$, so torque is the moment of the momentum flux carried by $\Tavg$ \cite{stratton1941,beth1936}. Because moments are taken about an origin, $\boldsymbol{\tau}_{\mathbf{r}_0}$ shifts when $\mathbf{r}_0$ is translated even if $(\E,\H)$ and $\Tavg$ on $S_m$ are unchanged (Eq.~\eqref{eq:ch1.stress-origin-shift-torque}).

When $\mathbf{F}$ and $\boldsymbol{\tau}_{\mathbf{r}_0}$ are fixed by the Maxwell solution on $S_m$, a polarizable body transduces that electromagnetic load into measurable mechanical exchange through its constitutive and elastic response \cite{beth1936,stratton1941}.

\subsection{Power flux, stress traction, and torque}
\label{sec:ch01.5.2}

On a source-enclosing shell $\Gamma$, outward $\Re\{\mathbf{S}_\mathrm{c}\cdot\hat{\mathbf{n}}\}$ integrates to the time-averaged radiated power of Eq.~\eqref{eq:ch1.power-balance-integral}; $\Im\{\mathbf{S}_\mathrm{c}\cdot\hat{\mathbf{n}}\}$ records reactive harmonic exchange on that same shell \cite{poynting1884,stratton1941}. In a lossless exterior, $\int_{\Gamma}\Re\{\mathbf{S}_\mathrm{c}\cdot\hat{\mathbf{n}}\}\,dS$ equals $P_{\mathrm{rad}}$ on every admissible enclosing surface, while $|P_{\mathrm{reac}}(r)|$ generally varies with radius (section~\ref{sec:ch01.2.4}).

Mechanical loading uses the outward normal $\hat{\mathbf{n}}$ of the \emph{material} face, which need not coincide with the field-domain normal on a source-enclosing sphere. Traction $\mathbf{t}=\Tavg\cdot\hat{\mathbf{n}}$ is already defined on $\partial V_b$ in Subsection~\ref{sec:ch01.5.1}; Beth's quartz plate measures torque where that traction couples to an anisotropic interface under circular polarization \cite{beth1936,stratton1941}. Subsection~\ref{sec:ch01.6.1} works the radiation example in full detail.

Comparisons of torque require a shared origin. When the reference point moves from $\mathbf{r}_0$ to $\mathbf{r}_0'=\mathbf{r}_0+\Delta\mathbf{r}$, the deterministic torque integral obeys
\begin{equation}
\label{eq:ch1.stress-origin-shift-torque}
\mathbf{T}_{\mathrm{det}}'
=
\mathbf{T}_{\mathrm{det}}-\Delta\mathbf{r}\times\mathbf{F}_{\mathrm{det}},
\end{equation}
with $\mathbf{F}_{\mathrm{det}}=\int_{S_m}\Tavg\cdot\hat{\mathbf{n}}\,dS$ and $\mathbf{T}_{\mathrm{det}}=\int_{S_m}(\mathbf{r}-\mathbf{r}_0)\times(\Tavg\cdot\hat{\mathbf{n}})\,dS$.
Eq.~\eqref{eq:ch1.stress-origin-shift-torque} is the surface statement of that origin dependence; it must be in force before torque predictions are compared \cite{stratton1941}.

\begin{example}[Why torque depends on the chosen origin]
\label{ex:ch1.origin-shift}
With $\mathbf{r}'=\mathbf{r}-\Delta\mathbf{r}$, stress torque transforms as $\boldsymbol{\tau}'=\boldsymbol{\tau}-\Delta\mathbf{r}\times\mathbf{F}$ (Eq.~\eqref{eq:ch1.stress-origin-shift-torque}). On a thin plate centred at $\mathbf{r}_0$, a shift $\Delta\mathbf{r}=2\,\mathrm{mm}\,\hat{\mathbf{x}}$ at fixed $\mathbf{F}_{\mathrm{det}}$ alters $\tau_z$ by $\Delta\tau_z=-F_y\,\Delta x$ when moments are taken about $\hat{\mathbf{z}}$ \cite{stratton1941}. Stress-integral and mechanical torque data are comparable only under a common reference origin.
\end{example}

The choice of integration surface belongs to the field problem itself \cite{stratton1941}. Two closed surfaces share the same force and torque resultants only if one can be deformed into the other without crossing sources, material, or singularities, if stress is evaluated on the same Maxwell-valid $(\E,\H)$ branch in the intervening source-free region, and if storage terms are handled consistently \cite{maxwell1873,stratton1941}. A partial patch $S_m$ cannot stand for a full closed $\partial V_b$ unless omitted segments and edge contributions are restored \cite{stratton1941}.

\phantomsection\label{prop:ch1.stress-torque-admissibility}
Stress-integral force and torque carry meaning in Chapter~\ref{ch:01} only when the integration surface, reference origin, and torque axis are declared alongside the field solution \cite{stratton1941,beth1936}.

Maxwell consistency delineates which field traces may transport flux on $\mathcal{F}_{\mathrm{rad}}$ \cite{stratton1941}.
\subsection{Field quantities, instrument readings, and inferred structure}
\label{sec:ch01.5.3}

Maxwell theory fixes $\E$ and $\H$ from sources, constitutive laws, boundary data, and the outgoing radiation condition \cite{maxwell1873,stratton1941}. At a sample point or on a finite patch, $\E$ and $\H$ are field quantities; closed-surface integrals of $\mathbf{S}_\mathrm{c}$, traction, and stress moment are observables tied to a declared surface, outward normal, and origin. A receiver couples to tangential fields through its geometry and constitutive response, and a dial reading is the mechanical or electronic transduction of that coupling---not an alternative definition of flux \cite{poynting1884,beth1936}.

Harmonic or modal coefficients of $\E_t$ arise only after one chooses a representation on an aperture or reference sphere. They compress angular structure for synthesis and diagnosis in later chapters, yet they do not replace Eqs.~\eqref{eq:ch1.power-balance-integral}, \eqref{eq:ch1.stress-force-integral}, or \eqref{eq:ch1.ex61-flux-from-stress} unless the partial data are supplemented by source specification, boundary conditions, and outgoing selection so that a unique Maxwell continuation exists \cite{hansen1988,stratton1941}.

\subsection{Apertures, angular charts, and later radiation mappings}
\label{sec:ch01.5.4}

On a large source-enclosing sphere $S_R$ in the source-free exterior, tangential $\E_t(\theta,\phi)$ and $\H_t(\theta,\phi)$ with outgoing radiation data fix the angle dependence of the field that sources radiate \cite{stratton1941}. A finite aperture samples that tangential field through an opening in a conducting screen; the aperture chart is a \emph{restriction} of the global solution, not an automatically complete radiating Maxwell state unless impressed currents behind the screen, material data on the rim, and continuity of tangential $\E$ and $\H$ across the aperture are stated together with outgoing selection in the half-space beyond \cite{hansen1988,harrington2001}.

Angular spectra and vector-harmonic expansions reorder the same tangential samples for propagation and pattern synthesis. They streamline the bookkeeping of direction and polarisation, but conservation of power, momentum, and angular momentum on a chosen shell still requires the stress and Poynting integrals of section~\ref{sec:ch01.1} evaluated on a Maxwell-valid $(\E,\H)$ pair \cite{stratton1941}. Chapter~\ref{ch:02} connects this angular organization to radiation operators and dyadic Green functions on the outgoing branch established in section~\ref{sec:ch01.1.3} \cite{harrington2001,tai1994}.

\phantomsection\label{prop:ch1.observable-inference-criterion}
Throughout Chapter~\ref{ch:01}, every reported value is either a Maxwell functional on a chosen surface---flux, traction, or torque---or a quantity inferred from partial field data through a stated representation; the distinction must be explicit \cite{stratton1941}.

\section{Examples and Counterexamples}
\label{sec:ch01.6}

The preceding sections stated general balances and their limitations. The examples below test those statements on concrete fields: radiation with circular polarization, modal superpositions that cancel net flux, gauge dependence of local partitions, and the radial shell structure of a Hertzian dipole including reactive power and the volume energy ratio $Q_{\mathrm{vol}}$ \cite{stratton1941,beth1936,allen1992,chu1948,mcclean1996}.

\subsection{Angular momentum conservation in radiation}
\label{sec:ch01.6.1}

Circular polarization fixes the handedness with which $\E$ and $\H$ rotate about the beam axis \cite{stratton1941,beth1936}. Mechanical torque on matter, and signed outward transport through a source-free shell, are read from the Maxwell stress integrals below---not from a polarization label or phase portrait alone. The axial flux definition is
\begin{equation}
\label{eq:ch1.ex61-flux-from-stress}
\dot J_{z,\mathrm{em,out}}
=-\oint_{S_R}\big[\mathbf{r}\times(\Tavg\cdot\hat{\mathbf{r}})\big]\cdot\hat{\mathbf{z}}\,dS,
\end{equation}
which is Eq.~\eqref{eq:ch1.outward-am-flux} with outward normal $\hat{\mathbf{n}}=\hat{\mathbf{r}}$ on a sphere $S_R$ in the source-free exterior \cite{stratton1941,beth1936}. For any Maxwell-valid pair $(\E,\H)$ on $S_R$,
\begin{equation}
\label{eq:ch1.ex61-power-def}
P_{\mathrm{out}}:=\oint_{S_R}\Re\{\mathbf{S}_\mathrm{c}\cdot\hat{\mathbf{r}}\}\,dS,\qquad
\mathbf{S}_\mathrm{c}=\tfrac12\E\times\H^*,
\end{equation}
with $\dot J_{z,\mathrm{em,out}}$ from Eq.~\eqref{eq:ch1.ex61-flux-from-stress}.

On a source-free sphere, outgoing spherical vector waves organize tangential data and foreshadow the radiation operators of Chapter~\ref{ch:02} \cite{stratton1941,hansen1988,jackson1999}. Fix $\exp(+\jj\omega t)$ and outgoing $h_\ell^{(2)}(kr)=j_\ell(kr)-\jj y_\ell(kr)$ \cite{stratton1941}. With $\sigma=\sqrt{\ell(\ell+1)}$ and $\vec L=\jj\mathbf r\times\nabla$,
\begin{equation}
\label{eq:ch1.ex61-vsh}
\vec Y_{\ell m}^{\ell m}=\frac{\vec L Y_\ell^m}{\sigma}
=\hat{\boldsymbol{\theta}}\,\frac{m Y_\ell^m}{\sigma\sin\theta}
+\hat{\boldsymbol{\phi}}\,\frac{\jj\,\partial_\theta Y_\ell^m}{\sigma},
\end{equation}
so that $\int_{4\pi}|\vec Y_{\ell m}^{\ell m}|^2\,d\Omega=1$ when $\int_{4\pi}|Y_\ell^m|^2\,d\Omega=1$ \cite{jackson1999}. The magnetic-type wave is defined by $\mathbf M_{\ell m}^+=h_\ell^{(2)}(kr)\vec Y_{\ell m}^{\ell m}$ and $\mathbf N_{\ell m}^+=(1/k)\nabla\times\mathbf M_{\ell m}^+$ \cite{stratton1941,jackson1999}. Coefficient templates that pair $\mathbf M_{\ell m}^+$ and $\mathbf N_{\ell m}^+$ with Poynting and stress fluxes still require a print-verified Stratton--Jackson--Hansen convention lock together with a physically consistent Maxwell pair $(\E,\H)$; that synthesis remains outside Chapter~\ref{ch:01}.

\begin{remark}[Hansen modal flux deferred to Chapter~\ref{ch:02}]
\label{rem:ch1.hansen-flux-deferred}
No single-mode Hansen power law or closed-form $\mathbf N_{\ell m}^+$ transport ratio is asserted as a Maxwell theorem in this chapter. Poynting and axial stress-moment integrals on $S_R$ remain defined by Eqs.~\eqref{eq:ch1.ex61-power-def} and~\eqref{eq:ch1.ex61-flux-from-stress} for any \emph{actual} source-free $(\E,\H)$ \cite{stratton1941,jackson1999,hansen1988}.
\end{remark}

Jackson's third-edition shell-content formulas for multipole energy and angular momentum in a spherical layer (Chapter~9, Eqs.~(9.136) and (9.143)--(9.144) in the citation record) are not quoted here without page-level verification. Because Eq.~\eqref{eq:ch1.ex61-flux-from-stress} is a Maxwell-stress \emph{surface} flux and Jackson's formulas are \emph{volume} shell-content identities under a different phasor branch, no Jackson shell equation is displayed in Chapter~\ref{ch:01}. A verified conversion between those objects is deferred to Chapter~\ref{ch:02} \cite{jackson1999}.

\begin{remark}[Vector generator and mixed multipoles]
\label{rem:ch1.am-power-target}
No Maxwell-derived transport law of the modal form
$\dot J_{z,\mathrm{em,out}}=(M/\omega)\,P_{\mathrm{out}}$
is asserted for arbitrary $\mathcal{J}_z$ eigenfields (compare Eq.~\eqref{eq:ch1.ex61-flux-from-stress}). Distinct $(\ell,m)$ contributions and electric--magnetic cross terms require separate control \cite{jackson1999}. Literal $\phi$-only modal templates without a Maxwell-valid $\theta$ envelope are bookkeeping only, not spherical vector modes.
\end{remark}

\paragraph{Equatorial ring source (separate construction).}
For a filamentary ring of radius $a$ in the plane $\theta=\pi/2$, use line current $I(\phi)$ in amperes and line charge density $\lambda(\phi)$ in coulombs per metre. Under $\exp(+\jj\omega t)$, charge conservation along the ring gives
\[
\frac{1}{a}\frac{\partial I}{\partial\phi}+\jj\omega\lambda=0.
\]
Let $I_m$ denote the $m$th azimuthal Fourier coefficient (amperes) in $I(\phi)=\sum_m I_m e^{\jj m\phi}$: it is a harmonic of this \emph{one-dimensional} model, not a sheet density. A finite-width conducting strip instead requires surface current $\mathbf K$ (A\,m$^{-1}$) and surface charge density $\sigma$ (C\,m$^{-2}$) with $\nabla_s\cdot\mathbf K+\jj\omega\sigma=0$ on the support; the two models must not be mixed when assigning units to $I_m$. A single $m\neq0$ harmonic does not, by itself, define a complete radiating Maxwell solution: impressed sources, continuity on the stated support measure, boundary data, and outgoing selection remain necessary \cite{stratton1941,tai1994}. Projection onto regular spherical vector waves is deferred to later chapters.
\begin{remark}[Modal $M=m+\sigma$ bookkeeping only]
\label{rem:ch1.m-plus-sigma}
With $\hat{\mathbf{e}}_\sigma=(\hat{\mathbf{x}}+\jj\sigma\hat{\mathbf{y}})/\sqrt{2}$ and helicity label $\sigma=\pm1$, the symbol $(m+\sigma)P/\omega$ is \emph{modal bookkeeping only} when a later chapter supplies a Maxwell-valid vector mode; it is not a Chapter~\ref{ch:01} theorem \cite{allen1992,cerjan2011,monteiro2009}.
\end{remark}

\paragraph{Plane-wave stress countercheck (symmetric open disk).}
For a uniform outgoing circularly polarized plane wave $\E=(E_0/\sqrt{2})(\hat{\mathbf{x}}+\jj\sigma\hat{\mathbf{y}})e^{-\jj kz}$, $\H=\eta_0^{-1}\hat{\mathbf{z}}\times\E$, the Maxwell stress tensor is spatially constant. On a finite transverse disk centered on the beam axis, $\oint[\mathbf{r}\times(\Tavg\cdot\hat{\mathbf{z}})]_z\,dS=0$ by azimuthal symmetry: power through the disk is $P=|E_0|^2 A/(2\eta_0)$, but the symmetric stress-moment integral does \emph{not} yield $\sigma P/\omega$ \cite{stratton1941,beth1936,stewart2005}. Torque on an anisotropic load in Beth's sense \cite{beth1936} is not recovered from this uniform-beam disk formula \cite{beth1936,stratton1941}.

\paragraph{Spin flux and ideal half-wave-plate scales (plane wave only).}
For the same CP plane wave with time-averaged power $P$ and helicity label $\sigma=\pm1$ in $\E\propto\hat{\mathbf{x}}+\jj\sigma\hat{\mathbf{y}}$, a helicity bookkeeping magnitude
\begin{equation}
\label{eq:ch1.cp-spin-flux-magnitude}
|\dot J_{z,\mathrm{spin,inc}}|=\frac{P}{\omega}
\end{equation}
is \emph{not} identified with the symmetric Maxwell-stress integral $\dot J_{z,\mathrm{em,out}}$ of Eq.~\eqref{eq:ch1.ex61-flux-from-stress} on a centred transverse disk (which vanishes for that symmetry). Eq.~\eqref{eq:ch1.cp-spin-flux-magnitude} records a conventional spin/helicity scale for uniform CP in vacuum, distinct from the governing stress-surface observable \cite{stratton1941,beth1936,belinfante1940}. A perfect absorber that removes the beam with definite helicity exchanges mechanical torque of magnitude $P/\omega$ about $\hat{\mathbf{z}}$ \cite{beth1936}. An ideal lossless half-wave plate that reverses helicity while transmitting the same power changes the field helicity by $\Delta\sigma=-2\sigma$; angular-momentum conservation then requires load torque of magnitude
\begin{equation}
\label{eq:ch1.hwp-torque-magnitude}
|\tau_{z,\mathrm{load}}|=\frac{2P}{\omega}
\end{equation}
for the ideal reversal case. Eq.~\eqref{eq:ch1.hwp-torque-magnitude} is a polarization-conversion scale, not Beth's quartz-plate measurement geometry and not the linearly polarized Hertzian dipole of Subsection~\ref{sec:ch01.6.4} \cite{beth1936,stratton1941}.

\FloatBarrier
Under monochromatic steady state with vanishing cycle-averaged storage rate, torque on the source support and outward field angular-momentum transport are related by opposite signs:
\begin{equation}
\label{eq:ch1.ex61-torque-flux-preview}
\boldsymbol{\tau}_{z,\mathrm{source}}=-\dot J_{z,\mathrm{em,out}},
\qquad
\boldsymbol{\tau}_{z,\mathrm{load}}=+\dot J_{z,\mathrm{em,out}}.
\end{equation}
Eq.~\eqref{eq:ch1.ex61-torque-flux-preview} states the same conservation law in rotational form as Eqs.~\eqref{eq:ch1.angular-balance-integral} and~\eqref{eq:ch1.outward-am-flux}. Beth \cite{beth1936} rendered the law tangible as mechanical torque on a quartz wave plate in circularly polarized light: because the plate alters the transmitted polarization, load torque follows from the net angular-momentum flux through the plate's control surface—incident minus outgoing—not from the source's isolated outward flux alone \cite{beth1936,stratton1941}. The torque is not read from the polarization ellipse; it is fixed by stress-moment balance at the anisotropic interface \cite{beth1936,stratton1941}.

The primary radiation illustration is a localized $z$-directed electric dipole on finite support, enclosed by a far-zone sphere $S_R$. Figure~\ref{fig:ch1.ex61-beth-helical-sketch} is a torque-balance schematic only: it neither specifies a rotating dipole moment nor supplies a closed-sphere stress proof for a modal transport law. On $S_R$ in the source-free exterior,
$P_{\mathrm{rad}}=\oint_{S_R}\Re\{\mathbf{S}_\mathrm{c}\cdot\hat{\mathbf{n}}\}\,dS$ is
finite and independent of $R$ on any enclosing sphere in the lossless source-free exterior
\cite{stratton1941,beth1936}. On a predominantly radiative shell, the field content is read most transparently as outward radiation. The dipole example uses $P_{\mathrm{rad}}$ from stress/Poynting integration on $S_R$ and does not assert a universal $(m+\sigma)P/\omega$ law \cite{stratton1941}.

\begin{figure}[b]
\centering
\ChOneFigBlock{%
\begin{tikzpicture}[ch1geo]
 \draw[ch1shell, dashed, line width=0.45pt] (0,0) circle (0.95);
 \node[ch1oncanvas, font=\scriptsize] at (0,1.12) {$S_R$};
 \draw[ch1axis, line width=0.5pt] (0,-0.15) -- (0,1.35);
 \node[ch1oncanvas, font=\scriptsize] at (0.12,1.35) {$\hat{\mathbf{z}}$};
 \draw[-{Latex[length=1.6mm,width=1mm]}, line width=0.5pt]
   (0.85,0.25) arc (15:75:0.85);
 \node[ch1oncanvas, font=\scriptsize] at (1.05,0.55) {$\mathbf{S}_\mathrm{c}$};
 \draw[line width=0.5pt] (-0.95,0) -- (-0.55,0);
 \draw[-{Latex[length=1.6mm,width=1mm]}, line width=0.55pt]
   (-0.75,0) -- (-0.75,0.4);
 \node[ch1oncanvas, font=\scriptsize] at (-1.05,0.5) {$\tau_z$};
 \node[ch1oncanvas, font=\scriptsize, align=center] at (0,-1.15)
 {torque on matter vs.\ flux through $S_R$};
\end{tikzpicture}%
}
\caption{Axisymmetric Hertzian dipole radiation (linear $\hat{\theta}$ polarization in the far zone) enclosed by far-zone sphere $S_R$: outward Poynting flux balances finite radiated power; for this $m{=}0$ pattern the net axial stress-moment flux on $S_R$ vanishes while $P_{\mathrm{rad}}$ is finite. Circular polarization and Beth-type load torque are separate examples (Fig.~\ref{fig:ch1.beth-schematic}) \cite{stratton1941,beth1936}.}
\label{fig:ch1.ex61-beth-helical-sketch}
\end{figure}

Figure~\ref{fig:ch1.ex61-beth-helical-sketch} depicts the localized $\hat{\mathbf{z}}$-directed \emph{electric} Hertzian dipole: in the far zone $\E_\theta\propto\sin\theta$ is linearly polarized in $\hat{\boldsymbol{\theta}}$, not circularly polarized \cite{stratton1941,harrington2001}. The sketch places torque--flux balance on $S_R$ at stand-off \(R=5\lambda\), where \(|P_{\mathrm{reac}}(R)|/P_{\mathrm{rad}}\ll 1\) in the sense of Subsection~\ref{sec:ch01.2.4}: reactive exchange no longer dominates the normal flux pattern, and the stress integral registers radiation rather than near-zone oscillation \cite{stratton1941}.

On a predominantly radiative shell, fields scale as \(|\E|\propto 1/R\) with the azimuthal phase of the dominant mode; Maxwell stress Eq.~\eqref{eq:ch1.maxwell-stress} with outward normal \(\hat{\mathbf{n}}=\hat{\mathbf{r}}\) supplies the integrand \([\mathbf{r}\times(\Tavg\cdot\hat{\mathbf{n}})]\cdot\hat{\mathbf{z}}\). Multimode superpositions require summing stress integrals on the actual $(\E,\H)$ field \cite{stratton1941}. The same balance evaluated at \(R=\lambda\) without stating \(|P_{\mathrm{reac}}(r)|/P_{\mathrm{rad}}\) can mis-attribute reactive-dominated storage exchange as far-zone transport \cite{stratton1941}.

Under strict monochromatic steady state at angular frequency $\omega$, the cycle-averaged stored field angular momentum $\langle\mathbf{J}_{\mathrm{em}}\rangle$ is $T$-periodic, so $\partial_t\langle\mathbf{j}_{\mathrm{em}}\rangle=\mathbf{0}$ in Eq.~\eqref{eq:ch1.angular-momentum-balance-averaged} and the storage contribution in Eq.~\eqref{eq:ch1.torque-stress-surface} vanishes on a radiative shell. The inherited balance then reduces to
\begin{equation}
\label{eq:ch1.ex61-torque-flux}
\dot J_{z,\mathrm{em,out}}
=
-\oint_{S_R}
\big[\mathbf{r}\times(\Tavg\cdot\hat{\mathbf{n}})\big]\cdot\hat{\mathbf{z}}\,dS
=-\mathcal{T}_{J,z},
\end{equation}
where \(\mathcal{T}_{J,z}\) is the $z$-component of the Maxwell-stress moment on \(S_R\). Outward transport \(\dot J_{z,\mathrm{em,out}}=\Phi_{J,z,\mathrm{out}}\) is the negative of the stress-moment integral under the convention of Eq.~\eqref{eq:ch1.outward-am-flux}.

For an absorbing load that intercepts the declared net angular-momentum flux through its enclosing control surface, \(\tau_{z,\mathrm{load}}=\dot J_{z,\mathrm{em,out}}\); for the enclosed source support, \(\tau_{z,\mathrm{source}}=-\dot J_{z,\mathrm{em,out}}=\mathcal{T}_{J,z}\) in steady operation \cite{beth1936}. A partially transmitting or polarization-changing load requires the full incoming and outgoing balance at that body’s surface; the source’s outward flux is not automatically equal to the load torque \cite{beth1936,stratton1941}. Modal-class symbols such as $(m+\sigma)P/\omega$ are not substituted for Eq.~\eqref{eq:ch1.ex61-torque-flux} without a Maxwell-valid mode proof in a later chapter \cite{stratton1941}.

When several \((m,\sigma)\) modes carry comparable power, no single-mode bookkeeping symbol replaces the stress integral in Eq.~\eqref{eq:ch1.ex61-torque-flux} \cite{stratton1941}.

When exchange is mediated at a material boundary, electromagnetic torque and angular-momentum flux are two faces of one conservation law \cite{beth1936}. Beth's experiment showed that circularly polarized radiation carries angular momentum whose mechanical counterpart is torque on matter exerted through Maxwell stress at the interface \cite{beth1936}.

\begin{figure}[t]
\centering
\VolIPhysFig{\VolIPhysBethSchematic}
\caption{Beth \cite{beth1936}: circularly polarized light incident on a wave plate on a torsion suspension; the plate receives torque through Maxwell stress at the interface.}
\label{fig:ch1.beth-schematic}
\end{figure}

\begin{example}[Electromagnetic torque scale]
\label{ex:ch1.beth-torque}
If one \emph{assumes} a modal bookkeeping relation $\dot J_z=P/\omega$ for circular polarization with $P=1\,\mathrm{W}$, then at $f=10\,\mathrm{GHz}$, $\omega=2\pi f$ and $P/\omega\approx1.59\times10^{-11}\,\mathrm{N\,m}$. This numeric illustration is \emph{not} a Maxwell-stress integral over a uniform plane wave on a symmetric disk (Subsection~\ref{sec:ch01.6.1}); Beth-type loads require the anisotropic boundary model of the original experiment \cite{beth1936,stratton1941}.
\end{example}

\subsection{Phase structure without angular-momentum transport}
\label{sec:ch01.6.2}

One may construct fields with rich angular phase structure yet vanishing net angular-momentum transport. Take two coherent components with azimuthal orders \(+m\) and \(-m\) on an open transverse plane \(S_o\). The forward power and forward angular-momentum flux associated with the outgoing modal components crossing \(S_o\) are forward-propagating modal statements on that aperture; the exact global conservation law remains the closed-surface Maxwell-stress balance on an enclosing surface \(\Gamma_R\) that bounds a control volume containing the source \cite{stratton1941,belinfante1940}. Their signed modal fluxes oppose; equal powers cancel signed orbital transport even while intensity and phase retain strong structure \cite{allen1992}. Signed orbital transport is proportional to the intensity difference \(I_+-I_-\), not to the visibility of azimuthal structure alone \cite{stratton1941}.

\begin{figure}[t]
\centering
\VolIPhysFig{\VolIPhysOpenPlaneVsClosedSphere}
\caption{Open transverse plane $S_o$ (modal flux and power through the aperture) versus closed enclosing sphere $\Gamma_R$ (exact Maxwell stress-moment conservation on a control surface) \cite{stratton1941,beth1936}.}
\label{fig:ch1.open-plane-vs-closed}
\end{figure}

With a common radial envelope \(A(\rho)\), the transverse field on \(S_o\) varies strongly in both amplitude and azimuthal phase. Allen \emph{et al.} \cite{allen1992} showed that a \emph{single} helical mode radiates definite orbital flux; equal counter-rotating weights remove net signed angular-momentum transport while retaining visible structure and substantial radiated power \cite{poynting1884}.

\begin{figure}[t]
\centering
\VolIPhysFig{\VolIPhysCounterRotatingAperture}
\caption{Counter-rotating components with $m=\pm2$ on open aperture $S_o$ (schematic): with matched amplitudes, $|E_t|^2\propto\cos^2(2\phi)$ shows four intensity lobes on the plane, while equal modal powers $P_+=P_-$ cancel signed forward orbital angular-momentum flux through $S_o$ ($\dot J_z=0$ on the open-surface modal bookkeeping, distinct from a closed-sphere stress integral) \cite{allen1992,poynting1884}.}
\label{fig:ch1.ex62-counter-rotating-aperture}
\end{figure}

Visible azimuthal structure alone does not fix net angular-momentum transport: when \(I_+=I_-\), signed flux through $S_o$ vanishes while $P_{\mathrm{tot}}=P_++P_-$ may remain large (Figure~\ref{fig:ch1.phase-vs-transport-sketch}) \cite{beth1936,stratton1941}. Component intensities follow from time-averaged Poynting flux:
\begin{equation}
\label{eq:ch1.ex62-intensity-definition}
I_\pm(\rho)=\tfrac12\Re\!\left\{\mathbf{E}_\pm\times\mathbf{H}_\pm^*\right\}\cdot\hat{\mathbf{z}},
\end{equation}
with $\mathbf{H}_\pm=(1/\eta_0)\,\hat{\mathbf{z}}\times\mathbf{E}_\pm$ on an \emph{outgoing paraxial/modal} branch (exact for a uniform plane wave; for spatially varying $\mathbf{E}_\pm$ it is a template, not a full-vector Maxwell solution unless longitudinal components and $\curl$ residuals are verified) and $\mathbf{S}_{\mathrm{c},\pm}=\tfrac12\mathbf{E}_\pm\times\mathbf{H}_\pm^*$ \cite{stratton1941,allen1992}.

\paragraph{Peak-amplitude convention (counter-rotating pair).}
Fix peak tangential phasor amplitudes $A_\pm(\rho)$ and linear polarization $\hat{\mathbf{e}}_t$. The $\pm m$ orbital components are
\begin{equation}
\label{eq:ch1.ex62-counter-rotating-components}
\mathbf{E}_\pm(\rho,\phi)=\frac{A_\pm(\rho)}{\sqrt{2}}\,e^{\pm\jj m\phi}\,\hat{\mathbf{e}}_t,
\qquad
\mathbf{H}_\pm=\frac{1}{\eta_0}\,\hat{\mathbf{z}}\times\mathbf{E}_\pm.
\end{equation}
Algebraically,
\begin{equation}
\label{eq:ch1.ex62-counter-rotating-sum-derive}
\mathbf{E}_++\mathbf{E}_-
=\frac{A_+ e^{\jj m\phi}+A_- e^{-\jj m\phi}}{\sqrt{2}}\,\hat{\mathbf{e}}_t;
\end{equation}
when $A_+=A_-=A$ this is $\sqrt{2}\,A(\rho)\cos(m\phi)\,\hat{\mathbf{e}}_t$, not $2A\cos(m\phi)\,\hat{\mathbf{e}}_t$.
With Eq.~\eqref{eq:ch1.ex62-intensity-definition}, each component carries
$I_\pm(\rho)=|A_\pm(\rho)|^2/(4\eta_0)$ on a uniform outgoing branch; matched peaks ($|A_+|=|A_-|$) give equal component powers $P_\pm=\int_{S_o}I_\pm\,dS$ while the coherent intensity $|E_t|^2$ includes the cross term $2\Re(\mathbf{E}_+\cdot\mathbf{E}_-^*)$ and is not $P_+/P_-$ alone.

The total transverse field is
\begin{equation}
\label{eq:ch1.ex62-counter-rotating-superposition}
\mathbf{E}_t(\rho,\phi)
=
\mathbf{E}_+(\rho,\phi)+\mathbf{E}_-(\rho,\phi),
\end{equation}
which equals $\sqrt{2}\,A(\rho)\cos(m\phi)\,\hat{\mathbf{e}}_t$ when $A_+=A_-=A$.
With matched peaks, $I_+=I_-=|A|^2/(4\eta_0)$. On a uniform outgoing branch with identical envelopes, component powers satisfy
\begin{equation}
\label{eq:ch1.ex62-power-amplitude-ratio}
\frac{P_+}{P_-}=\left(\frac{|A_+|}{|A_-|}\right)^2,
\qquad
P_\pm=\int_{S_o} I_\pm\,dS.
\end{equation}
For each orbital component, the time-averaged angular-momentum flux through \(S_o\)
scales as \((\pm m/\omega)\) times component power
\(P_\pm = \int_{S_o} I_\pm\,dS\) \cite{stratton1941}. Summing components yields the open-aperture modal bookkeeping (distinct from a closed-sphere Maxwell-stress integral for an arbitrary full-vector field)
\begin{equation}
\label{eq:ch1.ex62-channel-sum}
\dot J_{z,\mathrm{modal},S_o}
\;\stackrel{\mathrm{modal\ bookkeeping}}{=}\;
\frac{m}{\omega}(P_+ - P_-)
=
\frac{m}{\omega}\int_{S_o}(I_+-I_-)\,dS,
\end{equation}
which is Eq.~\eqref{eq:ch1.ex62-net-am-transport}. When \(I_+=I_-\), both \(P_+\) and
\(P_-\) are equal and the bookkeeping symbol vanishes even though the coherent intensity
\(|E_t|^2=|E_+|^2+|E_-|^2+2\Re(\mathbf{E}_+\cdot\mathbf{E}_-^*)\) exhibits a standing azimuthal pattern through the cross term; modal power budgets use \(P_\pm=\int I_\pm\,dS\), not \(|E_t|^2\) alone \cite{stratton1941,allen1992}.

\begin{example}[Equal $+m$ and $-m$ components]
\label{ex:ch1.counter-rotating-numeric}
Take $m=2$ with equal peak amplitudes $|A_+|=|A_-|$ so $P_+=P_-$. Then the open-aperture bookkeeping of Eq.~\eqref{eq:ch1.ex62-net-am-transport} vanishes. For a one-percent \emph{power} imbalance, $P_+=1.01P_-$, Eq.~\eqref{eq:ch1.ex62-power-amplitude-ratio} requires $|A_+|/|A_-|=\sqrt{1.01}\approx 1.0050$ at fixed $\eta_0$---not $|A_+|/|A_-|=1.01$, which would imply $P_+/P_-=1.0201$ (a $2.01\%$ power imbalance). The bookkeeping transport is $\dot J_z=(m/\omega)(P_+-P_-)=(m/\omega)(0.01P_-)$; e.g.\ if $P_-=1\,\mathrm{W}$ and $\omega=2\pi f$, then $\dot J_z=(2/\omega)(0.01\,\mathrm{W})$. This is not a closed-sphere Maxwell-stress evaluation \cite{allen1992}.
\end{example}

Peak amplitudes $A_\pm(\rho)$ in Eq.~\eqref{eq:ch1.ex62-counter-rotating-components} fix $I_\pm(\rho)=|A_\pm(\rho)|^2/(4\eta_0)$ on a uniform outgoing branch (Eq.~\eqref{eq:ch1.ex62-intensity-definition}). Let
\(\dot J_z\) be net angular-momentum transport rate through \(S_o\). The inherited
balance law gives
\begin{equation}
\label{eq:ch1.ex62-net-am-transport}
\dot J_z
=
\frac{m}{\omega}
\int_{S_o}\!\big(I_+-I_-\big)\,dS.
\end{equation}
Here \(m/\omega\) is the orbital angular-momentum-per-energy factor and \(I_+-I_-\) is
signed transport imbalance. Signed net transport follows \(I_+-I_-\), not \(|E_t|^2\) or its azimuthal derivative \cite{stratton1941}. The equal-amplitude counterexample is
\begin{equation}
\label{eq:ch1.ex62-equal-channel-condition}
I_+(\rho)=I_-(\rho)\quad\forall \rho\in S_o.
\end{equation}
Substitution into Eq.~\eqref{eq:ch1.ex62-net-am-transport} yields
\begin{equation}
\label{eq:ch1.ex62-zero-transport-result}
\dot J_z=0
\qquad\text{while}\qquad
E_t(\rho,\phi)=\sqrt{2}\,A(\rho)\cos(m\phi)\neq 0.
\end{equation}
If \(P_+\neq P_-\), a finite signed flux remains. In steady operation, section~\ref{sec:ch01.5} gives
\begin{equation}
\label{eq:ch1.ex62-zero-torque-prediction}
\tau_{z,\mathrm{source}}
=
-\dot J_z
=
0
\quad\text{under Eq.~\eqref{eq:ch1.ex62-equal-channel-condition}}.
\end{equation}
Let
\(I_+(\rho)=I_0(\rho)+\delta I(\rho)\) and
\(I_-(\rho)=I_0(\rho)-\delta I(\rho)\). Then
\begin{equation}
\label{eq:ch1.ex62-imbalance-sensitivity}
\dot J_z
=
\frac{2m}{\omega}\int_{S_o}\delta I(\rho)\,dS.
\end{equation}
\phantomsection\label{prop:ch1.ex62-phase-insufficiency}
Under Eq.~\eqref{eq:ch1.ex62-equal-channel-condition}, Eqs.~\eqref{eq:ch1.ex62-zero-transport-result} and~\eqref{eq:ch1.ex62-zero-torque-prediction} give \(\dot J_z=0\) and \(\tau_{z,\mathrm{source}}=0\) even when \(|E_t|\) and \(\partial_\phi E_t\) are large on \(S_o\) \cite{stratton1941,beth1936}. A finite signed flux requires \(I_+\neq I_-\), as in Eq.~\eqref{eq:ch1.ex62-imbalance-sensitivity}. Unequal modal powers, polarization-selective loss, incomplete modal orthogonality, or aperture truncation can break exact cancellation; gauge relabeling cannot restore transport removed by balanced superposition \cite{poynting1884,harrington2001}.
\subsection{Gauge dependence of local angular-momentum densities}
\label{sec:ch01.6.3}

Under \(\mathbf{A}'=\mathbf{A}+\nabla\chi\) and \(\phi'=\phi-\jj\omega\chi\), the fields \(\E=-\nabla\phi-\jj\omega\mathbf{A}\) and \(\mathbf{B}=\nabla\times\mathbf{A}\) are unchanged \cite{maxwell1873}; two potential descriptions may therefore label the same electromagnetic field. Any local density constructed from \(\phi\) or \(\mathbf{A}\) alone inherits gauge sensitivity; stress traction, Poynting flux, and total angular-momentum density built from \((\E,\H)\) do not \cite{belinfante1940,stratton1941}.

An ideal oscillating electric point dipole inside a gauge sphere $V_g$ of radius $R=2\lambda$ makes the distinction concrete. Coulomb and Lorenz gauges yield the same $(\E,\H)$ on $\partial V_g$ yet different $\nabla\phi$ at interior points, so any potential-linked local map differs pointwise between them. Mechanical torque $\oint_{\partial V_g}\mathbf{r}\times(\Tavg\cdot\hat{\mathbf{n}})\,dS$ and outward angular-momentum flux built from $(\E,\H)$ remain unchanged \cite{belinfante1940,stratton1941}.

\begin{figure}[t]
\centering
\ChOneFigBlock{%
\begin{tikzpicture}[ch1geo]
 \draw[ch1shell, dashed, line width=0.45pt] (0,0) circle (0.9);
 \node[ch1oncanvas, font=\scriptsize] at (0,1.05) {$\partial V_g$};
 \draw[gray!60!black, line width=0.55pt] (0,0.15) -- (0,0.75);
 \draw[gray!60!black, line width=0.55pt] (-0.08,0.45) -- (0.08,0.45);
 \node[ch1oncanvas, font=\scriptsize] at (0.25,0.55) {dipole};
 \draw[-{Latex[length=1.6mm,width=1mm]}, gray!40!black, line width=0.5pt]
   (0.65,0.35) arc (30:120:0.65);
 \node[ch1oncanvas, font=\scriptsize] at (0.85,0.55) {$\mathcal{L}_C$};
 \draw[-{Latex[length=1.6mm,width=1mm]}, gray!45!black, line width=0.5pt, dashed]
   (0.55,-0.25) arc (-30:-120:0.55);
 \node[ch1oncanvas, font=\scriptsize] at (0.75,-0.45) {$\mathcal{L}_L$};
 \draw[-{Latex[length=1.4mm,width=0.9mm]}, gray!50!black, line width=0.5pt]
   (0.9,0) -- (1.15,0);
 \node[ch1oncanvas, font=\scriptsize] at (1.25,0.12) {$\mathbf{M}\!\cdot\!\hat{\mathbf{n}}$};
 \node[ch1oncanvas, font=\scriptsize, align=center] at (0,-1.15)
 {same $\E$, $\H$; different local maps};
\end{tikzpicture}%
}
\caption{Ideal oscillating electric point dipole inside gauge sphere $V_g$: Coulomb and Lorenz gauges yield different gauge-dependent local maps at interior points while identical $(\E,\H)$ on $\partial V_g$ fixes the same stress-flux integral \cite{belinfante1940,stratton1941}.}
\label{fig:ch1.ex63-gauge-volume-sketch}
\end{figure}

Gauge relabeling moves potential-based local maps while leaving Maxwell fields and gauge-invariant surface integrals built from \((\E,\H)\) unchanged \cite{maxwell1873,belinfante1940}. Subsection~\ref{sec:ch01.6.4} holds the same dipole fixed while moving the observation shell from reactive to radiative radius \cite{stratton1941,chu1948}.
\subsection{Radial structure of a localized radiator}
\label{sec:ch01.6.4}

As observation moves outward from a localized source, what changes in the electromagnetic field? The source remains fixed; only the radius of observation grows. The outgoing Maxwell solution repartitions radial content: near the source, $1/r^3$ and $1/r^2$ terms dominate, and the $1/r$ radiative branch eventually sets the transverse far-zone form \cite{stratton1941,poynting1884}. Consider a $z$-directed electric dipole moment phasor $\mathbf{p}=p_0\hat{\mathbf{z}}$ at the origin with $\exp(+\jj\omega t)$ and outgoing factor $e^{-\jj kr}$. In spherical coordinates $(r,\theta,\phi)$ the exact source-free fields are \cite{stratton1941,harrington2001}
\begin{subequations}
\label{eq:ch1.dipole-fields-explicit}
\begin{align}
E_r&=\frac{p_0\cos\theta}{4\pi\varepsilon_0}\,e^{-\jj kr}
\left(\frac{2}{r^3}+\frac{2\jj k}{r^2}\right),
\label{eq:ch1.dipole-Er}\\[0.25em]
E_\theta&=\frac{p_0\sin\theta}{4\pi\varepsilon_0}\,e^{-\jj kr}
\left(\frac{1}{r^3}+\frac{\jj k}{r^2}-\frac{k^2}{r}\right),
\label{eq:ch1.dipole-Etheta}\\[0.25em]
H_\phi&=\frac{\jj k p_0\sin\theta}{4\pi\varepsilon_0\eta_0}\,e^{-\jj kr}
\left(\frac{1}{r^2}+\frac{\jj k}{r}\right),
\label{eq:ch1.dipole-Hphi}
\end{align}
\end{subequations}
with $k=\omega/c$. For a $z$-directed dipole moment $\mathbf{p}=p_0\hat{\mathbf{z}}$, the retarded Lorenz-gauge vector potential in the source-free exterior ($r>0$) is
\begin{equation}
\label{eq:ch1.dipole-A-retarded}
\mathbf{A}(\mathbf{r})=\jj\omega\mu_0\,\frac{p_0\hat{\mathbf{z}}\,e^{-\jj kr}}{4\pi r},
\qquad r>0,
\end{equation}
consistent with $\mathbf{A}(\mathbf{r})=\mu_0\int \Jimp(\mathbf{r}')\,G(\mathbf{r}-\mathbf{r}')\,d^{3}r'$ and the chapter Green function \cite{stratton1941,harrington2001,orfanidis2014}. Because $\E=-\jj\omega\mathbf{A}-\nabla\phi$ and $\H=(1/\mu_0)\curl\mathbf{A}$, each spatial derivative of $e^{-\jj kr}/r$ lowers the power of $r$ and, when acting on the oscillatory factor, introduces powers of $\jj k$. Collecting terms produces the $r^{-3}$, $r^{-2}$, and $r^{-1}$ hierarchy displayed above: the $r^{-3}$ part matches the static dipole field; the $r^{-2}$ terms are the induction fields; the $-k^2/r$ term in $E_\theta$ together with $H_\phi\propto k/r$ is the radiative branch \cite{stratton1941}. At fixed $kr$, several orders can be comparable; ``near,'' ``transition,'' and ``far'' therefore mark dominant contributions rather than disjoint shells \cite{stratton1941,poynting1884}.

\begin{figure}[t]
\centering
\VolIPhysFig{\VolIPhysDipoleThreeZones}
\caption{Concentric observation shells: near zone ($kr\ll1$, $r^{-3}$ and $r^{-2}$), transition ($kr\sim1$), and far zone ($kr\gg1$, $r^{-1}$ radiation) \cite{stratton1941,poynting1884}.}
\label{fig:ch1.dipole-three-zones}
\end{figure}

\begin{example}[Reading the three regions of a dipole field]
\label{ex:ch1.dipole-kr}
At $f=300\,\mathrm{MHz}$, $k=2\pi/\lambda\approx6.3\,\mathrm{rad/m}$. At $r=0.5\,\mathrm{m}$, $kr\approx3.1$: multipole terms of comparable weight (not the $kr\sim 1$ scale cited elsewhere as a transition \emph{indicator}---that scale is qualitative; here $|P_{\mathrm{reac}}|/P_{\mathrm{rad}}$ follows Eq.~\eqref{eq:ch1.ex64-power-functionals}). At $r=5\,\mathrm{m}$, $kr\approx31\gg1$ and the $r^{-1}$ terms in Eqs.~\eqref{eq:ch1.dipole-Etheta} and~\eqref{eq:ch1.dipole-Hphi} dominate, giving transverse radiation with $U(\theta)\propto\sin^2\theta$, $P_{\mathrm{rad}}$ from Eq.~\eqref{eq:ch1.ex64-power-invariance}, and $D(\theta)=(3/2)\sin^2\theta$ with $D_{\max}=(3/2)$ from Eqs.~\eqref{eq:ch1.radiation-intensity-Ur}--\eqref{eq:ch1.directivity-definition} and Eq.~\eqref{eq:ch1.ex64-power-invariance} \cite{stratton1941,harrington2001}.
\end{example}

For $r>0$, $\nabla\cdot\E=0$ follows from the spherical-coordinate divergence condition
\[
\partial_r E_r+\frac{2}{r}E_r+\frac{1}{r}\partial_\theta E_\theta+\frac{\cot\theta}{r}E_\theta=0
\]
when Eqs.~\eqref{eq:ch1.dipole-Er} and~\eqref{eq:ch1.dipole-Etheta} are substituted \cite{stratton1941}. For $r>0$, the displayed triad satisfies $\nabla\cdot\E=0$ and $\nabla\cdot\H=0$ with $\H=H_\phi\hat{\boldsymbol\phi}$ independent of $\phi$, so only $\partial_\phi H_\phi$ enters the spherical divergence and vanishes \cite{stratton1941,harrington2001}. The curl identities equivalent to $\curl\E=-\jj\omega\mu_0\H$ and $\curl\H=+\jj\omega\varepsilon_0\E$ hold with $\Jimp=\mathbf{0}$; the radial curl component $(\curl\H)_r=(r\sin\theta)^{-1}\partial_\theta(\sin\theta H_\phi)$ balances $\jj\omega\varepsilon_0 E_r$ even though $H_r=H_\theta=0$ \cite{stratton1941,harrington2001}. At $r=0$ the ideal point dipole requires distributional sources. For a $z$-directed electric dipole moment phasor $\mathbf{p}=p_0\hat{\mathbf{z}}$ ($p_0$ in C\,m), the charge density in the distributional sense is
\begin{equation}
\label{eq:ch1.dipole-rho-distributional}
\rho(\mathbf{r})=-\mathbf{p}\cdot\nabla\delta^{3}(\mathbf{r})
=-p_0\,\partial_z\delta^{3}(\mathbf{r}),
\end{equation}
so that $\int f(\mathbf{r})\,\rho(\mathbf{r})\,d^{3}r=+p_0\,\partial_z f|_{\mathbf{r}=\mathbf{0}}$ for smooth compactly supported test functions $f$ (integration by parts on $\partial_z\delta^{3}$) \cite{jackson1999,stratton1941}. Gauss's law $\divg\E=\rho/\varepsilon_0$ then requires $\divg\E=-(p_0/\varepsilon_0)\,\partial_z\delta^{3}$ at the origin. The compatible impressed current is
\begin{equation}
\label{eq:ch1.dipole-Jimp-distributional}
\Jimp(\mathbf{r})=\jj\omega p_0\,\hat{\mathbf{z}}\,\delta^{3}(\mathbf{r}),
\end{equation}
because $\divg\Jimp+\jj\omega\rho=0$ with $\divg\big(\hat{\mathbf{z}}\,\delta^{3}\big)=\partial_z\delta^{3}$ in $\mathcal{D}'(\mathbb{R}^{3})$ \cite{jackson1999,stratton1941}. With outgoing $G(\mathbf{r},\mathbf{r}')=e^{-\jj k|\mathbf{r}-\mathbf{r}'|}/(4\pi|\mathbf{r}-\mathbf{r}'|)$ and $(\nabla^2+k^2)G=-\delta^{3}(\mathbf{r}-\mathbf{r}')$ in $\mathcal{D}'(\mathbb{R}^{3})$, the Lorenz potentials of Eqs.~\eqref{eq:ch1.vector-potential-retarded} and~\eqref{eq:ch1.lorenz-scalar-potential-retarded} give
\begin{equation}
\label{eq:ch1.dipole-A-from-source}
\mathbf{A}(\mathbf{r})=\jj\omega\mu_0\,\mathbf{p}\,G(\mathbf{r},\mathbf{0}),
\qquad
\phi(\mathbf{r})=\frac{1}{\varepsilon_0}\int \rho(\mathbf{r}')\,G(\mathbf{r},\mathbf{r}')\,d^3r'
=-\frac{1}{\varepsilon_0}\,\mathbf{p}\cdot\nabla G(\mathbf{r},\mathbf{0}),
\end{equation}
using $\int (-\mathbf{p}\cdot\nabla'\delta^{3}(\mathbf{r}'))\,G(\mathbf{r},\mathbf{r}')\,d^3r'=\mathbf{p}\cdot\nabla' G(\mathbf{r},\mathbf{r}')|_{\mathbf{r}'=\mathbf{0}}=-\mathbf{p}\cdot\nabla G(\mathbf{r},\mathbf{0})$ \cite{jackson1999,stratton1941}. With $\E=-\jj\omega\mathbf{A}-\nabla\phi$ and $\H=\mu_0^{-1}\curl\mathbf{A}$ (Eq.~\eqref{eq:ch1.lorenz-field-reconstruction}),
\begin{equation}
\label{eq:ch1.dipole-E-dyadic-distributional}
\E=\epsilon_0^{-1}\left[k^{2}\mathbf{p}\,G+(\nabla\nabla G)\cdot\mathbf{p}\right],
\qquad
\H=\jj\omega\,\nabla G\times\mathbf{p},
\end{equation}
which reduce to Eqs.~\eqref{eq:ch1.dipole-Er}--\eqref{eq:ch1.dipole-Hphi} for $r>0$ \cite{stratton1941,harrington2001,yaghjian1982}. For $r>0$, successive differentiation gives $\partial_i(4\pi r)^{-1}=-\hat r_i/(4\pi r^2)$ and
\begin{equation}
\label{eq:ch1.hessian-classical-r-positive}
\partial_i\partial_j\frac{1}{4\pi r}
=\frac{3\hat r_i\hat r_j-\delta_{ij}}{4\pi r^{3}},
\qquad r>0,
\end{equation}
because $\partial_j\hat r_i=(\delta_{ij}-\hat r_i\hat r_j)/r$. In $\mathcal{D}'(\mathbb{R}^{3})$ the same second derivative carries the spherical principal-value kernel above plus the isotropic contact,
\begin{equation}
\label{eq:ch1.hessian-distributional}
\partial_i\partial_j\frac{1}{4\pi r}
=\operatorname{PV}\!\left[\frac{3\hat r_i\hat r_j-\delta_{ij}}{4\pi r^{3}}\right]
-\tfrac13\delta_{ij}\,\delta^{3}(\mathbf{r}),
\end{equation}
which is the standard spherical source-region structure cited in Yaghjian \cite{yaghjian1980dyadic,yaghjian1982}. Eq.~\eqref{eq:ch1.hessian-classical-r-positive} is the pointwise restriction of Eq.~\eqref{eq:ch1.hessian-distributional} away from the origin. Gauss's electric law uses $\mathbf{p}\cdot\nabla[(\nabla^2+k^2)G]=-\mathbf{p}\cdot\nabla\delta^{3}$ applied to the Lorenz structure; it must not be read as if the $k^{2}G\mathbf{p}$ term alone supplied an isotropic $\delta^{3}$ contact---that contact enters the Hessian term in Amp\`ere--Maxwell \cite{yaghjian1980dyadic,yaghjian1982}. For smooth compactly supported $\psi$, distributional Gauss's law follows from $(\nabla^2+k^2)G=-\delta^{3}$:
\[
\divg\E=\epsilon_0^{-1}\mathbf{p}\!\cdot\!\nabla\big[(\nabla^2+k^2)G\big]
=-\epsilon_0^{-1}\mathbf{p}\!\cdot\!\nabla\delta^{3},
\]
so
\[
\langle\divg\E,\psi\rangle
=-\epsilon_0^{-1}\sum_i p_i\langle\partial_i\delta^{3},\psi\rangle
=+\epsilon_0^{-1}\mathbf{p}\!\cdot\!\nabla\psi(\mathbf{0})
=\langle\rho/\epsilon_0,\psi\rangle,
\]
using $\langle\partial_i\delta^{3},\psi\rangle=-\partial_i\psi(\mathbf{0})$ and Eq.~\eqref{eq:ch1.dipole-rho-distributional} \cite{jackson1999,stratton1941}. Charge continuity $\divg\Jimp+\jj\omega\rho=0$ in $\mathcal{D}'(\mathbb{R}^{3})$ is immediate from the displayed pair. Lorenz gauge $\divg\mathbf{A}+\jj\omega\mu_0\varepsilon_0\phi=0$ follows from Eq.~\eqref{eq:ch1.dipole-A-from-source} because $\divg\mathbf{A}=\jj\omega\mu_0\mathbf{p}\cdot\nabla G$ and, with $\phi=-\varepsilon_0^{-1}\mathbf{p}\cdot\nabla G$, $\jj\omega\mu_0\varepsilon_0\phi=-\jj\omega\mu_0\mathbf{p}\cdot\nabla G$ in $\mathcal{D}'(\mathbb{R}^{3})$ \cite{jackson1999,stratton1941}:
\begin{equation}
\label{eq:ch1.dipole-lorenz-gauge}
\nabla\cdot\mathbf{A}+\jj\omega\mu_0\varepsilon_0\phi
=\jj\omega\mu_0\mathbf{p}\cdot\nabla G
-\jj\omega\mu_0\mathbf{p}\cdot\nabla G=0.
\end{equation} Faraday's law is closed by the curl of Eq.~\eqref{eq:ch1.dipole-E-dyadic-distributional}. Because $(\nabla\nabla G)\cdot\mathbf{p}=\nabla(\mathbf{p}\cdot\nabla G)$ for constant $\mathbf{p}$,
\[
\curl\bigl[(\nabla\nabla G)\cdot\mathbf{p}\bigr]=\mathbf{0}
\quad\text{in }\mathcal{D}'(\mathbb{R}^{3}),
\]
so
\begin{equation}
\label{eq:ch1.dipole-faraday-distributional}
\curl\E
=\frac{k^{2}}{\epsilon_0}\,\nabla G\times\mathbf{p}
=-\jj\omega\mu_0\,\mathbf{H},
\qquad
\mathbf{H}=\jj\omega\,\nabla G\times\mathbf{p},
\end{equation}
with $\mathbf{H}=\mu_0^{-1}\curl\mathbf{A}$ and $\mathbf{A}=\jj\omega\mu_0\mathbf{p}\,G$ \cite{jackson1999,stratton1941}. Gauss's magnetic law holds because $\H$ is solenoidal in the $\nabla G\times\mathbf{p}$ representation: $\langle\divg\H,\psi\rangle=0$ for every test $\psi$. Amp\`ere--Maxwell uses
\[
\curl(\nabla G\times\mathbf{p})
=\nabla(\mathbf{p}\cdot\nabla G)-\mathbf{p}\nabla^2 G
=\nabla(\mathbf{p}\cdot\nabla G)+k^2\mathbf{p}G+\mathbf{p}\delta^{3}
\]
in $\mathcal{D}'(\mathbb{R}^{3})$, with $\nabla^2G=-k^2G-\delta^{3}$ from $(\nabla^2+k^2)G=-\delta^{3}$ \cite{yaghjian1980dyadic,yaghjian1982}. With $\H=\jj\omega\nabla G\times\mathbf{p}$,
\begin{equation}
\label{eq:ch1.dipole-ampere-distributional}
\curl\H
=\jj\omega\left[
(\nabla\nabla G)\cdot\mathbf{p}
+k^{2}\mathbf{p}G
+\mathbf{p}\delta^{3}
\right]
=\jj\omega\epsilon_0\E+\Jimp,
\end{equation}
because $\E=\epsilon_0^{-1}[k^{2}\mathbf{p}G+(\nabla\nabla G)\cdot\mathbf{p}]$ and $\Jimp=\jj\omega\mathbf{p}\delta^{3}$. Contracting the isotropic contact in Eq.~\eqref{eq:ch1.hessian-distributional} with a constant $\mathbf{p}$ supplies the $-\tfrac13\mathbf{p}\delta^{3}$ contribution in the dyadic Amp\`ere--Maxwell closure \cite{yaghjian1980dyadic,yaghjian1982}. The \emph{distributional} Hessian identity in Eq.~\eqref{eq:ch1.hessian-distributional} is fixed for the spherical limit; what depends on the limiting exclusion geometry in Yaghjian's source-region theory is the depolarisation dyadic obtained when a finite principal volume is contracted to a point---not the statement that Maxwell's equations permit arbitrary contact tensors unrelated to the chosen limit \cite{yaghjian1980dyadic,yaghjian1982}. This chapter adopts the spherical contact term in Eq.~\eqref{eq:ch1.hessian-distributional}; nonspherical exclusion volumes yield a different depolarisation prescription and are not claimed here \cite{yaghjian1980dyadic,yaghjian1982}. Eqs.~\eqref{eq:ch1.dipole-Er}--\eqref{eq:ch1.dipole-Hphi} are the classical restriction of this distribution to $r>0$.
For a charge-continuous impressed current with fixed electric dipole moment $\mathbf{p}$, source scale $\ell\to 0$, and observation points at fixed distance from the support, the exterior fields of Eqs.~\eqref{eq:ch1.dipole-Er}--\eqref{eq:ch1.dipole-Hphi} follow from multipole theory when $\ell$ is small compared with both $k^{-1}$ and the observation distance; $k\ell\ll 1$ alone is not a sufficient convergence criterion without moment bounds and a stated regularization \cite{stratton1941,harrington2001,jackson1999,yaghjian1982}. The $E_r$--$E_\theta$--$H_\phi$ split is the outgoing Hertzian template with the published coefficients \cite{stratton1941,harrington2001}. The same ideal point dipole at
$f=1\,\mathrm{GHz}$ used in Subsection~\ref{sec:ch01.6.3} illustrates shell dependence on
two concentric spheres $S_{r_1}$ and $S_{r_2}$ with $kr_1=0.5$ and $kr_2=5$, as in Figure~\ref{fig:ch1.radial-field-zones}.

In the standard multipole description of an outgoing field,
\begin{equation}
\label{eq:ch1.ex64-multipole-radial}
\mathbf{E}(\mathbf{r})=\mathbf{E}_3(\theta,\phi)\,r^{-3}+\mathbf{E}_2(\theta,\phi)\,r^{-2}+\mathbf{E}_1(\theta,\phi)\,r^{-1}+\cdots,
\qquad
\mathbf{H}(\mathbf{r})=\mathbf{H}_2(\theta,\phi)\,r^{-2}+\mathbf{H}_1(\theta,\phi)\,r^{-1}+\cdots,
\end{equation}
where the angular factors $\mathbf{E}_n,\mathbf{H}_n$ depend on direction but not on $r$ in the source-free exterior \cite{stratton1941}. Near the source, $r^{-3}$ and $r^{-2}$ terms may rival or exceed the $r^{-1}$ radiation field. They describe stored and induction fields—local oscillation of electric and magnetic energy—while the normal Poynting vector may be largely imaginary even though the source already radiates definite outward power. At increasing distance the $r^{-1}$ terms dominate, $\E$ and $\H$ approach the transverse relation of an outgoing wave, and the real Poynting flux carries energy away from the source \cite{poynting1884}.

The radial intensity of that flux follows from the field itself. With
$\mathbf{S}_\mathrm{c}=\tfrac12\E\times\H^*$ and $H_\phi$ from
Eq.~\eqref{eq:ch1.dipole-Hphi} (within Eq.~\eqref{eq:ch1.dipole-fields-explicit}), the radial component is
\begin{equation}
\label{eq:ch1.dipole-sc-radial}
\Re\!\left\{(\mathbf{S}_\mathrm{c}\cdot\hat{\mathbf{r}})\right\}
=\tfrac12\Re\!\left\{E_\theta H_\phi^*\right\},
\end{equation}
since $\H$ has only a $\phi$-component \cite{stratton1941}.
On a sphere of radius $r$, the radiative part of $\Re(S_{\mathrm{c},r})$ scales as $1/r^2$ because both $\E$ and $\H$ carry the $1/r$ factor. The element of area grows as $r^2$, so
\begin{equation}
\label{eq:ch1.ex64-power-invariance}
P_{\mathrm{rad}}=\oint_{S_r}\Re\!\left\{\mathbf{S}_\mathrm{c}\cdot\hat{\mathbf{r}}\right\}dS
=\frac{\omega^4|p_0|^2}{12\pi\varepsilon_0 c^3},
\end{equation}
remains constant in a lossless source-free region between enclosing surfaces \cite{stratton1941}. Total radiated power is unchanged; with increasing $r$ only the local field composition through which that conserved power is conveyed changes.

In the far zone, $U(\theta,\phi)\propto\sin^2\theta$ with peak $U_{\max}$ at $\theta=\pi/2$. Eqs.~\eqref{eq:ch1.radiation-intensity-Ur}--\eqref{eq:ch1.directivity-definition} give
\[
D(\theta,\phi)=\frac{4\pi\,U(\theta,\phi)}{P_{\mathrm{rad}}}
=(3/2)\sin^2\theta,\qquad D_{\max}=(3/2),
\]
using Eq.~\eqref{eq:ch1.ex64-power-invariance} and $\int_{4\pi}\sin^2\theta\,d\Omega=8\pi/3$, so $\int_{4\pi}D(\theta,\phi)\,d\Omega=4\pi$ \cite{stratton1941,harrington2001}.

\begin{figure}[t]
\centering
\VolIPhysFig{\VolIPhysDipoleRadiationPattern}
\caption{Far-zone radiation intensity $U(\theta,\phi)\propto\sin^2\theta$ for an electric dipole aligned with $\hat{\mathbf{z}}$: intensity vanishes on axis and peaks in the equatorial plane. With Eq.~\eqref{eq:ch1.ex64-power-invariance}, $D(\theta)=(3/2)\sin^2\theta$ and $D_{\max}=(3/2)$ \cite{stratton1941,harrington2001}.}
\label{fig:ch1.dipole-radiation-pattern}
\end{figure}

Figure~\ref{fig:ch1.dipole-radiation-pattern} displays the angular side of Eq.~\eqref{eq:ch1.ex64-power-invariance}: $P_{\mathrm{rad}}$ is fixed while its solid-angle distribution is strongly non-uniform. Axisymmetric radiation carries finite power yet may carry zero net $z$-angular-momentum flux when the stress integrand is symmetric about $\hat{\mathbf{z}}$ \cite{beth1936}.

Let $x=kr$. The outgoing transverse factors in Eqs.~\eqref{eq:ch1.dipole-Etheta} and~\eqref{eq:ch1.dipole-Hphi} enter $\mathbf{S}_\mathrm{c}=\tfrac12\E\times\H^*$ as $(1+\jj x-x^2)$ and $(1+\jj x)$. Algebraically,
\begin{equation}
\label{eq:ch1.dipole-shell-algebra}
(1+\jj x-x^2)(1-\jj x)=1+\jj x^3,
\end{equation}
since the $\jj^2$ term cancels the $x^2$ contribution. Substituting Eqs.~\eqref{eq:ch1.dipole-Etheta},~\eqref{eq:ch1.dipole-Hphi}, and~\eqref{eq:ch1.dipole-sc-radial} into $\oint_{S_r}\mathbf{S}_\mathrm{c}\cdot\hat{\mathbf{r}}\,dS$ and performing the $\theta$ integration on $S_r$ gives, for the Hertzian electric dipole of Eqs.~\eqref{eq:ch1.dipole-fields-explicit},
\begin{equation}
\label{eq:ch1.dipole-complex-shell-flux}
\oint_{S_r}\mathbf{S}_\mathrm{c}\cdot\hat{\mathbf{r}}\,dS
=P_{\mathrm{rad}}\left(1-\frac{\jj}{(kr)^3}\right),
\end{equation}
with $P_{\mathrm{rad}}$ from Eq.~\eqref{eq:ch1.ex64-power-invariance}. Equivalently, with the definitions of Eq.~\eqref{eq:ch1.ex64-power-functionals},
\begin{equation}
\label{eq:ch1.ex64-complex-power}
\oint_{S_r}\mathbf{S}_\mathrm{c}\cdot\hat{\mathbf{r}}\,dS
=P_{\mathrm{rad}}+\jj P_{\mathrm{reac}}(r),
\qquad
P_{\mathrm{reac}}(r)=\Im\!\left\{\oint_{S_r}\mathbf{S}_\mathrm{c}\cdot\hat{\mathbf{r}}\,dS\right\}.
\end{equation}
Consistency with Eq.~\eqref{eq:ch1.dipole-complex-shell-flux} requires $P_{\mathrm{reac}}(r)=-P_{\mathrm{rad}}/(kr)^3$ (signed reactive power through $S_r$), so
\begin{equation}
\label{eq:ch1.dipole-reactive-ratio}
\frac{\bigl|P_{\mathrm{reac}}(r)\bigr|}{P_{\mathrm{rad}}}=\frac{1}{(kr)^3}.
\end{equation}
The shell integral $P_{\mathrm{reac}}(r)=\Im\{\oint_{S_r}\mathbf{S}_\mathrm{c}\cdot\hat{\mathbf{r}}\,dS\}$ measures reactive harmonic exchange through $S_r$. It is not a volume energy functional, and it is not Chu's equivalent-circuit $Q_n$ nor the volume ratio $Q_{\mathrm{vol}}(V)$ of Eqs.~\eqref{eq:ch1.chu-q-adopted} and~\eqref{eq:ch1.chu-q-lower-bound} \cite{poynting1884,stratton1941,chu1948,mcclean1996}. In a lossless exterior, $P_{\mathrm{rad}}$ is fixed on nested spheres while $\bigl|P_{\mathrm{reac}}(r)\bigr|$ varies with $r$; the ratio in Eq.~\eqref{eq:ch1.dipole-reactive-ratio} therefore marks shells that are storage-dominated when it is much greater than unity and predominantly radiative when it is much less than unity \cite{stratton1941}.

At $S_{r_1}$ with $kr_1=0.5$, reactive circulation still dominates the normal flux pattern although $P_{\mathrm{rad}}$ is unchanged. Integrating Eqs.~\eqref{eq:ch1.dipole-Etheta} and \eqref{eq:ch1.dipole-Hphi} on $S_{r_1}$ via Eq.~\eqref{eq:ch1.ex64-power-functionals} with $|p_0|=1\,\mathrm{C\,m}$ gives $|P_{\mathrm{reac}}(r_1)|/P_{\mathrm{rad}}\approx 8$. On $S_{r_2}$ with $kr_2=5$, the same integration gives $|P_{\mathrm{reac}}(r_2)|/P_{\mathrm{rad}}\approx 8\times 10^{-3}$, and $\E$--$\H$ approach the far-zone transverse relation
under the outgoing branch and radiation condition stated in section~\ref{sec:ch01.1} \cite{stratton1941}.

Angular-momentum transport demands the same distinction. Outward field angular-momentum flux through $S_r$ is
\begin{equation}
\label{eq:ch1.ex64-am-flux-integral}
\dot J_{z,\mathrm{em,out}}
=-\oint_{S_r}\big[\mathbf{r}\times(\Tavg\cdot\hat{\mathbf{r}})\big]\cdot\hat{\mathbf{z}}\,dS
\end{equation}
by Eqs.~\eqref{eq:ch1.outward-am-flux} and~\eqref{eq:ch1.stress-moment-integral}, and is conserved on nested source-free spheres in steady state when the same outgoing Maxwell branch is followed \cite{stratton1941}. Outward relocation of the enclosing surface does not create angular momentum; it alters the local stress distribution that carries the same signed flux. By contrast, the integral of the magnitude $\int_{S_r}|[\mathbf{r}\times(\Tavg\cdot\hat{\mathbf{r}})]\cdot\hat{\mathbf{z}}|\,dS$ may vary with $r$ because near-zone reactive structure contributes differently to the stress moment than the far-zone radiative field---without implying growth of the conserved net flux \cite{stratton1941,beth1936}.

Three quantities must remain distinct \cite{chu1948,mcclean1996,stratton1941}: the shell exchange $P_{\mathrm{reac}}(r)=\Im\{\oint_{S_r}\mathbf{S}_\mathrm{c}\cdot\hat{\mathbf{r}}\,dS\}$; Chu's equivalent-circuit radiation $Q_n$ and minimum-$Q$ bound for an ideal enclosed spherical radiator; and any finite-volume field-energy ratio defined from $(\E,\H)$.

For later bookkeeping only, define \emph{total} cycle-averaged electric and magnetic field energies in a control volume $V$:
\[
u_e=\tfrac14\varepsilon_0|\E|^2,\qquad u_m=\tfrac14\mu_0|\H|^2,
\qquad
W_e(V)=\int_V u_e\,dV,\qquad W_m(V)=\int_V u_m\,dV,
\]
under $\exp(+\jj\omega t)$ \cite{stratton1941,harrington2001}. These integrals sum the entire field supported in $V$, including propagating and reactive contributions, and they depend on the choice and size of $V$; they are \emph{not} McLean's or Chu's non-propagating stored energies, which require an explicit subtraction of travelling-wave energy \cite{chu1948,mcclean1996}. With $P_{\mathrm{rad}}$ the time-averaged outward real power through $\partial V$ on a predominantly radiative boundary, introduce the volume energy ratio
\begin{equation}
\label{eq:ch1.chu-q-adopted}
Q_{\mathrm{vol}}(V)=\frac{2\omega\max(W_e,W_m)}{P_{\mathrm{rad}}},
\end{equation}
a named finite-volume functional that must not be identified with Chu $Q_n$, antenna radiation $Q$, or $P_{\mathrm{reac}}(r)$ \cite{chu1948,mcclean1996}. If $W_e=W_m$ holds on the same $V$, then $Q_{\mathrm{vol}}=\omega(W_e+W_m)/P_{\mathrm{rad}}$. Chu \cite{chu1948} constructed modal equivalent circuits inside an enclosing sphere to separate non-propagating reactive storage from outward radiation; McLean \cite{mcclean1996} re-examined the exact finite-$ka$ minimum-$Q$ expression and distinguished it from the small-$ka$ expansion. Eq.~\eqref{eq:ch1.chu-q-lower-bound} quotes Chu's widely cited \emph{asymptotic} lower bound $(ka)^{-3}+(ka)^{-1}$ for the ideal minimum-$Q$ spherical model with $ka\ll1$; it is not asserted for arbitrary antennas, for $Q_{\mathrm{vol}}(V)$, or without Chu's interior idealizations \cite{chu1948,harrington2001,mcclean1996}. Under that stated ideal model in the limit $ka\ll1$,
\begin{equation}
\label{eq:ch1.chu-q-lower-bound}
Q_{\mathrm{LB}}=\frac{1}{(ka)^3}+\frac{1}{ka}.
\end{equation}
Local $r^{-3},r^{-2},r^{-1}$ weighting changes with radius; $P_{\mathrm{rad}}$ and nested-sphere $\boldsymbol{\Phi}_{J,\mathrm{out}}$ do not \cite{poynting1884,stratton1941}. Chapter~\ref{ch:02} develops angular-spectrum mappings for bounded sources \cite{harrington2001}.

% R10 BATCH 8 — problems source verified: build_package/ch01-problems.tex (R9 archive ch01-problems-final.tex; see R10_SOURCE_TO_PDF_RECONCILIATION.md)
\section{Summary}
\label{sec:ch01.7}

This chapter has kept angular questions tied to a completed Maxwell field. Impressed sources, constitutive laws, interface data, and outgoing selection fix $(\E,\H)$; on that pair the balances are complex Poynting power, symmetric Maxwell traction for force, and stress-moment flux $\boldsymbol{\Phi}_{J,\mathrm{out}}=-\boldsymbol{\mathcal{T}}_J$ for signed angular-momentum transport through a chosen surface \cite{maxwell1873,poynting1884,beth1936,stratton1941}. In applications, three ideas must remain separate: how the field is \emph{organized} in angle (phase, polarization, modal labels); how much angular momentum is \emph{associated} with the field in a volume or chart; and how much is actually \emph{transported} outward through a physical interface \cite{allen1992,stratton1941}.

The chapter's constructive results are the closed-surface observables, the reactive-versus-radiative shell integrals for a Hertzian dipole, and the scoped modal and helicity scales in Subsection~\ref{sec:ch01.6.1} \cite{stratton1941,harrington2001,chu1948}. Subsection~\ref{sec:ch01.1.3} derives a difference-field energy balance that forces a vanishing far-field pattern and tangential null data on a passive scalar sheet, then continues---via outgoing electric and magnetic vector spherical waves and a Rellich-type unique-continuation hypothesis for source-free Maxwell fields---to annihilate the homogeneous difference field in the connected vacuum exterior under the stated scalar passive impedance class; existence and broader impedance theorems remain cited, not reproduced \cite{sommerfeld1949,harrington2001}. Colton--Kress--Hsiao supplies existence--uniqueness theory for broader impedance classes whose hypotheses are not matched line by line here \cite{colton1981impedance}. Subsection~\ref{sec:ch01.6.4} states distributional Maxwell closure for the ideal electric dipole with the isotropic Hessian contact term and reconciles the $r>0$ Hertzian formulas with Yaghjian's source-region theory \cite{yaghjian1980dyadic,yaghjian1982}. Chu's original stored-energy partition, Jackson's shell-content formulas under this phasor branch, and a general Hansen modal transport law remain bounded explicitly and passed to Chapter~\ref{ch:02} \cite{jackson1999,chu1948,tai1994}. The following volume develops angular-spectrum and dyadic-Green operators on the same outgoing branch, so bounded sources can be analysed without conflating open-aperture modal accounts with closed-sphere stress theorems \cite{harrington2001}.

\phantomsection\label{sec:ch01.7.1}
\phantomsection\label{eq:ch1.continuity-admissibility}%
\phantomsection\label{eq:ch1.continuity-charge-continuous}%
\phantomsection\label{eq:ch1.master-admissibility-chain}%
\phantomsection\label{eq:ch1.master-charge-continuous-chain}%
\phantomsection\label{eq:ch1.report-class-lock}%
\phantomsection\label{sec:ch01.7.2}
\phantomsection\label{sec:ch01.7.3}
\phantomsection\label{eq:ch1.bridge-compatibility-criteria}%
\phantomsection\label{eq:ch1.compatibility-criteria}%

